% Configuration of network topology based on the Physical connection of computers — ICT Upper Sixth — Cours signé OnBuch+
% Official MINESEC syllabus. Compiler avec : tectonic ICT05-configuration-of-network-topology-based-on-the-physical-conn.tex
\newcommand{\DOCMATIERE}{ICT}
\newcommand{\DOCNIVEAU}{Upper Sixth}
\newcommand{\DOCMODULE}{Upper Sixth — ICT}
\newcommand{\DOCLECON}{5}
\newcommand{\DOCTITRE}{Configuration of network topology based on the Physical connection of computers}
% OnBuch+ — préambule « cours signés » ENRICHI (compilé avec tectonic / XeLaTeX)
% Système visuel « Pop » d'OnBuch+ : fond crème (jamais blanc), bordures encre
% épaisses, ombres « dures » décalées, titres Archivo Black en capitales.
% Version MATHÉMATIQUES : graphes (pgfplots/tikz), boîtes pédagogiques
% (définition, propriété, méthode, attention, le savais-tu, exercice résolu,
% exercice, TP, bilan), page de garde et signature OnBuch+ ancrée en bas.
%
% Macros attendues dans main.tex AVANT \input{preamble} :
%   \DOCMATIERE  \DOCNIVEAU  \DOCMODULE  \DOCLECON (numéro)  \DOCTITRE
\documentclass[11pt]{article}

\usepackage{fontspec}
\usepackage[english]{babel}
\usepackage[a4paper,margin=2cm,top=3.4cm,bottom=2.8cm,headheight=1.4cm,footskip=1.3cm]{geometry}
\usepackage{amsmath,amssymb,amsfonts}
\usepackage{mathtools} % \xmapsto, \coloneqq... souvent utilisés par les modèles
\usepackage{cancel} % simplifications barrées (bilans de forces)
% \abs{z} (module, valeur absolue) et \norm{u} (norme), utilisés par les modèles
\providecommand{\abs}{}\let\abs\relax\DeclarePairedDelimiter{\abs}{\lvert}{\rvert}
\providecommand{\norm}{}\let\norm\relax\DeclarePairedDelimiter{\norm}{\lVert}{\rVert}
\usepackage[table]{xcolor} % [table] : fournit aussi \rowcolors (alternance de lignes)
\usepackage{pagecolor}
\usepackage{tikz}
\usetikzlibrary{arrows.meta,positioning,calc,shapes.geometric,shapes.misc,shapes.arrows,decorations.pathmorphing,decorations.pathreplacing,decorations.markings,patterns,shadows,mindmap,trees,fit,backgrounds,matrix,intersections,angles,quotes}
\usepackage{pgfplots}
\usepackage[version=4]{mhchem} % équations nucléaires \ce{^{235}_{92}U}
\usepackage[european,siunitx]{circuitikz} % schémas électriques
\pgfplotsset{compat=1.18}
\usepgfplotslibrary{fillbetween,groupplots} % aires entre courbes (calcul intégral)
\usepackage{enumitem}
\usepackage{fancyhdr}
% [most] charge toutes les bibliothèques tcolorbox (listings, minted, raster,
% documentation...) et gonfle inutilement la mémoire TeX sur un document long
% (~300 boîtes) : on ne charge que ce qui est réellement utilisé.
\usepackage[skins,breakable]{tcolorbox}
\usepackage{graphicx}
\usepackage{colortbl}
\usepackage{booktabs}
\usepackage{tabularx}
\usepackage{diagbox} % case d'en-tête diagonale des tableaux à double entrée
% Colonnes extensibles centrée / gauche / droite (C, L, R), souvent utilisées par les modèles
\newcolumntype{C}{>{\centering\arraybackslash}X}
\newcolumntype{L}{>{\raggedright\arraybackslash}X}
\newcolumntype{R}{>{\raggedleft\arraybackslash}X}
\usepackage{array}
\usepackage{multicol}
\usepackage{siunitx}
% parse-numbers=false : accepte aussi \SI{2,5 \times 10^{-3}}{...}
\sisetup{locale=UK,output-decimal-marker={.},group-minimum-digits=4,parse-numbers=false}
\DeclareSIUnit\ohm{\ensuremath{\symup{\Omega}}} % U+2126 absent de la police
\DeclareSIUnit\division{div} % réglages d'oscilloscope (ms/div, V/div)
\DeclareSIUnit\tr{tr} % tours (vitesse de rotation, ex. \SI{1500}{\tr\per\minute})
\usepackage{qrcode}
% \degree : commande fréquemment utilisée par les modèles pour un angle ou une
% température en dehors de \SI{}{} (non fournie par les paquets chargés).
\providecommand{\degree}{^{\circ}}
% \qed (fin de démonstration), utilisé par les modèles sans amsthm
\providecommand{\qed}{\hfill$\square$}
% \barre : double barre des valeurs interdites dans un tableau de variations (tkz-tab)
\providecommand{\barre}{\ensuremath{\|}}
% environnement proof (amsthm non chargé) utilisé par les modèles
\ifdefined\proof\else\newenvironment{proof}[1][Proof]{\par\noindent\textit{#1.}\ }{\qed\par}\fi
\usepackage{lastpage}
\usepackage{hyperref}
\hypersetup{hidelinks}

% ============================================================
% Palette « Pop » OnBuch+ — mêmes hex que la classe P (plus_ui.dart)
% ============================================================
\definecolor{poporange}{HTML}{F59321}
\definecolor{popdark}{HTML}{E07A0C}
\definecolor{popcream}{HTML}{FFF6E8}
\definecolor{popink}{HTML}{1C1714}
\definecolor{poppurple}{HTML}{7A5AE0}
\definecolor{popgreen}{HTML}{2E9E6A}
\definecolor{poppink}{HTML}{E0457B}
\definecolor{popblue}{HTML}{2F7DE1}
\definecolor{popgold}{HTML}{C98A12}
\definecolor{popmuted}{HTML}{8A7F76}
\definecolor{popline}{HTML}{EADFD2}
\definecolor{popgray}{HTML}{8A7F76}
\colorlet{popgrey}{popgray}
% Alias tolérants (le modèle nomme parfois les couleurs différemment)
\colorlet{popwhite}{white}
\colorlet{popblack}{popink}
\colorlet{popred}{poppink}
\colorlet{popyellow}{popgold}
% Teintes claires pour fonds de tableaux / zones
\colorlet{popblueL}{popblue!10!white}
\colorlet{poporangeL}{poporange!14!white}
\colorlet{popgreenL}{popgreen!12!white}
\colorlet{poppurpleL}{poppurple!12!white}
\colorlet{poppinkL}{poppink!12!white}
\colorlet{popgoldL}{popgold!14!white}

\pagecolor{popcream}

% ============================================================
% Typographie
% ============================================================
\defaultfontfeatures{Ligatures=TeX}
\setmainfont{PlusJakartaSans-Medium.ttf}[Path=../../../fonts/,BoldFont=PlusJakartaSans-Bold.ttf]
% Maths en sans-serif (Fira Math) avec chiffres et lettres droites de Plus
% Jakarta, pour que les formules se fondent dans le texte.
\usepackage[math-style=ISO,bold-style=ISO]{unicode-math}
\setmathfont{FiraMath-Regular.otf}[Path=../../../fonts/]
\setmathfont{PlusJakartaSans-Medium.ttf}[Path=../../../fonts/,range={up/num,up/{Latin,latin},\mathup}]
% (icomma retiré : décimales avec point en anglais)
% Caractères mathématiques Unicode absents des polices de texte (Plus Jakarta,
% Archivo Black) : rendus par leur équivalent mathématique, en texte comme en
% formule — sinon ils s'affichent en carré vide (ex. « Arithmétique dans ℤ »).
\usepackage{newunicodechar}
\newunicodechar{ℝ}{\ensuremath{\mathbb{R}}}
\newunicodechar{ℤ}{\ensuremath{\mathbb{Z}}}
\newunicodechar{ℂ}{\ensuremath{\mathbb{C}}}
\newunicodechar{ℕ}{\ensuremath{\mathbb{N}}}
\newunicodechar{ℚ}{\ensuremath{\mathbb{Q}}}
\newunicodechar{𝔻}{\ensuremath{\mathbb{D}}}
\newunicodechar{α}{\ensuremath{\alpha}}
\newunicodechar{β}{\ensuremath{\beta}}
\newunicodechar{γ}{\ensuremath{\gamma}}
\newunicodechar{δ}{\ensuremath{\delta}}
\newunicodechar{ε}{\ensuremath{\varepsilon}}
\newunicodechar{θ}{\ensuremath{\theta}}
\newunicodechar{λ}{\ensuremath{\lambda}}
\newunicodechar{μ}{\ensuremath{\mu}}
\newunicodechar{σ}{\ensuremath{\sigma}}
\newunicodechar{τ}{\ensuremath{\tau}}
\newunicodechar{φ}{\ensuremath{\varphi}}
\newunicodechar{ψ}{\ensuremath{\psi}}
\newunicodechar{ω}{\ensuremath{\omega}}
\newunicodechar{Σ}{\ensuremath{\Sigma}}
% majuscules produites par \MakeUppercase dans les titres (α-aminés -> Α)
\newunicodechar{Α}{\ensuremath{\alpha}}
\newunicodechar{Β}{\ensuremath{\beta}}
\newunicodechar{Γ}{\ensuremath{\Gamma}}
\newunicodechar{Δ}{\ensuremath{\Delta}}
\newunicodechar{Ε}{\ensuremath{\varepsilon}}
\newunicodechar{Θ}{\ensuremath{\Theta}}
\newunicodechar{Λ}{\ensuremath{\Lambda}}
\newunicodechar{Μ}{\ensuremath{\mu}}
\newunicodechar{Τ}{\ensuremath{\tau}}
\newunicodechar{Φ}{\ensuremath{\Phi}}
\newunicodechar{Ψ}{\ensuremath{\Psi}}
\newunicodechar{Ω}{\ensuremath{\Omega}}
\newunicodechar{↦}{\ensuremath{\mapsto}}
\newunicodechar{∈}{\ensuremath{\in}}
\newunicodechar{∉}{\ensuremath{\notin}}
\newunicodechar{∧}{\ensuremath{\wedge}}
\newunicodechar{⇔}{\ensuremath{\Leftrightarrow}}
\newunicodechar{⟺}{\ensuremath{\Longleftrightarrow}}
\newunicodechar{⇒}{\ensuremath{\Rightarrow}}
\newunicodechar{⟹}{\ensuremath{\Longrightarrow}}
\newunicodechar{∘}{\ensuremath{\circ}}
\newunicodechar{∪}{\ensuremath{\cup}}
\newunicodechar{∩}{\ensuremath{\cap}}
\newunicodechar{≡}{\ensuremath{\equiv}}
\newunicodechar{⊂}{\ensuremath{\subset}}
\newunicodechar{⊥}{\ensuremath{\perp}}
\newunicodechar{∃}{\ensuremath{\exists}}
\newunicodechar{∀}{\ensuremath{\forall}}
\newunicodechar{∛}{\ensuremath{\sqrt[3]{\,}}}
\newunicodechar{∜}{\ensuremath{\sqrt[4]{\,}}}
\newunicodechar{⃗}{\ensuremath{{}^{\to}}}
\newunicodechar{ⁿ}{\ifmmode^{n}\else\textsuperscript{\ensuremath{n}}\fi}
\newunicodechar{ˣ}{\ifmmode^{x}\else\textsuperscript{\ensuremath{x}}\fi}
\newunicodechar{ᵇ}{\ifmmode^{b}\else\textsuperscript{\ensuremath{b}}\fi}
\newunicodechar{ʳ}{\ifmmode^{r}\else\textsuperscript{\ensuremath{r}}\fi}
\newunicodechar{ᵘ}{\ifmmode^{u}\else\textsuperscript{\ensuremath{u}}\fi}
\newunicodechar{ʸ}{\ifmmode^{y}\else\textsuperscript{\ensuremath{y}}\fi}
\newunicodechar{ᵃ}{\ifmmode^{a}\else\textsuperscript{\ensuremath{a}}\fi}
\newunicodechar{ᵉ}{\ifmmode^{e}\else\textsuperscript{\ensuremath{e}}\fi}
\newunicodechar{ᵏ}{\ifmmode^{k}\else\textsuperscript{\ensuremath{k}}\fi}
\newunicodechar{ᵖ}{\ifmmode^{p}\else\textsuperscript{\ensuremath{p}}\fi}
\newunicodechar{ᴺ}{\ifmmode^{N}\else\textsuperscript{\ensuremath{N}}\fi}
\newunicodechar{ᶜ}{\ifmmode^{c}\else\textsuperscript{\ensuremath{c}}\fi}
\newunicodechar{ʰ}{\ifmmode^{h}\else\textsuperscript{\ensuremath{h}}\fi}
\newunicodechar{ᵠ}{\ifmmode^{q}\else\textsuperscript{\ensuremath{q}}\fi}
\newunicodechar{ₙ}{\ifmmode_{n}\else\textsubscript{\ensuremath{n}}\fi}
\newunicodechar{ₐ}{\ifmmode_{a}\else\textsubscript{\ensuremath{a}}\fi}
\newunicodechar{ᵢ}{\ifmmode_{i}\else\textsubscript{\ensuremath{i}}\fi}
\newunicodechar{ᵣ}{\ifmmode_{r}\else\textsubscript{\ensuremath{r}}\fi}
\newunicodechar{ₖ}{\ifmmode_{k}\else\textsubscript{\ensuremath{k}}\fi}
\newunicodechar{ᵦ}{\ifmmode_{\beta}\else\textsubscript{\ensuremath{\beta}}\fi}
\newunicodechar{ₓ}{\ifmmode_{x}\else\textsubscript{\ensuremath{x}}\fi}
\newfontfamily\popdisplay{ArchivoBlack-Regular.ttf}[Path=../../../fonts/]
\newfontfamily\poplogo{SpaceGrotesk-Bold.ttf}[Path=../../../fonts/]
\newfontfamily\popsemi{PlusJakartaSans-SemiBold.ttf}[Path=../../../fonts/]
\newfontfamily\popxbold{PlusJakartaSans-ExtraBold.ttf}[Path=../../../fonts/]
\newfontfamily\popxboldls{PlusJakartaSans-ExtraBold.ttf}[Path=../../../fonts/,LetterSpace=3.0]

\setlist[enumerate]{leftmargin=1.7em,itemsep=5pt,topsep=4pt}
\setlist[itemize]{leftmargin=1.7em,itemsep=4pt,topsep=4pt,label={\color{poporange}\textbullet}}
\setlength{\parindent}{0pt}
\setlength{\parskip}{5pt}
\renewcommand{\arraystretch}{1.25}

% pgfplots : style Pop par défaut
\pgfplotsset{
  every axis/.append style={axis line style={popink,line width=1pt},tick style={popink},
    grid style={popline,line width=0.5pt},label style={font=\small},tick label style={font=\footnotesize}},
  popcurve/.style={poporange,line width=1.8pt,no markers,smooth},
}
% Même style disponible dans un \draw TikZ simple (hors pgfplots)
\tikzset{popcurve/.style={poporange,line width=1.8pt}}
\tikzset{popgrid/.style={popline,very thin}}
\colorlet{popcurve}{poporange} % le nom du style est parfois utilisé comme couleur

% ============================================================
% Pill / squiggle
% ============================================================
\newcommand{\poppill}[3]{%
  \tikz[baseline=-0.55ex]\node[draw=popink,line width=1.2pt,rounded corners=2.6pt,%
    fill=#2,inner xsep=6.5pt,inner ysep=3pt]{%
    \popxboldls\fontsize{7}{7}\selectfont\color{#3}\MakeUppercase{#1}};%
}
\newcommand{\popsquiggle}[1]{%
  \tikz[baseline=-0.4ex]{
    \draw[#1,line width=2.6pt,line cap=round]
      plot [smooth] coordinates {(0,0.09) (0.34,0.01) (0.68,0.21) (1.02,0.01) (1.36,0.10)};
  }%
}
% Mot-clé surligné (marqueur orange sous le texte)
\newcommand{\cle}[1]{{\popxbold\color{popdark}#1}}

% ============================================================
% En-tête / pied
% ============================================================
\pagestyle{fancy}
\fancyhf{}
\renewcommand{\headrulewidth}{0pt}
\renewcommand{\footrulewidth}{0pt}
\lhead{\raisebox{-0.15ex}{\poplogo\fontsize{13}{13}\selectfont\color{popink}OnBuch\color{poporange}+}}
\rhead{\poppill{\DOCMATIERE\ \textbullet\ \DOCNIVEAU\ \textbullet\ Chapter \DOCLECON}{white}{popink}}
\lfoot{\popxbold\fontsize{7.6}{7.6}\selectfont\color{popmuted}\MakeUppercase{Signed course OnBuch+}\ \textbullet\ {\popsemi\DOCTITRE}}
\rfoot{\tikz[baseline=-0.6ex]\node[rounded corners=8.5pt,fill=popink,minimum height=17pt,minimum width=17pt,inner xsep=5pt,inner sep=0]{\popxbold\fontsize{8}{8}\selectfont\color{popcream}\thepage\,/\,\pageref*{LastPage}};}

% ============================================================
% Titres
% ============================================================
\newcommand{\chaptitre}[2]{%
  \begin{tcolorbox}[enhanced,colback=poporange,colframe=popink,boxrule=2.6pt,arc=11pt,
      drop shadow=popink,left=15pt,right=15pt,top=12pt,bottom=11pt,before skip=3pt,after skip=18pt]
    \poppill{#2}{popink}{popcream}\par\vspace{9pt}
    {\raggedright\hyphenpenalty=10000\exhyphenpenalty=10000\popdisplay\fontsize{22}{24}\selectfont\color{popcream}\MakeUppercase{#1}\par}
  \end{tcolorbox}
}
% Section numérotée : pastille orange avec le numéro + titre Archivo
\newcounter{popsec}
\newcommand{\coursec}[1]{%
  \stepcounter{popsec}\setcounter{popsub}{0}%
  \par\vspace{16pt}\noindent
  \begin{minipage}{\linewidth}
  \tikz[baseline=-0.7ex]\node[draw=popink,line width=1.6pt,fill=poporange,rounded corners=4pt,minimum size=22pt,inner sep=0]{\popdisplay\fontsize{12}{12}\selectfont\color{popcream}\thepopsec};%
  \hspace{8pt}{\popdisplay\fontsize{15}{17}\selectfont\color{popink}\MakeUppercase{#1}}\par
  \vspace{4pt}\noindent{\color{poporange}\rule{36pt}{3.6pt}}
  \end{minipage}\par\nopagebreak\vspace{8pt}
}
\newcounter{popsub}
\newcommand{\courssub}[1]{%
  \stepcounter{popsub}%
  \par\vspace{9pt}\noindent{\popxbold\fontsize{11.5}{13}\selectfont\color{popink}{\color{poporange}\thepopsec.\thepopsub}\hspace{6pt}#1}\par\nopagebreak\vspace{4pt}
}

% ============================================================
% Boîtes pédagogiques — même recette : fond blanc, bordure épaisse colorée,
% ombre dure encre, badge plein. Titre optionnel [#1].
% ============================================================
\tcbset{popbase/.style={breakable,enhanced,colback=white,boxrule=2.3pt,arc=9pt,
  shadow={3pt}{-3pt}{0pt}{popink},left=14pt,right=14pt,top=11pt,bottom=11pt,before skip=11pt,after skip=15pt,
  fontupper=\popsemi\fontsize{10.6}{14.5}\selectfont\color{popink},
  fontlower=\popsemi\fontsize{10.6}{14.5}\selectfont\color{popink}}}
% Coche dessinée (le glyphe ✓ n'existe pas dans les polices du document)
\newcommand{\popcheck}[1]{\tikz[baseline=-0.5ex]\draw[#1,line width=1.6pt,line cap=round,line join=round](-0.13,0.0)--(-0.04,-0.09)--(0.13,0.1);}
\providecommand{\coche}{\popcheck{popgreen}}
\providecommand{\resultat}[1]{\par\smallskip\noindent\fcolorbox{popgreen}{popgreenL}{\parbox{\dimexpr\linewidth-2\fboxsep-2\fboxrule\relax}{\textbf{Result.} #1}}\par\smallskip}
\newcommand{\popboxhead}[3]{\poppill{#1}{#2}{white}\ifx\relax#3\relax\else\hspace{6pt}{\popxbold\fontsize{10.6}{12}\selectfont\color{popink}#3}\fi\par\vspace{7pt}}

\newtcolorbox{aretenir}[1][]{popbase,colframe=popgreen,before upper={\popboxhead{Key points}{popgreen}{#1}}}
\newtcolorbox{exemplebox}[1][]{popbase,colframe=poporange,before upper={\popboxhead{Exemple}{poporange}{#1}}}
\newtcolorbox{definition}[1][]{popbase,colframe=popblue,before upper={\popboxhead{Definition}{popblue}{#1}}}
\newtcolorbox{propriete}[1][]{popbase,colframe=poppurple,before upper={\popboxhead{Property}{poppurple}{#1}}}
\newtcolorbox{methode}[1][]{popbase,colframe=popgold,before upper={\popboxhead{Method}{popgold}{#1}}}
\newtcolorbox{attention}[1][]{popbase,colframe=poppink,before upper={\popboxhead{Warning}{poppink}{#1}}}
\newtcolorbox{experience}[1][]{popbase,colframe=popblue,colback=popblueL,before upper={\popboxhead{Experiment}{popblue}{#1}}}
% « Le savais-tu ? » : carte violette pleine (contexte, culture, Cameroun)
\newtcolorbox{savaistu}[1][]{popbase,colback=poppurple,colframe=popink,
  fontupper=\popsemi\fontsize{10.4}{14.2}\selectfont\color{white},
  before upper={\poppill{Did you know?}{white}{poppurple}\ifx\relax#1\relax\else\hspace{6pt}{\popxbold\color{white}#1}\fi\par\vspace{7pt}}}
% Exercice résolu : énoncé en haut, \tcblower puis la solution sur fond crème
\newcounter{popexo}\newcounter{popexr}
\newtcolorbox{exoresolu}[1][]{popbase,colframe=poporange,colbacklower=popcream,
  segmentation style={popink,line width=1.2pt,dashed},
  before upper={\stepcounter{popexr}\popboxhead{Worked example \thepopexr}{poporange}{#1}},
  before lower={\poppill{Solution}{popink}{popcream}\par\vspace{6pt}}}
% Exercice (non résolu en place) — la correction est dans la partie Corrigés
\newtcolorbox{exercice}[1][]{popbase,colframe=popink,
  before upper={\stepcounter{popexo}\popboxhead{Exercise \thepopexo}{popink}{#1}}}
% Corrigé
\newtcolorbox{corrige}[1][]{popbase,colframe=popgreen,colback=popgreenL,
  before upper={\popboxhead{Solution}{popgreen}{#1}}}
% Objectifs / prérequis / bilan
\newtcolorbox{objectifs}{popbase,colframe=popink,colback=poporangeL,
  before upper={\popboxhead{Chapter objectives}{popink}{}}}
\newtcolorbox{prerequis}{popbase,colframe=popink,colback=popblueL,
  before upper={\popboxhead{Prerequisites}{popblue}{}}}
\newtcolorbox{bilan}{popbase,colframe=popink,colback=white,boxrule=2.8pt,
  before upper={\popboxhead{Fiche bilan}{poporange}{L'essentiel à connaître pour l'examen}}}
% Figure encadrée avec légende
\newtcolorbox{popfigure}[1][]{enhanced,breakable=false,colback=white,colframe=popline,boxrule=1.4pt,arc=8pt,
  left=10pt,right=10pt,top=10pt,bottom=8pt,before skip=10pt,after skip=12pt,halign=center,
  lower separated=false,halign lower=center,
  fontlower=\popsemi\fontsize{9}{11.5}\selectfont\color{popmuted},
  #1}
\newcommand{\legende}[1]{\par\vspace{6pt}{\popsemi\fontsize{9}{11.5}\selectfont\color{popmuted}#1}}

% ============================================================
% Signature OnBuch+
% ============================================================
\newcommand{\poplogoinline}[1]{{\poplogo\fontsize{#1}{#1}\selectfont\color{popink}OnBuch\color{poporange}+}}
\newcommand{\signatureonbuch}{%
  \begin{tcolorbox}[enhanced,colback=popink,colframe=popink,boxrule=2.2pt,arc=10pt,
      drop shadow=poporange,left=14pt,right=14pt,top=10pt,bottom=10pt,before skip=0pt,after skip=0pt]
    \tikz[baseline=-0.7ex]\node[circle,fill=popgreen,draw=popcream,line width=1.3pt,minimum size=22pt,inner sep=0]{\popcheck{white}};%
    \hspace{9pt}\begin{minipage}[c]{0.62\linewidth}
      {\popxbold\fontsize{12}{13}\selectfont\color{popcream}Signed\ \textbullet\ {\poplogo OnBuch\color{poporange}+}}\par\vspace{2pt}
      {\raggedright\popsemi\fontsize{8}{10.5}\selectfont\color{popline}\MakeUppercase{OnBuch+ teaching team}\\\MakeUppercase{Follows the official MINESEC syllabus}\par}
    \end{minipage}\hfill
    {\poplogo\fontsize{22}{22}\selectfont\color{popcream}OnBuch\color{poporange}+}
  \end{tcolorbox}%
}

% ============================================================
% Page de garde
% #1 titre · #2 matière · #3 niveau · #4 module · #5 n° leçon · #6 durée ·
% #7 description / objectifs (texte ou itemize)
% ============================================================
\newcommand{\pagegarde}[7]{%
  \thispagestyle{empty}
  \newgeometry{a4paper,margin=1.7cm,top=1.6cm,bottom=1.4cm}%
  \begin{tikzpicture}[remember picture,overlay]
    \fill[poporange!18] ([xshift=-3.2cm,yshift=-3.6cm]current page.north east) circle (4.6cm);
    \fill[poppurple!14] ([xshift=2.4cm,yshift=7.6cm]current page.south west) circle (3.8cm);
  \end{tikzpicture}%
  {\poplogo\fontsize{30}{30}\selectfont\color{popink}OnBuch\color{poporange}+}\hfill
  \raisebox{0.6ex}{\poppill{Signed course \textbullet\ Premium}{popink}{popcream}}\par
  \vspace{1pt}\popsquiggle{poppurple}\par
  \vspace{8pt}
  \begin{tcolorbox}[enhanced,colback=poporange,colframe=popink,boxrule=2.8pt,arc=13pt,
      drop shadow=popink,left=18pt,right=18pt,top=12pt,bottom=13pt,after skip=10pt]
    \poppill{#2 \textbullet\ #3}{popink}{popcream}\hspace{5pt}\poppill{#4}{white}{popink}\par\vspace{8pt}
    {\popdisplay\fontsize{14}{14}\selectfont\color{popink}CHAPTER #5}\par\vspace{4pt}
    {\raggedright\hyphenpenalty=10000\exhyphenpenalty=10000\popdisplay\fontsize{25}{27}\selectfont\color{popcream}\MakeUppercase{#1}\par}
    \vspace{8pt}
    {\popxbold\fontsize{9.5}{9.5}\selectfont\color{popink}\MakeUppercase{Suggested time: #6}}
  \end{tcolorbox}
  \begin{tcolorbox}[enhanced,colback=white,colframe=popink,boxrule=2.2pt,arc=10pt,
      drop shadow=popink,left=16pt,right=16pt,top=10pt,bottom=10pt,after skip=8pt]
    \poppill{In this chapter}{poppurple}{white}\par\vspace{5pt}
    \popsemi\fontsize{9.3}{12.6}\selectfont\color{popink}#7
  \end{tcolorbox}
  \noindent\begin{minipage}[t]{0.487\linewidth}
    \begin{tcolorbox}[enhanced,colback=popgreen,colframe=popink,boxrule=2.2pt,arc=10pt,
        drop shadow=popink,left=11pt,right=11pt,top=10pt,bottom=10pt,height=3.1cm,valign=center]
      \tcbox[colback=white,colframe=popink,boxrule=1.3pt,arc=5pt,boxsep=0pt,nobeforeafter,
        left=3pt,right=3pt,top=3pt,bottom=3pt,valign=center]{\qrcode[height=1.9cm]{https://play.google.com/store/apps/details?id=com.luvvix.onbuch&hl=en}}%
      \hspace{8pt}\begin{minipage}[c]{0.5\linewidth}
        {\popdisplay\fontsize{11}{12.5}\selectfont\color{white}\MakeUppercase{Download OnBuch}}\par\vspace{4pt}
        {\raggedright\popsemi\fontsize{8.4}{11}\selectfont\color{white}Free \textbullet\ Google Play\\AI tutor, past papers, results\par}
      \end{minipage}
    \end{tcolorbox}
  \end{minipage}\hfill
  \begin{minipage}[t]{0.487\linewidth}
    \begin{tcolorbox}[enhanced,colback=poppurple,colframe=popink,boxrule=2.2pt,arc=10pt,
        drop shadow=popink,left=11pt,right=11pt,top=10pt,bottom=10pt,height=3.1cm,valign=center]
      \tikz[baseline=-0.7ex]\node[circle,fill=white,draw=popink,line width=1.3pt,minimum size=1.6cm,inner sep=0]{\popdisplay\fontsize{24}{24}\selectfont\color{poppurple}+};%
      \hspace{8pt}\begin{minipage}[c]{0.6\linewidth}
        {\popdisplay\fontsize{11}{12.5}\selectfont\color{white}\MakeUppercase{Go OnBuch+}}\par\vspace{4pt}
        {\raggedright\popsemi\fontsize{8.4}{11}\selectfont\color{white}All signed courses, unlimited AI tutor, PDF sheets — OnBuch+ tab in the app\par}
      \end{minipage}
    \end{tcolorbox}
  \end{minipage}
  \vspace{8pt}
  \popcontenu
  \vfill
  \signatureonbuch
  \restoregeometry
  \newpage
}

% Rangée « Ce cours contient » (page de garde)
\newcommand{\poptile}[3]{% #1 couleur #2 gros texte #3 légende
  \begin{tcolorbox}[enhanced,colback=#1,colframe=popink,boxrule=2pt,arc=9pt,shadow={2.5pt}{-2.5pt}{0pt}{popink},
      left=6pt,right=6pt,top=9pt,bottom=9pt,halign=center,valign=center,height=1.75cm,width=\linewidth,nobeforeafter]
    {\popdisplay\fontsize{12.5}{13}\selectfont\color{white}\MakeUppercase{#2}}\par\vspace{3pt}
    {\centering\popsemi\fontsize{7.8}{9.5}\selectfont\color{white}#3\par}
  \end{tcolorbox}}
\newcommand{\popcontenu}{%
  \par\noindent{\popxboldls\fontsize{7.5}{7.5}\selectfont\color{popmuted}THIS SIGNED COURSE CONTAINS}\par\vspace{7pt}
  \noindent\begin{minipage}[t]{0.235\linewidth}\poptile{popblue}{Course}{complete, illustrated, step by step}\end{minipage}\hfill
  \begin{minipage}[t]{0.235\linewidth}\poptile{poporange}{Methods}{and worked examples}\end{minipage}\hfill
  \begin{minipage}[t]{0.235\linewidth}\poptile{poppink}{Exercises}{exam style + solutions}\end{minipage}\hfill
  \begin{minipage}[t]{0.235\linewidth}\poptile{popgreen}{Summary sheet}{the essentials to remember}\end{minipage}\par}

% Sommaire
\newtcolorbox{sommaire}{enhanced,colback=white,colframe=popink,boxrule=2.2pt,arc=10pt,
  drop shadow=popink,left=18pt,right=18pt,top=13pt,bottom=13pt,before skip=6pt,after skip=20pt,
  before upper={\poppill{Contents}{poppurple}{white}\par\vspace{9pt}\popsemi\fontsize{10.6}{14.5}\selectfont\color{popink}}}

% Signature de fin de cours — poussée tout en bas de la dernière page
\newcommand{\findecours}{%
  \par\vfill\nopagebreak
  \begin{center}\popsquiggle{poporange}\end{center}\vspace{4pt}
  \signatureonbuch
}
\AtBeginDocument{\def\llbracket{\mathopen{[\mkern-3.5mu[}}\def\rrbracket{\mathclose{]\mkern-3.5mu]}}} % intervalles d'entiers (glyphe ⟦ absent de la police)
% Glyphes absents de Fira Math : redéfinis à partir de glyphes disponibles
\AtBeginDocument{%
  \def\setminus{\mathbin{\backslash}}\let\smallsetminus\setminus
  \def\vdots{\mathord{\vbox{\baselineskip3.2pt\lineskiplimit0pt\kern4pt\hbox{.}\hbox{.}\hbox{.}}}}%
  \def\ddots{\mathinner{\mkern1mu\raise6.5pt\hbox{.}\mkern2mu\raise3.6pt\hbox{.}\mkern2mu\raise0.7pt\hbox{.}\mkern1mu}}%
  \def\triangle{\mathord{\scalebox{0.9}{$\symup{\Delta}$}}}%
  \def\nabla{\mathord{\raisebox{\depth}{\scalebox{1}[-1]{$\symup{\Delta}$}}}}%
  \def\checkmark{\popcheck{popdark}}%
  \def\star{\mathbin{\ast}}\def\bigstar{\mathord{\ast}}%
  \def\diamond{\mathbin{\scalebox{0.6}{$\Diamond$}}}%
  \def\bigcup{\mathop{\mathchoice{\raisebox{-0.3em}{\scalebox{1.7}{$\cup$}}}{\scalebox{1.25}{$\cup$}}{\cup}{\cup}}\displaylimits}%
  \def\bigcap{\mathop{\mathchoice{\raisebox{-0.3em}{\scalebox{1.7}{$\cap$}}}{\scalebox{1.25}{$\cap$}}{\cap}{\cap}}\displaylimits}%
  \def\lesseqqgtr{\mathrel{\lessgtr}}\def\bigoplus{\mathop{\oplus}\displaylimits}%
}
\providecommand{\ssi}{if and only if} % abréviation fréquente
\AtBeginDocument{\providecommand{\wideparen}{\overparen}} % arc de cercle

% Fonctions usuelles que les modèles utilisent sans que amsmath les fournisse
\providecommand{\arsinh}{\operatorname{arsinh}}\providecommand{\arcosh}{\operatorname{arcosh}}
\providecommand{\artanh}{\operatorname{artanh}}\providecommand{\arsech}{\operatorname{arsech}}
\providecommand{\arcsch}{\operatorname{arcsch}}\providecommand{\arcoth}{\operatorname{arcoth}}
\providecommand{\arcsinh}{\operatorname{arsinh}}\providecommand{\arccosh}{\operatorname{arcosh}}
\providecommand{\arctanh}{\operatorname{artanh}}\providecommand{\sech}{\operatorname{sech}}
\providecommand{\csch}{\operatorname{csch}}\providecommand{\arcsec}{\operatorname{arcsec}}
\providecommand{\arccsc}{\operatorname{arccsc}}\providecommand{\arccot}{\operatorname{arccot}}
\providecommand{\sgn}{\operatorname{sgn}}\providecommand{\lcm}{\operatorname{lcm}}
\providecommand{\Var}{\operatorname{Var}}\providecommand{\Cov}{\operatorname{Cov}}

%% FIGFIX-BEGIN v4 — correctifs globaux des figures (docs/onbuch-generation/tools/figfix)
\usepackage{adjustbox}\usepackage{etoolbox}
% toute figure TikZ/pgfplots plus large que la ligne est réduite à la largeur disponible
\BeforeBeginEnvironment{tikzpicture}{\begin{adjustbox}{max width=\linewidth}}
\AfterEndEnvironment{tikzpicture}{\end{adjustbox}}
% légendes pgfplots lisibles par-dessus les courbes
\pgfplotsset{every axis/.append style={legend style={fill=white,fill opacity=0.92,text opacity=1,draw=black!15,inner sep=3pt}}}
% étiquettes d'arêtes (syntaxe edge["..."]) avec fond blanc
\tikzset{every edge quotes/.append style={fill=white,inner sep=1.2pt}}
%% FIGFIX-END
\begin{document}

\pagegarde{Configuration of network topology based on the Physical connection of computers}{ICT}{Upper Sixth}{Upper Sixth — ICT}{5}{7 periods}{This chapter explores how computers are physically interconnected to form networks: the main topologies, their strengths and weaknesses, and the criteria for choosing among them. It then moves from theory to practice by presenting network equipment and guiding the learner through the design and setting up of a complete organisational network.
\vspace{4pt}
\begin{multicols}{2}\small
\begin{itemize}
\item By the end of this chapter I can define a network topology and distinguish physical from logical topology.
\item By the end of this chapter I can describe the structure and functioning of bus, star, ring, mesh and tree (hierarchical) topologies.
\item By the end of this chapter I can compare topologies using criteria such as cost, reliability, scalability and ease of troubleshooting.
\item By the end of this chapter I can justify the choice of a topology for a given real-life situation (school, cybercafé, bank, ministry).
\item By the end of this chapter I can identify and state the role of common network equipment: NIC, hub, switch, router, modem, access point, repeater, cables and connectors.
\item By the end of this chapter I can select appropriate equipment and cabling for a small organisational network given a budget and requirements.
\end{itemize}\end{multicols}}

\begin{sommaire}
\begin{multicols}{2}
\begin{enumerate}
\item Physical Network Topologies
\item Comparing and Choosing a Topology
\item Network Equipment
\item Setting Up a Computer Network
\item Selecting Equipment for Organisational Networks
\item Integration activity
\item Methods and worked examples
\item Exercises
\item Solutions to the exercises
\item Summary sheet
\end{enumerate}
\end{multicols}
\end{sommaire}

\begin{objectifs}
\begin{itemize}
\item By the end of this chapter I can define a network topology and distinguish physical from logical topology.
\item By the end of this chapter I can describe the structure and functioning of bus, star, ring, mesh and tree (hierarchical) topologies.
\item By the end of this chapter I can compare topologies using criteria such as cost, reliability, scalability and ease of troubleshooting.
\item By the end of this chapter I can justify the choice of a topology for a given real-life situation (school, cybercafé, bank, ministry).
\item By the end of this chapter I can identify and state the role of common network equipment: NIC, hub, switch, router, modem, access point, repeater, cables and connectors.
\item By the end of this chapter I can select appropriate equipment and cabling for a small organisational network given a budget and requirements.
\item By the end of this chapter I can plan and carry out the physical setting up of a simple peer-to-peer or client-server network.
\item By the end of this chapter I can test, troubleshoot and document a small network installation.
\end{itemize}
\end{objectifs}

\begin{prerequis}
\begin{itemize}
\item Definition, advantages and types of computer networks (LAN, MAN, WAN) seen in Lower Sixth
\item Basic computer hardware components (system unit, ports, peripherals)
\item Transmission media basics: guided vs unguided media (Lower Sixth introduction to networks)
\item Notion of data transmission, bandwidth and simple units (bit/s, Mb/s)
\item Elementary electricity and safety precautions when handling equipment
\end{itemize}
\end{prerequis}

\newpage

\coursec{Physical Network Topologies}

At a busy cybercafé in Buea, the owner complains that whenever one cable is accidentally unplugged, the whole network goes down and customers lose their work. At a bank branch in Douala, however, a single faulty computer never interrupts the other tellers. Why do two networks built from similar computers behave so differently? The answer lies in how the machines are physically connected — the \cle{network topology} — and in the equipment chosen to link them. In this section, you will learn to identify, analyse and compare the fundamental physical topologies so that you can later design, set up and troubleshoot networks suited to real Cameroonian institutions.

\courssub{What Is a Network Topology?}

\begin{definition}[Network Topology]
The \textbf{network topology} is the arrangement (layout) of nodes (computers, printers, servers, etc.) and the connecting links (cables, wireless channels) in a computer network. It describes the geometric pattern formed by the interconnections.
\end{definition}

\begin{definition}[Physical vs. Logical Topology]
\textbf{Physical topology} refers to the actual layout of cables, devices and connection points — what you see when you walk into the server room or trace the wires under the desks. \\
\textbf{Logical topology} refers to the way data signals actually travel through the network from source to destination, regardless of the physical wiring. \\
\textit{Example:} A modern office wired with twisted-pair cables to a central switch has a \textbf{physical star} topology. If that switch is a \emph{hub} (legacy), data is broadcast to every port $\rightarrow$ \textbf{logical bus}. If the switch is a modern \emph{switch}, data is forwarded only to the destination port $\rightarrow$ \textbf{logical star} (point-to-point).
\end{definition}

\begin{attention}[Common Confusion]
Do \textbf{not} confuse the physical shape with the logical behaviour. A network that \emph{looks} like a star (all cables meeting at a central box) may behave like a bus if the central box is a hub. Always check the \textbf{equipment} (hub vs. switch) to decide the logical topology.
\end{attention}

\courssub{Bus Topology}

In a \textbf{bus topology}, all nodes connect to a single central cable — the \cle{backbone} or \cle{bus} — via drop lines and taps. Signals travel in both directions along the backbone. Each end of the cable must be fitted with a \cle{terminator} (a resistor matching the cable's characteristic impedance) to absorb signals and prevent reflections that would corrupt data.

\begin{popfigure}
\begin{tikzpicture}[scale=0.9, every node/.style={font=\small}]
% Backbone
\draw[popdark, very thick] (0,0) -- (10,0);
% Terminators
\node[draw, circle, fill=popred!20, inner sep=2pt] at (0,0) {T};
\node[draw, circle, fill=popred!20, inner sep=2pt] at (10,0) {T};
% Nodes
\foreach \x/\name in {1.5/PC1, 3.5/PC2, 5.5/PC3, 7.5/PC4, 9/Printer} {
  \draw[popblue, thick] (\x,0) -- (\x,1.5);
  \node[draw, rectangle, rounded corners, fill=popblue!10, minimum width=1.8cm, minimum height=0.8cm] at (\x,2.2) {\name};
}
% Labels
\node[popmuted] at (5,-0.8) {Backbone cable (coaxial or twisted pair)};
\node[popmuted] at (0.5,-1.5) {Terminator (absorbs signal)};
\end{tikzpicture}
\legende{Bus topology: single backbone with terminators at both ends.}
\end{popfigure}

\textbf{How data moves:} When a node transmits, the signal propagates in \emph{both} directions along the backbone. Every node receives the frame, but only the intended destination (matched by MAC address) processes it; others discard it. This is a \textbf{broadcast} medium — only one node can effectively transmit at a time (half-duplex), typically managed by CSMA/CD (Carrier Sense Multiple Access with Collision Detection).

\textbf{Effect of a break:} A single cable break \emph{anywhere} on the backbone splits the network into two isolated segments. \textbf{All} communication stops. This is the cybercafé scenario in our hook.

\begin{aretenir}[Bus Topology — Key Points]
\begin{itemize}
\item Simple, cheap cabling (less cable than star).
\item Easy to extend by adding a tap.
\item \textbf{Single point of failure:} backbone break kills the whole network.
\item Difficult to isolate faults; performance drops as more nodes contend for the bus.
\item Legacy: 10BASE2 (thinnet), 10BASE5 (thicknet). Rarely used for new installations.
\end{itemize}
\end{aretenir}

\courssub{Star Topology}

In a \textbf{star topology}, every node connects via a dedicated point-to-point link to a central device — historically a \cle{hub}, today almost always a \cle{switch}. The central device acts as a traffic controller.

\begin{popfigure}
\begin{tikzpicture}[scale=0.9, every node/.style={font=\small}]
% Central switch
\node[draw, rectangle, rounded corners, fill=poporange!20, minimum width=2.5cm, minimum height=1.5cm] (SW) at (5,2.5) {Switch};
% Nodes around
\foreach \angle/\name/\pos in {45/PC1/right, 135/PC2/left, 225/Printer/left, 315/Server/right} {
  \node[draw, rectangle, rounded corners, fill=popgreen!10, minimum width=1.6cm, minimum height=0.7cm] (N) at (\angle:3) {\name};
  \draw[popblue, thick] (SW) -- (N);
}
% Legend
\node[popmuted] at (5,-0.5) {Point-to-point links (UTP Cat 5e/6)};
\end{tikzpicture}
\legende{Star topology: all nodes connect to a central switch.}
\end{popfigure}

\textbf{Failure modes:}
\begin{itemize}
\item \textbf{One link / one node fails:} Only that node is isolated. The rest of the network continues unaffected (bank teller scenario).
\item \textbf{Central switch fails:} \textbf{Entire network goes down.} This makes the switch a critical single point of failure. Redundancy (dual power supplies, stacking, or a backup switch) is essential in production networks.
\end{itemize}

\begin{attention}[Hub vs. Switch — Physical Star, Different Logical Behaviour]
\begin{itemize}
\item \textbf{Hub (Layer 1 repeater):} Electrically repeats every incoming bit out of \emph{all} other ports. Logical topology = \textbf{bus}. Collisions propagate everywhere. Half-duplex only.
\item \textbf{Switch (Layer 2 bridge):} Learns MAC addresses, forwards frames \emph{only} to the destination port. Logical topology = \textbf{star} (point-to-point). Full-duplex possible. No collisions between different ports.
\end{itemize}
\textbf{Exam tip:} If a question says "star topology" without specifying the central device, assume a modern \textbf{switch}. If it says "hub", explicitly state the logical bus behaviour.
\end{attention}

\courssub{Ring Topology}

In a \textbf{ring topology}, each node connects to exactly two neighbours, forming a closed loop. Data travels in one direction (uni-directional ring) or both directions (bi-directional / dual ring). Access is controlled by a \cle{token} — a special control frame that circulates when the network is idle. A node may transmit only when it captures the token.

\begin{popfigure}
\begin{tikzpicture}[scale=0.9, every node/.style={font=\small}]
% Nodes on a circle
\foreach \i/\name in {0/PC1, 1/PC2, 2/PC3, 3/PC4, 4/Server, 5/Printer} {
  \node[draw, circle, fill=poppurple!10, minimum width=1.6cm] (N\i) at ({90-\i*60}:3) {\name};
}
% Ring links with arrows
\foreach \i in {0,...,5} {
  \pgfmathtruncatemacro{\j}{mod(\i+1,6)}
  \draw[popblue, thick, ->, >=stealth] (N\i) to[bend left=15] (N\j);
}
% Token illustration
\node[draw, circle, fill=popgold, inner sep=1pt, font=\tiny] at (90:3.8) {Token};
\node[popmuted] at (0,-3.5) {Uni-directional ring with token passing};
\end{tikzpicture}
\legende{Ring topology: token circulates; only token holder transmits.}
\end{popfigure}

\textbf{Single vs. Dual Ring:}
\begin{itemize}
\item \textbf{Single ring:} One cable break breaks the whole ring. Used in Token Ring (IEEE 802.5) and FDDI (single attachment stations).
\item \textbf{Dual ring (FDDI, some resilient Ethernet rings):} Two counter-rotating rings. On a break, the ring \emph{wraps} at the neighbours of the fault, restoring a single logical ring. High availability.
\end{itemize}

\begin{aretenir}[Ring Topology — Key Points]
\begin{itemize}
\item Deterministic access (token) $\rightarrow$ predictable latency, good for real-time.
\item No collisions.
\item Adding/removing nodes disrupts the ring (unless using a MAU — Multistation Access Unit — which bypasses failed ports).
\item Largely replaced by switched Ethernet, but concepts appear in resilient metropolitan rings (e.g., CAMTEL backbone rings).
\end{itemize}
\end{aretenir}

\courssub{Mesh Topology}

In a \textbf{mesh topology}, nodes are interconnected with multiple redundant paths. We distinguish:

\begin{itemize}
\item \textbf{Full mesh:} Every node connects directly to every other node.
\item \textbf{Partial mesh:} Only some pairs have direct links; others route through intermediate nodes.
\end{itemize}

\begin{propriete}[Number of Links in a Full Mesh]
For $n$ nodes, a full mesh requires exactly
\[
L = \frac{n(n-1)}{2}
\]
dedicated links. This grows quadratically ($O(n^2)$), making full mesh practical only for small $n$ (e.g., core router interconnects, small critical clusters).
\end{propriete}

\begin{exemplebox}[Calculating Mesh Links]
A ministry department in Yaoundé wants a full mesh among 6 servers for maximum redundancy. How many direct links are needed?
\[
L = \frac{6 \times 5}{2} = 15 \text{ links.}
\]
If they add a 7th server, $L = 21$ links — 6 new cables just for that one server. This is why full mesh is reserved for the \textbf{core layer} of large networks, not the access layer.
\end{exemplebox}

\begin{popfigure}
\begin{tikzpicture}[scale=0.8, every node/.style={font=\small}]
% Full mesh (left)
\begin{scope}[xshift=-4cm]
\node[popmuted] at (0,3.2) {Full mesh ($n=5$, $L=10$)};
\foreach \i/\angle in {0/90, 1/18, 2/-54, 3/-126, 4/-198} {
  \node[draw, circle, fill=popblue!10, minimum width=1.2cm] (N\i) at (\angle:2) {\i};
}
\foreach \i [evaluate=\i as \start using int(\i+1)] in {0,...,3} {
  \foreach \j in {\start,...,4} {
    \draw[popblue, thin] (N\i) -- (N\j);
  }
}
\end{scope}
% Partial mesh (right)
\begin{scope}[xshift=4cm]
\node[popmuted] at (0,3.2) {Partial mesh};
\foreach \i/\angle in {0/90, 1/18, 2/-54, 3/-126, 4/-198} {
  \node[draw, circle, fill=popgreen!10, minimum width=1.2cm] (M\i) at (\angle:2) {\i};
}
\draw[popgreen, thick] (M0) -- (M1) -- (M2) -- (M3) -- (M4) -- (M0);
\draw[popgreen, thick] (M0) -- (M2);
\draw[popgreen, thick] (M1) -- (M3);
\end{scope}
\end{tikzpicture}
\legende{Full mesh vs. partial mesh (5 nodes).}
\end{popfigure}

\textbf{Redundancy \& Routing:} Multiple paths allow automatic failover (via routing protocols like OSPF, or spanning tree at Layer 2). Partial mesh is the standard for WANs and enterprise core networks (e.g., MTN Cameroon backbone).

\courssub{Tree (Hierarchical) Topology}

A \textbf{tree topology} (or \cle{hierarchical topology}) combines multiple star networks by connecting their central switches in a bus or higher-level star arrangement. It organises the network in \textbf{layers}:

\begin{itemize}
\item \textbf{Core layer:} High-speed backbone switches (often partial mesh).
\item \textbf{Distribution (aggregation) layer:} Switches that aggregate access-layer switches, enforce policies, routing between VLANs.
\item \textbf{Access layer:} Switches to which end devices (PCs, APs, printers) connect — pure star.
\end{itemize}

\begin{popfigure}
\begin{tikzpicture}[scale=0.85, every node/.style={font=\small}]
% Core
\node[draw, rectangle, rounded corners, fill=popred!20, minimum width=3cm, minimum height=1.2cm] (Core) at (5,6) {Core Switches (Mesh)};
% Distribution
\foreach \i/\x/\name in {1/1.5/Dist-SW1, 2/5/Dist-SW2, 3/8.5/Dist-SW3} {
  \node[draw, rectangle, rounded corners, fill=poporange!20, minimum width=2.5cm, minimum height=1cm] (D\i) at (\x,4) {\name};
  \draw[popdark, thick] (Core) -- (D\i);
}
% Access
\foreach \i/\dx/\base in {1/1.5/1, 2/5/4, 3/8.5/7} {
  \foreach \j in {0,1} {
    \pgfmathtruncatemacro{\idx}{\base+\j}
    \node[draw, rectangle, rounded corners, fill=popblue!10, minimum width=2cm, minimum height=0.8cm] (A\idx) at ({\dx+\j*1.2},2) {Access-SW\idx};
    \draw[popblue, thick] (D\i) -- (A\idx);
    % End devices
    \foreach \k/\dev in {0/PC, 1/Printer} {
      \node[draw, circle, fill=popgreen!10, minimum width=1cm, font=\tiny] at ({\dx+\j*1.2 + (\k-0.5)*0.8},1) {\dev};
      \draw[popmuted, thin] (A\idx) -- ({\dx+\j*1.2 + (\k-0.5)*0.8},1);
    }
  }
}
% Layer labels
\node[popmuted, font=\small\bfseries] at (0.5,6) {Core};
\node[popmuted, font=\small\bfseries] at (0.5,4) {Distribution};
\node[popmuted, font=\small\bfseries] at (0.5,2) {Access};
\end{tikzpicture}
\legende{Three-layer hierarchical (tree) topology — standard for campus/enterprise LANs.}
\end{popfigure}

\textbf{Use in Cameroon:} University campuses (e.g., University of Buea, University of Yaoundé I), ministry buildings in Yaoundé, large secondary school computer labs. The hierarchy limits broadcast domains, simplifies troubleshooting, and allows modular growth.

\courssub{Hybrid Topologies and Real-World Cameroonian Examples}

A \textbf{hybrid topology} combines two or more basic topologies. In practice, \emph{almost every real network is hybrid}. Recognising the dominant topology in each segment is the key skill.

\begin{experience}[Identifying Topologies in Cameroonian Scenarios]
\textbf{Scenario 1: School Computer Lab (30 PCs)}\\
\textit{Observation:} All PCs wired to a 48-port switch in a locked cabinet.\\
\textit{Physical topology:} \textbf{Star}.\\
\textit{Logical topology:} \textbf{Star} (switch forwards selectively).\\
\textit{Resilience:} One cut cable affects one PC. Switch failure kills lab $\rightarrow$ keep a spare switch.

\textbf{Scenario 2: Cybercafé in Buea (Hook Revisited)}\\
\textit{Observation:} Single coaxial cable (or daisy-chained hubs) running along the wall, PCs tapped off it.\\
\textit{Physical topology:} \textbf{Bus} (or extended star with hubs $\rightarrow$ logical bus).\\
\textit{Failure:} One unplugged cable / broken coax $\rightarrow$ whole café down.\\
\textit{Upgrade path:} Rewire to star with a switch.

\textbf{Scenario 3: Ministry LAN (Multi-building)}\\
\textit{Observation:} Each building has an access switch (star). Building switches connect to a distribution switch in the main block (star of stars = tree). Distribution switches interconnect in a partial mesh for redundancy.\\
\textit{Overall:} \textbf{Hybrid} (Tree at access/distribution, Mesh at core).
\end{experience}

\courssub{Contrasting Physical and Logical Topology: A Visual Summary}

The following diagram shows the same physical star wiring used with a hub (logical bus) versus a switch (logical star).

\begin{popfigure}
\begin{tikzpicture}[scale=0.9, every node/.style={font=\small}]
% LEFT: Hub -> Logical Bus
\begin{scope}[xshift=-5cm]
\node[draw, rectangle, rounded corners, fill=popred!10, minimum width=2.2cm, minimum height=1.2cm] (HUB) at (0,0) {HUB (Layer 1)};
\foreach \angle/\name in {45/PC1, 135/PC2, 225/PC3, 315/PC4} {
  \node[draw, rectangle, rounded corners, fill=popblue!10, minimum width=1.4cm, minimum height=0.6cm] (N) at (\angle:2.5) {\name};
  \draw[popblue, thick] (HUB) -- (N);
}
% Logical flow: broadcast arrows
\draw[popred, thick, ->, >=stealth, dashed] (HUB) -- ++(0,1.5) node[above, font=\tiny, text width=2.5cm, align=center] {Broadcast to ALL ports (Logical Bus)};
\node[popmuted] at (0,-3) {Physical Star \textbackslash Logical Bus};
\end{scope}

% RIGHT: Switch -> Logical Star
\begin{scope}[xshift=5cm]
\node[draw, rectangle, rounded corners, fill=popgreen!10, minimum width=2.2cm, minimum height=1.2cm] (SW) at (0,0) {SWITCH (Layer 2)};
\foreach \angle/\name in {45/PC1, 135/PC2, 225/PC3, 315/PC4} {
  \node[draw, rectangle, rounded corners, fill=popblue!10, minimum width=1.4cm, minimum height=0.6cm] (N) at (\angle:2.5) {\name};
  \draw[popblue, thick] (SW) -- (N);
}
% Logical flow: unicast arrow
\draw[popgreen, thick, ->, >=stealth, dashed] (SW) -- ++(0,1.5) node[above, font=\tiny, text width=2.5cm, align=center] {Forward ONLY to dest. port (Logical Star)};
\node[popmuted] at (0,-3) {Physical Star \textbackslash Logical Star};
\end{scope}
\end{tikzpicture}
\legende{Same physical star cabling; logical topology depends on central device (Hub vs. Switch).}
\end{popfigure}

\courssub{Method: Recognising a Topology from a Network Diagram}

\begin{methode}[Identify the Physical Topology in 3 Steps]
\begin{enumerate}
\item \textbf{Locate the central point(s).} Is there a single device (switch/hub) to which all nodes connect? $\rightarrow$ \textbf{Star}. Is there a single cable running past all nodes? $\rightarrow$ \textbf{Bus}. Do nodes form a closed loop? $\rightarrow$ \textbf{Ring}. Are there many interconnections? $\rightarrow$ \textbf{Mesh}.
\item \textbf{Count the links per node.}
  \begin{itemize}
  \item 1 link (end device) + 1 central device with many ports $\rightarrow$ Star.
  \item 2 links (daisy-chained) $\rightarrow$ Bus or Ring (check if loop closes).
  \item $n-1$ links $\rightarrow$ Full Mesh.
  \end{itemize}
\item \textbf{Check for hierarchy.} If you see distinct layers (core $\rightarrow$ distribution $\rightarrow$ access), it is a \textbf{Tree} (hierarchical) topology, possibly with a meshed core.
\end{enumerate}
\end{methode}

\begin{exoresolu}[Topology Identification]
\textbf{Statement:} A network diagram shows 4 switches labelled SW-A, SW-B, SW-C, SW-D. SW-A connects to SW-B, SW-C and SW-D. SW-B, SW-C, SW-D each connect to 20 PCs. SW-B, SW-C, SW-D are \emph{not} directly connected to each other. Identify the topology.

\textbf{Solution:}
\begin{enumerate}
\item \textbf{Access layer:} Each of SW-B, SW-C, SW-D forms a \textbf{star} with its 20 PCs.
\item \textbf{Distribution/Core:} SW-A connects to the three switches in a \textbf{star} pattern (SW-A is the centre).
\item \textbf{Overall:} A \textbf{Tree (hierarchical) topology} with two levels (core SW-A, access SW-B/C/D). No mesh, no ring, no bus.
\end{enumerate}
\end{exoresolu}

\courssub{Comparing and Choosing a Network Topology (Lesson 15)}

Choosing a topology depends on the organisation's requirements: cost, scalability, fault tolerance, performance, and ease of maintenance. The table below summarises the trade-offs.

\begin{center}
\begin{tabularx}{\linewidth}{|l|X|X|X|X|}\hline
\rowcolor{popblueL}
\textbf{Criterion} & \textbf{Bus} & \textbf{Star} & \textbf{Ring} & \textbf{Mesh (Partial)} \\\hline
\textbf{Cabling cost} & Low (single backbone) & Moderate (dedicated runs) & Moderate & High (multiple links) \\\hline
\textbf{Fault tolerance} & Very poor (backbone break = total outage) & Good (link/PC failure isolated); switch = single point of failure & Poor (single ring); Excellent (dual ring) & Excellent (multiple paths) \\\hline
\textbf{Scalability} & Poor (performance drops with nodes) & Excellent (add ports/switches) & Poor (disrupts ring to add nodes) & Good (add links selectively) \\\hline
\textbf{Performance} & Low (shared medium, collisions) & High (dedicated bandwidth per port, full-duplex) & Deterministic (token), moderate bandwidth & High (load balancing over paths) \\\hline
\textbf{Troubleshooting} & Difficult (fault isolation hard) & Easy (test per link/port) & Moderate (need to locate break) & Complex (requires routing/spanning tree) \\\hline
\textbf{Typical use today} & Legacy only & \textbf{Standard for LAN access} & Metropolitan fibre rings, industrial & Core/backbone, WAN, data centres \\\hline
\end{tabularx}
\end{center}

\begin{methode}[Choosing a Topology — Decision Guide]
\begin{enumerate}
\item \textbf{Small office / school lab ($<$ 50 nodes, budget tight):} \textbf{Star} with a single switch. Low cost, easy to manage, sufficient resilience.
\item \textbf{Multi-floor building / campus:} \textbf{Tree (hierarchical)} — access switches per floor/building, distribution switches, meshed core.
\item \textbf{High-availability server cluster (4–8 nodes):} \textbf{Full mesh} (or dual-ring) for sub-millisecond failover.
\item \textbf{ISP / Telco backbone (MTN, Orange, CAMTEL):} \textbf{Partial mesh} with dynamic routing (OSPF/IS-IS) and MPLS.
\item \textbf{Industrial control / real-time:} \textbf{Dual ring} (or Ethernet ring protection switching — ERPS) for deterministic recovery $<$ 50 ms.
\end{enumerate}
\end{methode}

\begin{attention}[Exam Traps — Choosing Topology]
\begin{itemize}
\item \textbf{"Cheapest for 10 PCs in one room"} $\rightarrow$ Star (UTP + switch), not Bus (coaxial is obsolete, hard to source, no NICs support it).
\item \textbf{"No single point of failure"} $\rightarrow$ Requires \textbf{redundancy}: dual core switches, dual uplinks, partial mesh. A simple star has a SPOF (the switch).
\item \textbf{"Deterministic latency for VoIP/video"} $\rightarrow$ Switched star with QoS is standard today; token ring is legacy. Do not recommend ring for new LANs.
\end{itemize}
\end{attention}

\courssub{Integration Activities (Lesson 16)}

\begin{experience}[Guided Design: Upgrading a Cybercafé Network]
\textbf{Context:} A cybercafé in Limbé has 20 PCs connected via a single coaxial bus (10BASE2). The owner wants to upgrade to support 30 PCs, enable Internet sharing, and eliminate total outages when one cable fails.

\textbf{Step 1 — Requirements analysis:}
\begin{itemize}
\item 30 endpoints $\rightarrow$ need 48-port switch (or 2 $\times$ 24-port stacked).
\item Internet gateway $\rightarrow$ router/firewall (e.g., MikroTik, Ubiquiti).
\item Wireless for customers $\rightarrow$ 2–3 Access Points (APs) connected to switch.
\item Power protection $\rightarrow$ UPS for switch, router, APs.
\end{itemize}

\textbf{Step 2 — Topology selection:} \textbf{Star} (wired) + \textbf{Star} (wireless via APs) = \textbf{Hybrid Star}. All wired runs terminate at a patch panel in a locked cabinet $\rightarrow$ switch.

\textbf{Step 3 — Equipment list (Cameroon market typical):}
\begin{itemize}
\item 1 $\times$ 48-port Gigabit managed switch (TP-Link TL-SG3428 / Cisco SG350 / Ubiquiti USW-48).
\item 1 $\times$ Router/Firewall (MikroTik hEX / Ubiquiti EdgeRouter).
\item 3 $\times$ Ceiling-mount APs (Ubiquiti U6-Lite / TP-Link EAP650).
\item 30 $\times$ Cat 6 UTP cables (3 m drops), 1 $\times$ 24-port patch panel, 1 $\times$ 9U wall cabinet.
\item 1 $\times$ 1.5 kVA UPS (APC / Mercury).
\end{itemize}

\textbf{Step 4 — Implementation checklist:}
\begin{enumerate}
\item Run all cables, label both ends (patch panel port $\leftrightarrow$ wall jack).
\item Terminate on patch panel (T568B standard), test with cable certifier.
\item Connect patch panel to switch ports; connect router WAN to ISP modem, LAN to switch.
\item Configure switch: VLAN 10 (wired PCs), VLAN 20 (wireless guests), VLAN 99 (management).
\item Configure router: DHCP per VLAN, firewall rules (guest VLAN no access to admin VLAN), bandwidth limiting per user.
\item Configure APs: same SSID, WPA2-Enterprise (RADIUS) or captive portal for guests.
\item Document: network diagram, IP plan, VLAN table, switch port map.
\end{enumerate}

\textbf{Step 5 — Verification:} Ping test from each PC, roaming test between APs, failover test (unplug one AP $\rightarrow$ clients roam), bandwidth test (Speedtest).
\end{experience}

\courssub{Network Equipment Overview (Lesson 17 — Introduction)}

Understanding topology requires knowing the devices that implement it. This section introduces the key equipment; detailed configuration is covered in later chapters.

\begin{definition}[Network Equipment Categories]
\begin{itemize}
\item \textbf{End devices (Hosts):} PCs, laptops, smartphones, printers, IP phones, servers — sources/destinations of data.
\item \textbf{Intermediary devices (Network devices):} Forward, filter, or direct data traffic.
  \begin{itemize}
  \item \textbf{Hub (Layer 1):} Multi-port repeater. Obsolete for LANs. Creates one collision domain.
  \item \textbf{Switch (Layer 2):} MAC-address learning, forwards frames per port. Creates separate collision domains per port. Modern standard.
  \item \textbf{Wireless Access Point (AP, Layer 2):} Bridges wireless (802.11) to wired Ethernet. Acts like a hub for RF medium (CSMA/CA).
  \item \textbf{Router (Layer 3):} Interconnects different networks (IP subnets), makes forwarding decisions based on IP address. Essential for Internet access and inter-VLAN routing.
  \item \textbf{Firewall (Layer 3/4/7):} Enforces security policies (allow/deny/inspect). Can be dedicated appliance or router feature.
  \item \textbf{Modem / ONT:} Demarcation device converting ISP signal (fibre GPON, xDSL, 4G/5G) to Ethernet.
  \end{itemize}
\item \textbf{Media:} UTP Cat 5e/6/6a (100 m max), Fibre optic (single-mode / multi-mode, km range), Wireless (2.4/5/6 GHz).
\end{itemize}
\end{definition}

\begin{aretenir}[Equipment–Topology Mapping]
\begin{itemize}
\item \textbf{Bus:} Coaxial cable + BNC T-connectors + terminators (legacy). No active central device.
\item \textbf{Star:} Central \textbf{switch} (wired) or \textbf{AP} (wireless).
\item \textbf{Ring:} Each node has two interfaces; or \textbf{MAU} (Multistation Access Unit) for Token Ring; or fibre switches in ring with ERPS.
\item \textbf{Mesh:} \textbf{Routers} or \textbf{Layer 3 switches} with multiple interfaces running dynamic routing (OSPF, BGP).
\item \textbf{Tree:} \textbf{Access switches} (Layer 2) $\rightarrow$ \textbf{Distribution switches} (Layer 2/3) $\rightarrow$ \textbf{Core switches} (Layer 3, high throughput).
\end{itemize}
\end{aretenir}

\courssub{Summary of Physical Topologies}

\begin{center}
\begin{tabularx}{\linewidth}{|l|X|X|X|}\hline
\rowcolor{popblueL}
\textbf{Topology} & \textbf{Key Physical Trait} & \textbf{Failure Impact} & \textbf{Typical Use in Cameroon} \\\hline
Bus & Single backbone, terminators & Backbone break = total outage & Legacy cybercafés, old school labs \\\hline
Star & Central switch, point-to-point links & Link/PC fail = isolated; Switch fail = total outage & Modern labs, offices, SME LANs \\\hline
Ring & Closed loop, token passing & Single ring: break = total outage; Dual ring: self-heals & Metropolitan fibre rings (CAMTEL), some industrial \\\hline
Full Mesh & Every node $\leftrightarrow$ every node ($n(n-1)/2$ links) & Highly resilient; multiple paths & Core router interconnects, critical server clusters \\\hline
Partial Mesh & Selected redundant links & Good resilience, scalable & Enterprise core, ISP backbone (MTN, Orange) \\\hline
Tree (Hierarchical) & Stars of stars in layers (Core/Dist/Access) & Localised failures; core failure major & University campuses, ministries, large schools \\\hline
\end{tabularx}
\end{center}

\begin{savaistu}[Cameroon Context: ANTIC and Network Standards]
The \textbf{Agence Nationale des Technologies de l'Information et de la Communication (ANTIC)} promotes secure and resilient network architectures for public administration. Their guidelines for government LANs (e.g., in the \emph{Plan National de Développement du Numérique}) recommend \textbf{hierarchical tree topologies} with a \textbf{meshed core}, VLAN segmentation, and redundant links — exactly the hybrid model studied here. When you design a network for a ministry or council, align with this model.
\end{savaistu}

\begin{exercice}[Practice: Topology Analysis]
\begin{enumerate}
\item A small NGO office in Bamenda has 8 laptops connected to a single 8-port \textbf{hub}. Draw the physical topology. What is the logical topology? What happens if the hub fails?
\item A university campus connects 5 faculty buildings. Each building has a star-wired LAN. The 5 building switches are connected in a ring using fibre. Name the overall topology. If the fibre between Building 2 and Building 3 is cut, can Building 1 still reach Building 4? Explain.
\item Calculate the number of direct links required for a full mesh among 12 core routers. Why is a partial mesh preferred for this scale?
\item \textbf{Exam-style:} "A network uses a central device to which all computers are connected by individual cables. If one cable is cut, only that computer loses connectivity. However, if the central device fails, the whole network stops. Name the physical topology and the most likely central device used today." (3 marks)
\item \textbf{Design question:} A secondary school in Douala wants to network 3 computer labs (25 PCs each), an admin block (10 PCs), and a library (15 PCs). The labs are in separate buildings 50 m apart. The admin block is 100 m from the nearest lab. Propose a topology, justify your choice, and list the main equipment needed (switches, fibre/copper, router). (6 marks)
\end{enumerate}
\end{exercice}

\begin{corrige}[Exercise 1]
\begin{enumerate}
\item Physical: \textbf{Star}. Logical: \textbf{Bus} (hub broadcasts). Hub failure $\rightarrow$ total outage.
\item Overall: \textbf{Hybrid (Tree of Stars + Ring)}. Yes, traffic goes the other way around the ring (Building 1 $\rightarrow$ Building 5 $\rightarrow$ Building 4).
\item $L = 12 \times 11 / 2 = 66$ links. Partial mesh preferred because 66 cables/ports is expensive and hard to manage; routing protocols provide redundancy with fewer links.
\item Physical topology: \textbf{Star}. Central device: \textbf{Switch} (modern) or \textbf{Hub} (legacy). Today: \textbf{Switch}.
\item \textbf{Proposed topology:} \textbf{Tree (hierarchical)}.\\
\textbf{Justification:} Separate buildings $\rightarrow$ each building = access layer (star). Building switches uplink to a distribution switch in admin block (or central lab). Fibre for inter-building links (distance $>$ 100 m, immunity to lightning). Copper (Cat 6) inside buildings. Router at distribution for Internet/VLAN routing.\\
\textbf{Equipment:}
\begin{itemize}
\item 3 $\times$ 24-port PoE+ access switches (labs + library), 1 $\times$ 12-port switch (admin).
\item 1 $\times$ 24-port Layer 3 distribution switch (core) with SFP+ uplinks.
\item 4 $\times$ SFP+ modules + 4 $\times$ OM3 multi-mode fibre patch cords (inter-building).
\item Cat 6 UTP for all horizontal runs (patch panels, wall jacks).
\item 1 $\times$ Router/Firewall (Internet gateway, VLAN routing).
\item 1 $\times$ UPS for distribution switch + router.
\end{itemize}
\end{enumerate}
\end{corrige}

\coursec{Comparing and Choosing a Topology}

In the previous section we described the five main \cle{physical topologies} --- bus, star, ring, mesh and tree --- and how computers are physically connected in each of them. Knowing the shapes, however, is not enough. A network engineer in Douala or Bamenda who is asked to network a school computer laboratory, a bank branch or a hospital must answer a harder question: \emph{which topology should we actually install, and why?} This section gives you the comparison criteria, a detailed comparison of the five topologies, and a rigorous method for justifying a choice --- exactly what GCE Advanced Level examiners expect.

\courssub{Comparison Criteria}

\begin{definition}[Comparison criteria for network topologies]
A \cle{comparison criterion} is a measurable or observable property used to evaluate and rank topologies. The five standard criteria are: \textbf{cost} (cabling and equipment), \textbf{reliability / fault tolerance} (what happens when a cable or device fails), \textbf{performance} (speed, collisions, delays), \textbf{scalability} (ease of adding new stations), and \textbf{ease of installation and maintenance} (including fault detection).
\end{definition}

Let us build the intuition behind each criterion before comparing.

\textbf{Cost.} Cabling is often the largest hidden expense of a network. A bus uses one single backbone cable, so it is cheap in cable length; a full mesh with $n$ computers needs $\dfrac{n(n-1)}{2}$ separate links, which explodes quickly. Equipment cost also matters: a star needs a central switch, while a bus needs only terminators.

\textbf{Reliability / fault tolerance.} Ask: \emph{if one cable is cut, does the whole network stop?} In a bus, cutting the backbone kills everything. In a star, cutting one branch cable disconnects only one computer --- but if the \emph{switch} fails, everything stops (single point of failure). In a mesh, traffic simply takes another path.

\textbf{Performance and collisions.} On a shared medium (bus, or a ring without proper token control), two stations transmitting simultaneously cause a \cle{collision}, and both must retransmit. The more stations, the more collisions, and the slower the effective throughput. A switched star gives each computer a dedicated link to the switch, so collisions are practically eliminated.

\textbf{Scalability.} Adding a computer to a star means plugging one cable into a free switch port --- seconds of work. Adding one to a bus may require shutting the network down and extending the backbone; adding one to a full mesh requires a new cable to \emph{every} existing machine.

\textbf{Ease of installation and maintenance.} Centralised topologies (star, tree) make fault-finding easy because all connections meet at one point. Distributed topologies (bus, mesh) force the technician to hunt along walls and ceilings for the faulty segment.

\courssub{Detailed Comparison of the Five Topologies}

\begin{popfigure}
\begin{center}
\small
\begin{tabularx}{\linewidth}{|l|X|X|X|X|X|}
\hline
\rowcolor{popblueL} \textbf{Criterion} & \textbf{Bus} & \textbf{Star} & \textbf{Ring} & \textbf{Mesh} & \textbf{Tree} \\
\hline
Cost of cabling \& equipment & Very low: one backbone, no central device & Moderate: one cable per host + switch & Moderate: ring cable + MAU & Very high: $\frac{n(n-1)}{2}$ links, many ports & High: hierarchical cabling + several switches \\
\hline
Reliability / fault tolerance & Poor: backbone break stops all & Good per link; switch is single point of failure & Poor: one break stops ring (unless dual ring) & Excellent: multiple alternative paths & Good: failure isolated to one branch; root switch critical \\
\hline
Performance / collisions & Degrades fast: shared medium, collisions & High: dedicated links, no collisions with switches & Fair: token passing gives orderly access, delay grows with $n$ & Very high: direct dedicated links & High within branches; backbone can bottleneck \\
\hline
Scalability & Poor: adding hosts is disruptive & Excellent: plug into a free port & Poor: ring must be opened & Very poor: new links to all nodes & Excellent: add a branch or a switch \\
\hline
Installation \& maintenance & Simple to install, hard to troubleshoot & Easy to install; faults found centrally & Moderate; fault location tricky & Very complex to install and manage & Moderate; structured, manageable per branch \\
\hline
\end{tabularx}
\end{center}
\legende{Comparison of the five physical topologies across the five standard criteria.}
\end{popfigure}

\begin{aretenir}[Advantages and disadvantages of each topology]
\textbf{Bus} --- Advantages: cheapest cabling, simple, uses little cable. Disadvantages: backbone failure stops everything, collisions, hard fault-finding, limited size.\par
\textbf{Star} --- Advantages: easy to install, extend and troubleshoot; a cable failure affects only one host; excellent performance with switches. Disadvantages: needs a central device (cost); switch failure stops the whole network; more cable than a bus.\par
\textbf{Ring} --- Advantages: orderly access (token), no collisions, predictable delay. Disadvantages: one break can stop the ring, adding/removing hosts is disruptive, largely obsolete.\par
\textbf{Mesh} --- Advantages: outstanding fault tolerance, dedicated high-speed links, secure point-to-point paths. Disadvantages: extremely expensive in cable and ports, complex to install and configure.\par
\textbf{Tree} --- Advantages: highly scalable, groups departments naturally, faults isolated per branch. Disadvantages: depends on backbone and root switch, more cabling and equipment than a simple star.
\end{aretenir}

\begin{propriete}[Dominance of the star topology]
For the great majority of local area networks, the \cle{star topology} offers the best balance of cost, reliability and scalability: a single cable failure isolates only one machine, new stations are added in seconds, and modern switches eliminate collisions at moderate cost. This is why the star (and its extension, the tree) dominates virtually all modern Ethernet LANs. \emph{(Justification discussed in class; not a formal proof to reproduce in the exam.)}
\end{propriete}

\courssub{Decision Factors in Context}

Choosing a topology is never done in the abstract. The correct answer always depends on the \emph{scenario}:

\begin{itemize}
\item \textbf{Size of the organisation.} A handful of machines can live with almost anything; hundreds of machines demand the structure and scalability of a star or tree.
\item \textbf{Budget.} A cybercafé in Buea with limited capital will avoid mesh; a bank can afford redundancy.
\item \textbf{Expected growth.} If the organisation plans to double its computers in two years, scalability (star/tree) outweighs initial savings (bus).
\item \textbf{Criticality of service.} In a hospital, where a network failure can endanger lives, or in a bank processing mobile-money transactions, fault tolerance dominates the decision. In a small cybercafé, an hour of downtime is annoying but not catastrophic, so cost may dominate.
\end{itemize}

\begin{attention}[The mesh trap]
A very common mistake --- in the exam and in real life --- is to answer ``mesh'' whenever reliability is mentioned, without considering its cabling cost and complexity. With $n$ computers, a full mesh needs $\dfrac{n(n-1)}{2}$ links: 30 computers already require $435$ cables! Mesh is justified only for a \emph{small number} of critical interconnection points (e.g. linking backbone routers). For end-user machines, a star with a backup switch usually gives sufficient reliability at a fraction of the cost.
\end{attention}

\courssub{Case Analyses}

\begin{exemplebox}[Case 1 --- A 20-computer school laboratory]
Needs: low budget, one room, easy maintenance by one teacher, moderate reliability. \textbf{Choice: star.} One 24-port switch at the teacher's desk, one cable per computer. A faulty cable disconnects only one student machine; adding computers later means plugging into free ports. A bus would be cheaper in cable but a single backbone fault would stop the whole class --- unacceptable during a practical exam session.
\end{exemplebox}

\begin{exemplebox}[Case 2 --- A bank branch]
Needs: high availability (transactions must never stop), security, moderate number of stations. \textbf{Choice: star with redundancy} --- a star per floor or department, with a backup switch and dual links to the branch server; the interconnection between the critical core devices may use a partial mesh. Pure mesh for every teller machine would be needlessly expensive; a bus or ring would risk total outage.
\end{exemplebox}

\begin{exemplebox}[Case 3 --- Connecting three buildings of a campus]
Needs: each building has its own LAN; buildings must communicate; the network must grow. \textbf{Choice: tree (hierarchical star-of-stars).} Each building keeps an internal star with its own switch; the three building switches connect to a central backbone switch (often by fibre optic). A fault in one building does not affect the others, and a fourth building can be added later by connecting one more branch.
\end{exemplebox}

\courssub{A Rigorous Method for Justifying a Choice}

\begin{methode}[Four-step procedure to justify a topology choice]
\begin{enumerate}
\item \textbf{Identify the needs} from the scenario: number of stations, budget, expected growth, criticality of service, physical layout (one room? several buildings?).
\item \textbf{List the candidate topologies} that could physically fit (usually two or three of bus, star, ring, mesh, tree).
\item \textbf{Compare the candidates on at least three criteria} (cost, reliability, performance, scalability, maintenance), linking each criterion explicitly to the scenario.
\item \textbf{Conclude and justify}: state the chosen topology and explain why it beats the rejected candidates on the decisive criteria, acknowledging its disadvantages and how they are mitigated.
\end{enumerate}
\end{methode}

\begin{popfigure}
\begin{center}
\begin{tikzpicture}[
node distance=0.55cm,
startstop/.style={rectangle, rounded corners, draw=popdark, fill=popblueL, text width=4.6cm, align=center, minimum height=0.9cm, font=\small},
decision/.style={diamond, draw=popdark, fill=popgreenL, text width=3.4cm, align=center, inner sep=1pt, font=\small},
process/.style={rectangle, draw=popdark, fill=poporange!25, text width=3.6cm, align=center, minimum height=0.9cm, font=\small},
arrow/.style={-{Stealth}, thick, popdark}]
\node[startstop] (start) {Read scenario: size, budget, growth, criticality};
\node[decision, below=of start] (crit) {Service highly critical?};
\node[decision, below left=0.9cm and -0.4cm of crit] (many) {Many critical core links?};
\node[decision, below right=0.9cm and -0.4cm of crit] (budget) {Very tight budget, few PCs?};
\node[process, below=of many] (mesh) {Partial mesh for core + star for users};
\node[process, below=of budget] (bus) {Bus (small, temporary) or simple star};
\node[decision, below=1.1cm of crit] (grow) {Growth expected / several rooms?};
\node[process, below left=0.7cm and -0.6cm of grow] (tree) {Tree (star of stars)};
\node[process, below right=0.7cm and -0.6cm of grow] (star) {Star with switch};
\draw[arrow] (start) -- (crit);
\draw[arrow] (crit) -| node[above, font=\scriptsize]{yes} (many);
\draw[arrow] (crit) -| node[above, font=\scriptsize]{no} (budget);
\draw[arrow] (many) -- node[right, font=\scriptsize]{yes} (mesh);
\draw[arrow] (budget) -- node[right, font=\scriptsize]{yes} (bus);
\draw[arrow] (many.south) |- node[near start, above, font=\scriptsize]{no} (grow.west);
\draw[arrow] (budget.south) |- node[near start, above, font=\scriptsize]{no} (grow.east);
\draw[arrow] (grow) -| node[above, font=\scriptsize]{yes} (tree);
\draw[arrow] (grow) -| node[above, font=\scriptsize]{no} (star);
\end{tikzpicture}
\end{center}
\legende{Decision flowchart: from requirements to a justified topology choice.}
\end{popfigure}

\begin{exoresolu}[GCE-style question]
A new regional hospital in Bamenda has 45 computers spread over three wards. The network carries patient records and must remain available at all times. The hospital expects to open a fourth ward next year. Recommend a physical topology, justifying your answer with at least three criteria.
\tcblower
\textbf{Step 1 --- Needs:} 45 stations, several buildings/wards, high criticality (patient records, 24/7 availability), planned growth (a fourth ward). Budget is secondary to reliability.\par
\textbf{Step 2 --- Candidates:} star (per ward), tree (campus-wide), mesh (for core links). Bus and ring are rejected immediately: a single break would stop a ward.\par
\textbf{Step 3 --- Comparison on criteria:}
\begin{itemize}
\item \emph{Reliability:} a tree isolates a fault to one ward, so the other wards keep working --- essential for a hospital. A single flat star would make the whole hospital depend on one switch.
\item \emph{Scalability:} the fourth ward is added next year by connecting one new branch switch to the backbone --- a tree handles this naturally; a bus or ring would require disruptive rewiring.
\item \emph{Performance:} each ward's star gives dedicated switched links (no collisions), and traffic between wards uses the backbone only when needed.
\item \emph{Cost:} a full mesh for 45 machines would need $\frac{45\times 44}{2}=990$ links --- absurdly expensive; the tree achieves the required reliability at a fraction of that cost.
\end{itemize}
\textbf{Step 4 --- Conclusion:} we recommend a \textbf{tree topology} (a star per ward, ward switches linked to a central backbone switch, with a backup backbone link for extra fault tolerance). It best satisfies reliability, scalability and performance for this critical, growing hospital network, at a justifiable cost.
\end{exoresolu}

\begin{attention}[Exam tip --- link criteria to the scenario]
In GCE questions, marks are awarded for \emph{argued} criteria, not for reciting definitions. Never write ``star is reliable'' in isolation. Write: ``a star is chosen \textbf{because} in this laboratory a cable fault must not stop the other 19 students during practical examinations.'' Every criterion you cite must be tied to a fact given in the scenario (size, budget, growth, criticality). A correct topology with no scenario-linked justification earns only partial credit.
\end{attention}

\begin{exercice}[Choosing and justifying]
A cybercafé in Limbe has 8 computers in one small room and a very limited budget; the owner may add 4 more computers next year. A microfinance bank in Yaoundé has 25 workstations and must guarantee continuous service for mobile-money operations. For each organisation, recommend a topology and justify it with at least three criteria linked to the scenario.
\end{exercice}

\begin{corrige}[Exercise 1]
\textbf{Cybercafé:} star with a small 16-port switch. Criteria: (i) \emph{cost} --- one inexpensive switch and short cables suit the limited budget; (ii) \emph{scalability} --- the 4 future computers plug into the free ports, so growth is effortless; (iii) \emph{maintenance} --- the owner can identify a faulty machine instantly at the switch, important since there is no paid technician. A bus would be slightly cheaper but a backbone fault would close the whole café.\par
\textbf{Microfinance bank:} star with redundancy (backup switch, dual link to the server), possibly partial mesh between the core devices. Criteria: (i) \emph{reliability} --- continuous service for mobile-money transactions is critical, so no single cable or switch failure may stop operations; (ii) \emph{performance} --- switched dedicated links avoid collisions and keep transaction times short; (iii) \emph{fault isolation} --- a faulty teller cable disconnects only that teller. Full mesh for all 25 stations (300 links) is rejected as unnecessarily costly.
\end{corrige}

\coursec{Network Equipment}

In the previous section we compared the physical topologies (bus, ring, star, mesh) and concluded that the \cle{star topology} is the standard choice for modern local area networks. But a topology is only a \emph{plan} on paper: to transform that plan into a working network, we need concrete \cle{network equipment} --- the cards, cables, connectors and devices that physically carry the data. In this section we study each piece of equipment in turn: what it is, what it does, and where it is used. At the GCE Advanced Level, examiners frequently present a network diagram and ask: \emph{``State the role of device X in this network.''} By the end of this section you will be able to answer such questions precisely.

\courssub{The Network Interface Card (NIC)}

\begin{definition}[Network Interface Card]
A \cle{Network Interface Card} (\cle{NIC}) is the hardware component that connects a computer (or printer, or any other device) to a network. It prepares the computer's data for transmission over the medium and controls the flow of data between the computer and the cable or radio channel.
\end{definition}

\textbf{Intuition.} Think of the NIC as the ``mouth and ears'' of the computer on the network. Without it, a computer simply cannot speak to or hear other machines. On modern motherboards the NIC is usually built in (onboard); on older machines it is a separate expansion card (PCI, PCIe, or USB).

Every NIC manufactured in the world carries a unique physical address called the \cle{MAC address} (Media Access Control address). It is a 48-bit number, usually written as six pairs of hexadecimal digits, for example \texttt{3C:5A:B4:0F:2E:91}. The first three octets identify the manufacturer (OUI), the last three are a unique serial number. The MAC address is used by switches to identify exactly which device is attached to which port.

There are two main variants:
\begin{itemize}
\item \textbf{Wired NIC:} has an RJ45 port for an Ethernet cable; typical speeds are 100 Mbit/s (Fast Ethernet), 1 Gbit/s (Gigabit Ethernet), or 2.5/10 Gbit/s on newer hardware.
\item \textbf{Wireless NIC (Wi-Fi card):} has an antenna and communicates by radio waves with a wireless access point, following the IEEE 802.11 standards (Wi-Fi 4/5/6/6E/7).
\end{itemize}

\begin{attention}[NIC drivers and binding]
A NIC requires a \emph{device driver} in the operating system. In the GCE practical, you may be asked to verify that the NIC is recognised (e.g. \texttt{ipconfig /all} on Windows or \texttt{ip link} on Linux) and that it has a valid MAC address.
\end{attention}

\courssub{Transmission Media and Connectors}

Data must travel on a physical medium. The three media you must know for the examination are the following.

\textbf{Twisted pair cable.} This is the most common LAN cable. It contains four pairs of copper wires twisted together to reduce electromagnetic interference (EMI) and crosstalk. Two versions exist:
\begin{itemize}
\item \cle{UTP} (Unshielded Twisted Pair): no metallic shield; cheap, flexible, easy to install; used in offices, schools and homes.
\item \cle{STP} (Shielded Twisted Pair): each pair or the whole cable is wrapped in a metallic foil/braid; more resistant to interference; used in factories, near heavy electrical equipment, or for outdoor runs.
\end{itemize}
Twisted pair cables are terminated with an \cle{RJ45 connector} (8P8C) and are classified into \cle{categories}:
\begin{center}
\begin{tabularx}{\linewidth}{|l|X|X|}\hline
\rowcolor{popblueL}
\textbf{Category} & \textbf{Max Data Rate (typical)} & \textbf{Max Distance} \\ \hline
Cat5e & 1 Gbit/s & 100 m \\ \hline
Cat6 & 1 Gbit/s (100 m), 10 Gbit/s (up to 55 m) & 100 m \\ \hline
Cat6a & 10 Gbit/s & 100 m \\ \hline
Cat7 / Cat7a & 10--40 Gbit/s & 100 m \\ \hline
\end{tabularx}
\end{center}
The maximum recommended length of a single twisted pair segment (horizontal cabling) is \textbf{100 m} (90 m solid cable + 10 m patch cords).

\begin{attention}[Straight vs crossover cables]
A \cle{straight-through cable} (T568B at both ends) has identical wire order at both ends and connects \emph{different} types of devices (PC to switch, switch to router). A \cle{crossover cable} (T568A at one end, T568B at the other) swaps the transmit and receive pairs (pins 1\&2 with 3\&6) and connects \emph{similar} devices directly (PC to PC, switch to switch on old equipment without auto-MDIX). Modern Gigabit interfaces implement \textbf{auto-MDIX} and automatically detect the cable type, making crossover cables rarely needed today. Confusing the two is a classic practical-exam mistake.
\end{attention}

\begin{popfigure}
\begin{center}
\begin{tikzpicture}[scale=0.9, font=\small]
% RJ45 plug body
\draw[fill=popblueL, draw=popdark, thick] (0,0) rectangle (3.4,2.2);
\draw[fill=white, draw=popdark] (0.4,-0.5) rectangle (3.0,0);
\node[popdark] at (1.7,-0.25) {locking clip};
% pins with T568B colors
\foreach \i/\col in {0/poporange, 1/poporange, 2/popgreen, 3/popblue, 4/popblue, 5/popgreen, 6/brown, 7/brown}{
  \pgfmathsetmacro{\x}{0.25+\i*0.4}
  \draw[fill=\col, draw=popdark] (\x,2.2) rectangle (\x+0.28,2.6);
}
\node[popdark, align=center] at (1.7,3.0) {RJ45 plug (pins up, clip down)};
% colour labels
\foreach \i/\name/\col in {0/{1 W-O}/poporange, 1/{2 O}/poporange, 2/{3 W-G}/popgreen, 3/{4 Bl}/popblue, 4/{5 W-Bl}/popblue, 5/{6 G}/popgreen, 6/{7 W-Br}/brown, 7/{8 Br}/brown}{
  \pgfmathsetmacro{\x}{0.39+\i*0.4}
  \node[rotate=90, anchor=east, font=\scriptsize, popdark] at (\x,1.1) {\name};
}
% cable
\draw[line width=4pt, popblue] (3.4,1.1) -- (5.2,1.1);
\node[popdark, anchor=west] at (5.3,1.1) {UTP cable (4 twisted pairs)};
\end{tikzpicture}
\end{center}
\legende{RJ45 connector pinout (T568B standard) viewed with the locking clip facing down. Pins 1 \& 2 = transmit pair (orange), pins 3 \& 6 = receive pair (green).}
\end{popfigure}

\textbf{Coaxial cable.} A central copper conductor surrounded by insulation, a metallic braid (shield) and an outer jacket. It was used in old bus networks (10BASE2 ``thinnet'', 10BASE5 ``thicknet'') and is still used for cable television (CATV) and some broadband Internet (DOCSIS). It is largely obsolete in LANs but you may encounter it in legacy installations.

\textbf{Fibre optic cable.} Transmits data as pulses of light through a glass or plastic core. It offers very high speeds, immunity to electromagnetic interference and distances of several kilometres. Two main types:
\begin{itemize}
\item \textbf{Single-mode fibre (SMF):} small core (8--10 \textmu m), laser source, distances up to 100+ km; used for ISP backbones (CAMTEL, MTN fibre links between towns).
\item \textbf{Multi-mode fibre (MMF):} larger core (50/62.5 \textmu m), LED/VCSEL source, distances up to 550 m (10 Gbit/s); used for campus backbones, data centres.
\end{itemize}
Common connectors: SC, LC, ST, FC. In Cameroon, FTTH (Fibre To The Home) deployments by CAMTEL, MTN and Orange typically use an \cle{ONT} (Optical Network Terminal) that converts light to Ethernet.

\courssub{The Hub}

\begin{definition}[Hub]
A \cle{hub} (multiport repeater) is a Layer 1 device that connects several computers in a physical star topology. When it receives bits on one port, it simply \textbf{regenerates and broadcasts} those bits to \emph{all} its other ports, regardless of the destination.
\end{definition}

\textbf{Intuition.} A hub behaves like a town crier in the market square: it shouts every message to everybody, and each person must decide whether the message concerns them. This creates two serious limitations:
\begin{itemize}
\item \textbf{Collisions:} all ports share a single \cle{collision domain}; if two computers transmit at the same time, the signals collide and must be re-sent (CSMA/CD), wasting time.
\item \textbf{Wasted bandwidth and poor security:} every machine receives traffic meant for others (promiscuous mode can capture all frames).
\end{itemize}
For these reasons the hub is now \cle{obsolete} and has been replaced by the switch. It survives in the syllabus mainly so that you can \emph{contrast} it with the switch.

\courssub{The Switch}

\begin{definition}[Switch]
A \cle{switch} (multiport bridge) is a Layer 2 device that connects several devices in a star topology and forwards data frames \textbf{only to the port of the destination device}, using the destination \cle{MAC address} contained in each frame.
\end{definition}

\textbf{How it works.} The switch maintains a \emph{MAC address table} (CAM table) that records which MAC address is reachable through which port. When a frame arrives, the switch reads the destination MAC address, consults its table, and sends the frame out of the correct port only (\emph{unicast}). If the address is not yet known (\emph{unknown unicast}), the switch temporarily floods the frame to all ports except the incoming one, then learns the sender's location from the reply. Each port is its own collision domain (full-duplex operation eliminates collisions entirely).

\begin{propriete}[Switch vs Hub -- examinable comparison]
\begin{itemize}
\item \textbf{Hub:} Layer 1, one collision domain, half-duplex, broadcasts to all ports, 10/100 Mbit/s shared bandwidth.
\item \textbf{Switch:} Layer 2, one collision domain \emph{per port}, full-duplex, forwards/floods based on MAC table, dedicated bandwidth per port (100 Mbit/s, 1 Gbit/s, 10 Gbit/s).
\end{itemize}
\end{propriete}

\begin{attention}[Hub $\neq$ switch]
A very common mistake is to treat a hub and a switch as identical. They are not: \textbf{a hub broadcasts to all ports; a switch selects the destination port.} The switch therefore reduces collisions, increases effective speed and improves confidentiality.
\end{attention}

\textbf{Managed vs Unmanaged.} An \emph{unmanaged switch} works out of the box (plug-and-play). A \emph{managed switch} offers VLANs, port mirroring, QoS, SNMP, link aggregation (LACP), and can be configured via CLI/Web/SSH. School labs and enterprises use managed switches; homes and small cybercafés often use unmanaged.

The switch is the \emph{central device} of the modern star topology studied in the previous sections: every computer in the lab or office is wired back to it.

\courssub{The Router}

\begin{definition}[Router]
A \cle{router} is a Layer 3 device that \textbf{interconnects different networks} and routes data packets between them, choosing the best path towards the destination using \cle{IP addresses}. The router that links a local network to the outside world is called the \cle{default gateway}.
\end{definition}

\textbf{Intuition.} If the switch is the internal postman of one town (one network), the router is the sorting office that dispatches mail \emph{between} towns (between networks). When a student in Buea opens a website hosted abroad, the packet leaves the school LAN through the router (default gateway), which forwards it towards the Internet.

The \cle{home router} (often called ``box'' or ``gateway'') supplied by Internet providers (MTN, Orange, CAMTEL) is actually several devices in one box:
\begin{itemize}
\item a \textbf{router} (WAN port, NAT, firewall, DHCP server, DNS proxy);
\item a small \textbf{switch} (typically 4 LAN ports);
\item a \textbf{wireless access point} (Wi-Fi 5/6);
\item often a \textbf{modem} (ADSL, VDSL, or fibre ONT integrated or separate).
\end{itemize}

\begin{aretenir}[Router functions in a SOHO device]
\begin{enumerate}
\item \textbf{Inter-network routing:} forwards packets between LAN and WAN using IP addresses.
\item \textbf{NAT (Network Address Translation):} allows multiple private IPs (e.g. 192.168.1.x) to share one public IP.
\item \textbf{DHCP server:} automatically assigns IP configuration to LAN clients.
\item \textbf{Firewall / SPI:} blocks unsolicited incoming traffic from the Internet.
\item \textbf{Wireless AP:} provides Wi-Fi access (see below).
\end{enumerate}
\end{aretenir}

\courssub{The Modem}

\begin{definition}[Modem]
A \cle{modem} (MOdulator--DEModulator) is a device that converts the digital signals of a computer network into signals suitable for a transmission line (telephone line, coaxial cable, fibre) and back again. It connects a LAN to an Internet Service Provider (ISP).
\end{definition}

\begin{itemize}
\item \textbf{ADSL/VDSL modem:} uses the ordinary copper telephone line (PSTN); typical of older home connections; asymmetric speeds (download > upload).
\item \textbf{Cable modem (DOCSIS):} uses the coaxial TV cable network; higher speeds.
\item \textbf{Fibre modem / ONT (Optical Network Terminal):} converts light pulses from a fibre optic line (GPON, XGS-PON) into Ethernet signals; used with fibre-to-the-home offers in Douala, Yaound\'e and major towns.
\end{itemize}
In modern FTTH installations, the ISP provides an ONT (modem) and a separate router, or a combined gateway.

\courssub{Wireless Access Point, Repeater and Bridge}

\begin{itemize}
\item \textbf{Wireless Access Point (WAP / AP):} allows Wi-Fi devices (laptops, phones, tablets) to join a wired network. It is the wireless equivalent of a switch port. It bridges 802.11 (Wi-Fi) frames to 802.3 (Ethernet) frames. Multiple APs with the same SSID form an \emph{Extended Service Set (ESS)} enabling roaming.
\item \textbf{Repeater / Range extender:} receives a Wi-Fi signal and retransmits it to extend coverage, for example to carry Wi-Fi to the far end of a school building. It halves the available bandwidth (half-duplex relay). Modern \emph{mesh systems} use dedicated backhaul radios to avoid this penalty.
\item \textbf{Bridge:} connects two similar network segments (usually wired) so that they behave as one Layer 2 network, for example linking the wired LAN of two nearby classrooms or buildings via a wireless bridge (point-to-point link).
\end{itemize}

\courssub{Other Installation Hardware (Structured Cabling)}

\begin{itemize}
\item \textbf{Patch panel:} a panel of RJ45 sockets (usually 24 or 48 ports), mounted in a 19-inch cabinet/rack, where all the horizontal cables of a building terminate; short \emph{patch cords} then link it to the switch. Provides flexibility and protects solid-core cable from repeated plugging.
\item \textbf{Wall socket (keystone jack / faceplate):} the RJ45 outlet fixed in the wall or floor box, into which the user plugs a computer via a patch cord.
\item \textbf{Crimping tool:} the pliers used to fix an RJ45 connector onto the wires of a patch cable (stranded wire). \emph{Not} for punching down solid cable onto a patch panel or keystone (use a \emph{punch-down tool} / Krone tool).
\item \textbf{Cable tester:} a small device (main unit + remote) that checks whether all eight wires of a cable are correctly connected, detects opens, shorts, crossed pairs, and split pairs. Essential for practical exams.
\item \textbf{PoE (Power over Ethernet) injector / PoE switch:} delivers DC power (up to 90 W with 802.3bt) over the same Ethernet cable to devices like APs, IP cameras, VoIP phones --- no separate power adapter needed.
\end{itemize}

\begin{popfigure}
\begin{center}
\begin{tikzpicture}[font=\small, node distance=1.1cm]
\tikzset{dev/.style={draw=popdark, fill=popblueL, rounded corners, thick, minimum width=2.1cm, minimum height=0.9cm, align=center}}
\node[dev] (pc1) at (0,1.4) {PC 1};
\node[dev] (pc2) at (0,0) {PC 2};
\node[dev] (pc3) at (0,-1.4) {Laptop (Wi-Fi)};
\node[dev, fill=popgreenL, align=center] (sw) at (3.6,0){Switch\\ \scriptsize forwards frames by MAC};
\node[dev, fill=popgreenL, align=center] (ap) at (3.6,-2.6){Access Point\\ \scriptsize Wi-Fi $\leftrightarrow$ Ethernet bridge};
\node[dev, fill=popblueL, align=center] (rt) at (7.2,0){Router\\ \scriptsize Gateway, NAT, DHCP, FW};
\node[dev, align=center] (md) at (10.2,0){ONT / Modem\\ \scriptsize Fibre / ADSL / Cable};
\draw[thick, poppurple] (12.6,0) circle (1.05);
\node[poppurple, align=center] at (12.6,0) {Internet};
\draw[thick, popdark] (pc1.east) -- node[above,font=\scriptsize] {UTP} (sw.170);
\draw[thick, popdark] (pc2.east) -- node[above,font=\scriptsize] {UTP} (sw.west);
\draw[thick, popdark, dashed] (pc3.east) -- node[above,font=\scriptsize] {Wi-Fi} (ap.west);
\draw[thick, popdark] (ap.north) -- node[right,font=\scriptsize] {UTP} (sw.south);
\draw[thick, popdark] (sw.east) -- node[above,font=\scriptsize] {UTP} (rt.west);
\draw[thick, popdark] (rt.east) -- node[above,font=\scriptsize] {WAN} (md.west);
\draw[thick, popdark] (md.east) -- node[above,font=\scriptsize] {Fibre/Coax/Phone} (11.55,0);
\end{tikzpicture}
\end{center}
\legende{A small organisational LAN: PCs connect to a switch via UTP, the switch to a router (gateway), the router to an ONT/modem, and the modem to the ISP network. Wi-Fi clients join via an Access Point wired to the switch.}
\end{popfigure}

\begin{methode}[Distinguishing hub, switch and router in exam questions]
Ask one question: \emph{``Who receives the data?''}
\begin{enumerate}
\item If the device sends the data to \textbf{everybody} on the local network (Layer 1 broadcast) $\Rightarrow$ it is a \textbf{hub}.
\item If it sends the data only to the \textbf{one destination machine} inside the same network (Layer 2, MAC address) $\Rightarrow$ it is a \textbf{switch}.
\item If it sends the data \textbf{out of the local network} towards another network or the Internet (Layer 3, IP address) $\Rightarrow$ it is a \textbf{router}.
\end{enumerate}
\end{methode}

\begin{exoresolu}[Identifying devices in a school network]
The computer laboratory of GBHS Bamenda has 20 PCs cabled to a central device which sends each frame only to the intended PC. This device is connected to a second device which gives access to the Internet through the school's CAMTEL fibre line. (a) Name the two devices. (b) State the role of each. (c) Why is the first device preferred to a hub?
\tcblower
\textbf{(a)} The first device is a \textbf{switch}; the second is a \textbf{router} (acting as default gateway, with a fibre ONT towards CAMTEL).\\
\textbf{(b)} The switch connects the 20 PCs in a star topology and forwards each frame only to the port of the destination PC, using MAC addresses. The router interconnects the school LAN with the Internet and routes packets between the two networks using IP addresses; it also performs NAT so that all 20 PCs share the school's single public IP.\\
\textbf{(c)} A hub would broadcast every frame to all 20 PCs, causing collisions, wasting bandwidth and letting every PC see traffic meant for others. The switch selects only the destination port, so communication is faster, more efficient and more private.
\end{exoresolu}

\begin{experience}[Guided practical: Build and test a patch cable]
\textbf{Objective:} Terminate a Cat6 UTP cable with two RJ45 plugs (T568B) and verify continuity.
\textbf{Equipment:} Cat6 UTP cable (1 m), 2 RJ45 plugs, crimping tool, cable tester, wire stripper/cutter.
\textbf{Steps:}
\begin{enumerate}
\item Strip $\approx$ 25 mm of outer jacket; untwist pairs only as needed.
\item Arrange wires in T568B order: White-Orange, Orange, White-Green, Blue, White-Blue, Green, White-Brown, Brown.
\item Trim wires flush so $\approx$ 12 mm protrude; insert fully into RJ45 plug (each wire in its channel).
\item Crimp firmly; repeat for the other end (same order for straight-through).
\item Test: plug each end into the tester; all 8 LEDs should light in sequence. If any LED fails, re-crimp.
\end{enumerate}
\textbf{Exam tip:} In the practical, you may be asked to identify a faulty cable (split pair, reversed pair, open circuit) using the tester.
\end{experience}

\begin{exercice}[Roles of network equipment]
For each situation, name the most appropriate device and justify your choice:
\begin{enumerate}
\item Connecting 30 desktop computers of a cybercafé in Douala into one LAN.
\item Giving Wi-Fi access to students' phones in that LAN.
\item Linking the LAN to the Internet via a CAMTEL fibre subscription.
\item Extending Wi-Fi coverage to a room 80 m away from the main router.
\item Interconnecting the wired LANs of two buildings 200 m apart using a wireless link.
\end{enumerate}
\end{exercice}

\begin{corrige}[Exercise 1]
\begin{enumerate}
\item A \textbf{switch} (preferably managed, 48-port Gigabit): it connects the PCs in a star and forwards frames only to the destination port, providing dedicated bandwidth.
\item A \textbf{wireless access point} (PoE-powered, ceiling-mounted): it lets Wi-Fi devices join the wired LAN; multiple APs with same SSID allow roaming.
\item A \textbf{router} (with a separate \textbf{ONT} provided by CAMTEL): the ONT converts fibre to Ethernet; the router performs NAT, DHCP, firewall and connects the LAN to the ISP.
\item A \textbf{Wi-Fi repeater/extender} (or better, a second AP wired via UTP if distance $\le$ 100 m, or a mesh system): it regenerates the Wi-Fi signal to cover the distant room.
\item A \textbf{wireless bridge} (point-to-point link, e.g. 5 GHz directional antennas): it connects the two building switches transparently at Layer 2.
\end{enumerate}
\end{corrige}

\begin{aretenir}[Exam tip: one role + one example per device]
GCE questions very often ask: \emph{``State the role of device X in this diagram.''} For each device, memorise \textbf{one precise role and one example of use}:
\begin{itemize}
\item \textbf{NIC} --- connects a computer to the network and holds the MAC address (e.g. laptop Wi-Fi card).
\item \textbf{Switch} --- forwards frames to the destination port only using MAC addresses (e.g. school lab centre).
\item \textbf{Router} --- interconnects different networks, default gateway to the Internet, performs NAT (e.g. home gateway).
\item \textbf{Modem / ONT} --- converts signals for the ISP line (ADSL, cable, fibre) (e.g. CAMTEL ONT).
\item \textbf{Access Point} --- bridges Wi-Fi clients to the wired LAN (e.g. classroom ceiling AP).
\item \textbf{Hub} --- obsolete Layer 1 broadcaster (contrast with switch).
\item \textbf{Repeater/Extender} --- extends Wi-Fi coverage (half-duplex relay).
\item \textbf{Bridge} --- connects two LAN segments at Layer 2 (e.g. building-to-building wireless link).
\item \textbf{Patch panel / Keystone} --- structured cabling termination points.
\item \textbf{Crimping tool / Punch-down tool / Cable tester} --- installation and verification tools.
\end{itemize}
Never write that a hub and a switch do the same job.
\end{aretenir}

Now that we know the equipment, the next section will put it into practice: physically \emph{setting up a computer network}, from cabling the RJ45 connectors to configuring the machines.

\coursec{Setting Up a Computer Network}

In the previous section we studied the \cle{network equipment} --- cables, connectors, switches, routers, network interface cards --- and how to choose them. Owning good equipment, however, is not the same as having a working network. A box of switches and cables in the store of a school in Bamenda connects nobody. This section is the \emph{practical heart} of the chapter: we learn, step by step, how to move from an empty computer room to a fully working \cle{local area network} (LAN) --- planning, cabling, connecting, configuring, testing and documenting. This is exactly the kind of task an ICT technician performs when a school, a cybercaf\'e or an office in Douala asks for a network to be installed.

\courssub{Overview of the Setup Process}

Setting up a network is a project, and like every project it follows an ordered process. Skipping a step almost always means coming back later to fix avoidable mistakes.

\begin{methode}[The 6-step procedure for setting up a small LAN]
\begin{enumerate}
\item \textbf{Plan}: assess needs, survey the site, draw the network plan.
\item \textbf{Gather equipment}: buy or collect cables, connectors, switch, router, tools, according to the plan.
\item \textbf{Cable and crimp}: measure, cut and terminate the cables with RJ45 plugs; test each cable.
\item \textbf{Connect}: plug computers into the switch, the switch into the router; label everything.
\item \textbf{Configure}: assign IP addresses, computer names, workgroup; enable sharing.
\item \textbf{Test and document}: verify connectivity (link lights, ping), fix faults, write the installation report.
\end{enumerate}
\end{methode}

\begin{popfigure}
\begin{center}
\begin{tikzpicture}[
box/.style={draw=popblue, fill=popblueL, rounded corners=2pt, minimum width=2.6cm, minimum height=0.9cm, align=center, font=\small},
arr/.style={-{Stealth[length=2.5mm]}, thick, popdark}]
\node[box, align=center] (plan) at (0,0){\textbf{1. Plan}\\needs, survey, map};
\node[box, align=center] (gather) at (3.4,0){\textbf{2. Gather}\\equipment \& tools};
\node[box, align=center] (cable) at (6.8,0){\textbf{3. Cable}\\cut, crimp, test};
\node[box, align=center] (connect) at (6.8,-1.8){\textbf{4. Connect}\\PCs, switch, router};
\node[box, align=center] (config) at (3.4,-1.8){\textbf{5. Configure}\\IP, names, sharing};
\node[box, align=center] (test) at (0,-1.8){\textbf{6. Test}\\ping, document};
\draw[arr] (plan) -- (gather);
\draw[arr] (gather) -- (cable);
\draw[arr] (cable) -- (connect);
\draw[arr] (connect) -- (config);
\draw[arr] (config) -- (test);
\draw[arr, dashed, poporange] (test.west) .. controls (-1.6,-1.8) and (-1.6,0) .. (plan.west) node[midway, left, font=\scriptsize, text=poporange, align=center]{fault?\\rework};
\end{tikzpicture}
\end{center}
\legende{The network setup process. If testing fails, return to the step where the fault originated.}
\end{popfigure}

\courssub{Phase 1: Planning the Network}

\begin{definition}[Needs assessment]
The \cle{needs assessment} is the first stage of planning, in which we determine \emph{what the network must do}: the number of users and devices, the services required (file sharing, printing, Internet access), the expected growth, and the available budget.
\end{definition}

A good technician asks questions before buying anything. For a school computer laboratory the typical answers are: 15 student computers plus 1 teacher computer; shared printer; Internet access through a CAMTEL or MTN router; budget limited; possible extension to 30 computers next year.

The second planning activity is the \cle{site survey}: visiting the room, measuring distances, locating electrical sockets, noting where cables can run (trunking along walls, ceiling, under desks) and choosing a central, dry, ventilated position for the switch. From the survey we draw the \cle{network plan} --- a scaled sketch showing each computer, the switch, the router and the cable runs. This plan is our working document and later becomes part of the documentation.

\begin{exemplebox}[Planning a 15-computer school laboratory]
Room: $10\ \mathrm{m} \times 7\ \mathrm{m}$. Computers arranged in three rows of five along the walls. Switch placed on a wall shelf at the centre of the room to minimise cable lengths (star topology). Router near the switch, connected to the ISP line. Longest cable run estimated at $12\ \mathrm{m}$; we add $2\ \mathrm{m}$ slack per cable, so we prepare cables of about $14\ \mathrm{m}$ maximum.
\end{exemplebox}

\begin{popfigure}
\begin{center}
\begin{tikzpicture}[scale=0.62, every node/.style={font=\scriptsize}]
\draw[thick, popdark] (0,0) rectangle (12,8);
\node[below left, popmuted] at (0,8) {Lab $10\ \mathrm{m}\times 7\ \mathrm{m}$};
\foreach \x in {1,3,5,7,9} {
  \draw[fill=popblueL, draw=popblue] (\x,7.1) rectangle (\x+1.2,7.8);
  \draw[fill=popblueL, draw=popblue] (\x,0.2) rectangle (\x+1.2,0.9);
}
\foreach \y in {2.4,3.8,5.2} {
  \draw[fill=popblueL, draw=popblue] (0.2,\y) rectangle (0.9,\y+1.2);
  \draw[fill=popblueL, draw=popblue] (11.1,\y) rectangle (11.8,\y+1.2);
}
\draw[fill=popgreenL, draw=popgreen, thick] (10.6,4.6) rectangle (11.4,5.4);
\node at (11,5) {SW};
\draw[fill=popblueL, draw=popgold, thick] (10.6,3.4) rectangle (11.4,4.2);
\node at (11,3.8) {RTR};
\draw[poporange, dashed] (11,4.6) -- (11,4.2);
\foreach \x in {1.6,3.6,5.6,7.6,9.6} {
  \draw[popline, thin] (\x+0.6,7.1) -- (11,5.4);
  \draw[popline, thin] (\x+0.6,0.9) -- (11,4.6);
}
\foreach \y in {3,4.4,5.8} {
  \draw[popline, thin] (0.9,\y) -- (10.6,4.9);
  \draw[popline, thin] (11.1,\y) -- (11.4,5);
}
\draw[fill=popink] (5.4,3.6) rectangle (6.6,4.4);
\node[white] at (6,4) {Tchr};
\draw[popline, thin] (6.6,4) -- (10.6,4.8);
\end{tikzpicture}
\end{center}
\legende{Floor-plan sketch of the 15-computer laboratory: star topology, switch (SW) and router (RTR) on a central wall shelf, cable runs in trunking along the walls.}
\end{popfigure}

\courssub{Phase 2: Choosing Topology and Equipment}

From the previous sections, the \cle{star topology} is the natural choice for a wired lab: easy to extend, a single cable failure affects only one machine, and switches are affordable. The equipment list follows directly from the plan.

\begin{center}
\begin{tabularx}{\linewidth}{|l|c|X|}
\hline
\rowcolor{popblueL} \textbf{Item} & \textbf{Qty} & \textbf{Specification / remark} \\
\hline
Desktop computers with NIC & 16 & 15 students + 1 teacher \\
\hline
Switch & 1 & 24 ports (16 used, room to grow) \\
\hline
Router / modem & 1 & Provided by ISP (CAMTEL, MTN\ldots) \\
\hline
UTP cable Cat 5e/6 & 1 box & $305\ \mathrm{m}$ roll, solid or stranded \\
\hline
RJ45 connectors & 50 & More than 32 needed (spares for bad crimps) \\
\hline
Crimping tool, stripper & 1 set & For RJ45 \\
\hline
Cable tester & 1 & Checks continuity and wire order \\
\hline
Cable ties, trunking, labels & --- & Cable management and identification \\
\hline
\end{tabularx}
\end{center}

\courssub{Phase 3: Cabling and Crimping}

A UTP cable contains four twisted pairs, i.e.\ eight coloured wires. To connect it to a plug or socket, the wires must be inserted in a precise, standardised order. The standard used almost everywhere today is \cle{T568B}.

\begin{methode}[Crimping an RJ45 plug (T568B)]
\begin{enumerate}
\item Measure the cable run, add about $1\ \mathrm{m}$ of slack, and cut the cable cleanly.
\item Strip about $2\ \mathrm{cm}$ of the outer jacket with the stripper, without nicking the inner wires.
\item Untwist the pairs and arrange the eight wires flat, in the T568B order (pin 1 to pin 8): \textbf{white-orange, orange, white-green, blue, white-blue, green, white-brown, brown}.
\item Trim the wires straight so about $1.2\ \mathrm{cm}$ remains; keep the order tight.
\item Push the wires fully into the RJ45 plug (clip facing down, cable entering away from you); check through the plug that every wire reaches the end and the jacket enters the plug.
\item Insert the plug into the crimping tool and squeeze firmly once. Repeat for the other end, then test the cable with the cable tester: all eight lights must light in sequence on both ends.
\end{enumerate}
\end{methode}

\begin{attention}[Straight-through vs crossover cable]
A \cle{straight-through} cable (both ends T568B) connects \emph{unlike} devices: computer to switch, switch to router. A \cle{crossover} cable (one end T568A, the other T568B) connects \emph{like} devices: computer to computer, or switch to switch without an uplink port. Many modern devices auto-sense (Auto-MDIX) and accept either cable, but in the GCE examination you must state the rule correctly: \textbf{unlike devices $\to$ straight-through; like devices $\to$ crossover}.
\end{attention}

\begin{attention}[Electrical safety]
Never crimp, punch down or dress cables on equipment that is powered on. Switch off and unplug devices before working inside them, keep cables away from mains power lines to avoid interference and shock risks, and never work with wet hands. A crimping tool and a cable knife are sharp: cut away from your body.
\end{attention}

\courssub{Phase 4: Physical Installation}

With tested cables ready, installation is systematic:

\begin{enumerate}
\item Mount the switch on its shelf and the router beside it; connect their power adapters but keep them off until cabling is done.
\item Run each cable through the trunking from a computer's NIC to the switch; plug it into a switch port.
\item \textbf{Label both ends of every cable} and the corresponding switch port (e.g.\ PC01 $\to$ port 1). Labelling saves hours during troubleshooting.
\item Connect the switch to the router's LAN port with a straight-through cable; connect the router to the ISP line (telephone line, fibre box or 4G modem).
\item Power on the router, then the switch, then the computers. Check that the \cle{link light} on each switch port and each NIC is on (usually green) or blinking: this confirms the physical connection.
\end{enumerate}

\courssub{Phase 5: Basic Configuration}

Physical connection is not enough: each machine needs a logical identity on the network --- an \cle{IP address}. There are two ways to obtain one.

\begin{definition}[DHCP]
\cle{DHCP} (Dynamic Host Configuration Protocol) is a network service that \emph{automatically assigns} IP addresses (together with the subnet mask, default gateway and DNS server) to devices when they join a network. The alternative is \cle{static addressing}, where the administrator types the address manually on each machine.
\end{definition}

For our lab, the router usually acts as the DHCP server, so the simplest configuration is to set every computer to ``Obtain an IP address automatically''. Static addresses are preferred for machines that must always be found at the same address (the teacher's PC sharing the printer, a server). A typical private addressing scheme is $192.168.1.x$ with mask $255.255.255.0$: the router takes $192.168.1.1$, the teacher's PC $192.168.1.10$ (static), and the students receive $192.168.1.100$ and above from DHCP.

Each computer also receives a unique \cle{computer name} (PC01, PC02, \ldots, TEACHER) and joins the same \cle{workgroup} (e.g.\ LABGROUP) so that machines can see one another. Finally, on the teacher's machine we enable \cle{file and printer sharing} and share the printer and a folder ``Course\_Documents'' so students can retrieve handouts and print.

\courssub{Phase 6: Testing and Troubleshooting}

The reference tool for testing connectivity is the \cle{ping} command: it sends small echo-request packets to a target address and waits for replies. Open a command prompt and type, for example, \texttt{ping 192.168.1.10} from PC01.

\begin{popfigure}
\begin{center}
\begin{tikzpicture}[
pc/.style={draw=popblue, fill=popblueL, rounded corners=2pt, minimum width=1.7cm, minimum height=1.1cm, align=center, font=\small},
sw/.style={draw=popgreen, fill=popgreenL, rounded corners=2pt, minimum width=2.6cm, minimum height=0.8cm, font=\small},
msg/.style={-{Stealth[length=2.2mm]}, thick, poporange},
font=\scriptsize]
\node[pc, align=center] (a) at (0,0){PC01\\192.168.1.101};
\node[sw] (s) at (4.2,0) {Switch};
\node[pc, align=center] (b) at (8.4,0){TEACHER\\192.168.1.10};
\draw[thick, popdark] (a) -- (s);
\draw[thick, popdark] (s) -- (b);
\draw[msg] (1.1,0.9) -- (7.3,0.9) node[midway, above]{echo request: ping 192.168.1.10};
\draw[msg] (7.3,-0.9) -- (1.1,-0.9) node[midway, below]{echo reply: Reply from 192.168.1.10: time$<$1ms};
\node[align=left, anchor=west, text=popdark] at (0,-1.9) {Expected result: 4 replies received, 0\% loss.};
\node[align=left, anchor=west, text=popink] at (0,-2.5) {``Request timed out'' or ``Destination unreachable'' $\Rightarrow$ fault to locate.};
\end{tikzpicture}
\end{center}
\legende{Ping test between PC01 and the teacher's computer through the switch, with the expected replies.}
\end{popfigure}

\begin{methode}[Locating a fault systematically]
\begin{enumerate}
\item Check the \textbf{link lights} on the NIC and the switch port. No light $\Rightarrow$ physical problem (cable, crimping, port).
\item Swap the cable with a known-good one, or re-test it with the cable tester. A bad crimp (wire in wrong order, wire not reaching the end) is the most common fault: cut the plug and crimp again.
\item Check the \textbf{cable type}: straight-through for PC-to-switch.
\item \texttt{ping 127.0.0.1} tests the machine's own TCP/IP stack; then ping the gateway, then another PC.
\item Check the IP configuration: two machines with the \textbf{same static IP} cause an \cle{IP conflict} --- both may lose connectivity. Give a different address or use DHCP.
\end{enumerate}
\end{methode}

\begin{exoresolu}[Diagnosing a dead connection in the lab]
After installation, PC07 cannot reach the network. Its switch port link light is off, while all other ports are lit. The cable tester shows that wires 3 and 6 light up swapped. What is the fault and the remedy?
\tcblower
The link light being off indicates a \emph{physical} fault, not a configuration problem. The tester shows wires 3 and 6 (white-green and orange in a wrong crimp) are exchanged at one end: the RJ45 plug was crimped with two wires in the wrong order, so the pairs are broken and no link can be established. \textbf{Remedy}: cut off the faulty plug, re-arrange the eight wires in the correct T568B order (white-orange, orange, white-green, blue, white-blue, green, white-brown, brown), crimp a new plug, re-test with the cable tester until all eight wires light in sequence, then reconnect and verify with \texttt{ping 192.168.1.1}.
\end{exoresolu}

\courssub{Safety, Cable Management and Documentation}

A professional installation is recognised by its finishing. Cables are grouped with \cle{cable ties} and hidden in \cle{trunking}; nothing crosses walkways where people can trip; network cables are kept away from mains cables; the switch and router sit in a ventilated, dust-free position, ideally protected by a surge protector or UPS against the power cuts and surges common in our environment. Finally, the technician writes the \cle{documentation}: the network plan, the list of equipment, the IP addressing table, the labelling scheme and the test results. This document is indispensable for maintenance and for any future extension of the network.

\begin{aretenir}[Key points]
\begin{itemize}
\item Setup follows six steps: plan, gather, cable, connect, configure, test.
\item T568B wire order: white-orange, orange, white-green, blue, white-blue, green, white-brown, brown.
\item Straight-through cable for unlike devices; crossover for like devices.
\item DHCP assigns addresses automatically; static addresses for fixed machines.
\item Link lights and \texttt{ping} are the first troubleshooting tools; bad crimping and IP conflicts are the classic faults.
\item Label everything, manage cables safely, and document the installation.
\end{itemize}
\end{aretenir}

\begin{exercice}[Setting up a cybercaf\'e network]
A cybercaf\'e in Buea has 8 customer computers, 1 cashier computer, one shared printer and a 4G router. (a) Draw the network plan and state the topology. (b) Give the equipment list. (c) Propose an IP addressing scheme using DHCP with a static address for the cashier's PC. (d) After setup, PC05 replies to \texttt{ping 127.0.0.1} but not to \texttt{ping 192.168.1.1}. Suggest two possible causes and their remedies.
\end{exercice}

\begin{corrige}[Exercise 1]
(a) Star topology: all 9 computers connect to a central switch; the switch connects to the 4G router; the printer is shared from the cashier's PC (or a network printer on the switch). (b) Equipment: 9 computers with NICs, one 16-port switch, one 4G router, UTP Cat 5e/6 cable, at least 20 RJ45 plugs, crimping tool, cable tester, trunking, labels. (c) Network $192.168.1.0/24$: router $192.168.1.1$ as DHCP server (range $192.168.1.100$--$192.168.1.150$), cashier PC static $192.168.1.10$, mask $255.255.255.0$, gateway $192.168.1.1$. (d) Since the loopback ping works, TCP/IP is fine. Possible causes: (i) faulty cable or bad crimp between PC05 and the switch --- test the cable and re-crimp; (ii) wrong IP configuration (wrong gateway or IP conflict) --- set the PC to obtain an address automatically via DHCP, or correct the static settings, then ping the gateway again.
\end{corrige}

\coursec{Selecting Equipment for Organisational Networks}

In the previous section we learnt how to physically set up a small computer network: cabling, connecting devices, configuring IP addresses and testing connectivity. But before any cable is pulled, a harder question must be answered: \emph{which equipment should we buy, and how much of it?} A school bursar, a cybercafé owner in Bamenda or a hospital administrator in Yaoundé does not buy a switch ``because it looks nice''. Equipment is selected to match \cle{needs}, \cle{growth} and \cle{budget}. This section teaches you the professional reasoning behind equipment selection, ending with two complete case studies and an integration activity in which you produce a full \cle{network dossier} — exactly the kind of scenario question the GCE Advanced Level examiners love.

\courssub{Needs Analysis: The Starting Point}

Imagine a tailor sewing a suit without measuring the customer. The result may be beautiful, but it will not fit. Selecting network equipment without analysing needs is the same mistake. A \cle{needs analysis} is a systematic study of what the organisation actually requires from its network.

\begin{definition}[Needs analysis]
A \cle{needs analysis} is the systematic study, carried out before any purchase, of the users, services, traffic, growth prospects and physical constraints of an organisation, in order to determine the network equipment that must be acquired.
\end{definition}

\begin{methode}[Carrying out a needs analysis]
\begin{enumerate}
\item \textbf{Count and classify the users.} How many people will use the network? Are they students, secretaries, administrators, customers? Each class of user has different demands.
\item \textbf{List the services required.} File sharing, Internet access, shared printing, video streaming, VoIP telephony, a local server (e.g.\ school management software), video surveillance.
\item \textbf{Estimate the traffic.} Ten users typing documents generate little traffic; twenty-five users streaming video generate a lot. Traffic determines the required \cle{bandwidth}.
\item \textbf{Forecast growth.} Will the organisation add users, computers or services in the next 3 to 5 years? A college that plans to open a computer science laboratory next year must plan for it today.
\item \textbf{Note the physical constraints.} Number of rooms and floors, distances between buildings, availability of electrical power, security of the equipment room.
\end{enumerate}
\end{methode}

\begin{exemplebox}[Needs analysis for a small secondary school]
A school in Buea has: 1 administrative office (4 computers, 1 printer), 1 staff room (6 laptops), and 1 computer laboratory (30 student PCs + 1 teacher PC). Services needed: Internet sharing, file/print sharing, and a school-management database. Growth: a second laboratory is planned within 3 years. Conclusion: the network must serve about 42 stations today, scale to about 75, and carry moderate traffic (web + database, little video).
\end{exemplebox}

\begin{attention}[Common mistake]
Analysing only \emph{today's} situation. Equipment bought for exactly today's needs becomes undersized within two years, forcing a costly complete replacement. Always include a growth forecast in the needs analysis.
\end{attention}

\courssub{Sizing the Equipment}

Once needs are known, each piece of equipment is \emph{sized} — its capacity is chosen to match the demand.

\textbf{Switches: port count and speed.}\par
The two key figures are the \cle{number of ports} and the \cle{speed per port} (10/100/1000~Mb/s, i.e.\ Fast Ethernet or Gigabit Ethernet). Every wired station, printer, server and access point consumes one port; inter-switch links (uplinks) consume more.

\begin{propriete}[The 20\% spare-port rule]
A switch (or a set of switches) should provide \textbf{at least 20\% more ports than the number of devices connected today}, to allow future growth without replacing equipment. (This rule is a design guideline, not a theorem to prove; you must be able to \emph{apply} it in calculations.)
\end{propriete}

\begin{exemplebox}[Applying the 20\% rule]
A laboratory has 30 student PCs, 1 teacher PC, 1 printer and 1 server: 33 devices, plus 1 uplink port to the rest of the network $= 34$ ports used. Adding 20\%: $34 \times 1.2 = 40.8$, so at least 41 ports are required. Two 24-port switches ($48$ ports) satisfy the rule; a single 48-port switch also works but offers no redundancy if it fails.
\end{exemplebox}

\textbf{Router capacity.}\par
The router must handle the organisation's Internet traffic and the number of simultaneous users. For a small office, a combined router/modem from the ISP (CAMTEL, MTN, Orange) may suffice; a campus needs a router with higher throughput, multiple interfaces and possibly firewall and VPN features.

\textbf{Wireless coverage.}\par
Wireless access points (APs) have a limited indoor range (typically 20--50~m through walls). Sizing wireless means counting how many APs are needed to cover all rooms, and checking that each AP can serve the expected number of simultaneous clients.

\begin{attention}[Common mistake]
Buying exactly the number of switch ports needed today. With zero spare ports, adding a single new computer forces the purchase of a whole new switch, plus re-cabling. The 20\% rule exists precisely to avoid this trap.
\end{attention}

\courssub{Budgeting and Total Cost of Ownership}

A network project has four cost families: \textbf{cabling} (cables, conduits, sockets, patch panels), \textbf{devices} (switches, router, APs, server, UPS), \textbf{labour} (technicians, configuration), and \textbf{operation} (electricity, ISP subscription, maintenance, replacements). The sum of purchase cost and operating cost over the equipment's lifetime is called the \cle{total cost of ownership} (TCO). A cheap switch that fails every year may cost more in TCO than a reliable branded one. The art of budgeting is \emph{balancing performance against cost}: Gigabit everywhere is wasteful for a typing pool, but Fast Ethernet is a bottleneck for a video-editing room.

\courssub{Structured Cabling for Multi-Room and Multi-Floor Organisations}

When a network spans several rooms or floors, cables cannot be run randomly from desk to desk. The professional solution is \cle{structured cabling}.

\begin{definition}[Structured cabling]
\cle{Structured cabling} is a standardised way of organising a building's network cabling into two subsystems: the \cle{backbone} (vertical cabling linking equipment rooms and floors, often optical fibre) and the \cle{horizontal cabling} (links from each floor's distribution switch to the wall sockets of individual rooms). Connections are made through \cle{patch panels} in an \cle{equipment room}, so any socket can be connected to any switch port simply by moving a patch cord.
\end{definition}

The advantages are enormous: faults are located quickly, moving a user to another office requires no new cable, and the backbone can be upgraded (e.g.\ to faster fibre) without touching the horizontal cabling.

\begin{popfigure}
\begin{tikzpicture}[scale=0.95, every node/.style={font=\small}]
% Equipment room
\draw[fill=popblueL, draw=popblue, thick, rounded corners] (0,0) rectangle (3.4,1.6);
\node[align=center] at (1.7,0.8) {\textbf{Equipment room}\\ router + main switch\\ patch panels};
% Floor switches
\draw[fill=popgreenL, draw=popgreen, thick, rounded corners] (6,2.6) rectangle (9.4,3.8);
\node[align=center] at (7.7,3.2) {Floor 2 switch\\ (distribution)};
\draw[fill=popgreenL, draw=popgreen, thick, rounded corners] (6,0.2) rectangle (9.4,1.4);
\node[align=center] at (7.7,0.8) {Floor 1 switch\\ (distribution)};
% Backbone
\draw[very thick, poporange] (3.4,1.1) -- (6,0.8);
\draw[very thick, poporange] (3.4,1.3) -- (4.6,2.9) -- (6,3.2);
\node[poporange, align=center] at (4.4,2.0) {backbone\\ (fibre)};
% Wall sockets floor 1
\foreach \x/\lab in {10.6/{R1}, 12.2/{R2}, 13.8/{R3}}{
  \draw[fill=white, draw=popdark] (\x-0.35,0.55) rectangle (\x+0.35,1.05);
  \node at (\x,0.8) {\lab};
  \draw[popline] (9.4,0.8) -- (\x-0.35,0.8);
}
% Wall sockets floor 2
\foreach \x/\lab in {10.6/{R4}, 12.2/{R5}, 13.8/{R6}}{
  \draw[fill=white, draw=popdark] (\x-0.35,2.95) rectangle (\x+0.35,3.45);
  \node at (\x,3.2) {\lab};
  \draw[popline] (9.4,3.2) -- (\x-0.35,3.2);
}
\node[align=center, font=\footnotesize\itshape] at (12.2,0.1) {horizontal cabling to wall sockets};
\end{tikzpicture}
\legende{Structured cabling: equipment room, fibre backbone, floor distribution switches and horizontal cabling to wall sockets.}
\end{popfigure}

\courssub{Case Study 1: A 25-Machine Cybercafé in Bamenda}

\begin{experience}[Equipping a cybercafé]
\textbf{Needs analysis.} 25 customer PCs + 1 cashier PC + 1 shared printer + 1 server (billing software) $= 27$ wired devices. Services: Internet browsing, printing, light video. Growth: up to 30 PCs within 2 years.

\textbf{Topology.} A single room $\Rightarrow$ \cle{star topology} around central switches (simple, cheap, one failure affects one machine).

\textbf{Sizing.} Ports needed today: 27 devices + 1 uplink between switches + 1 port to the router $= 29$. With 20\% spare: $29 \times 1.2 = 34.8$, so at least 35 ports $\Rightarrow$ two 24-port Gigabit switches (48 ports). One router with firewall; one wireless AP for customers' phones; a UPS for the server and cashier.
\end{experience}

\begin{center}
\begin{tabularx}{\linewidth}{|l|l|X|}\hline
\rowcolor{popblueL} \textbf{Item} & \textbf{Qty} & \textbf{Justification} \\ \hline
24-port Gigabit switch & 2 & 48 ports $\geq$ 35 required (20\% rule) \\ \hline
Router with firewall & 1 & Shares the ISP link, protects the LAN \\ \hline
Wireless access point & 1 & Wi-Fi for customers' phones \\ \hline
Server + billing software & 1 & Controls sessions and payments \\ \hline
Cat6 cable (boxes) + sockets & 3 & Star cabling to 27 points \\ \hline
UPS & 2 & Protects server/cashier from power cuts \\ \hline
Labour (cabling + configuration) & --- & Certified technician, testing \\ \hline
\end{tabularx}
\end{center}

The estimated budget is then built by assigning a unit price to each line (obtained from local suppliers, e.g.\ in Bamenda's Commercial Avenue) and adding 10\% contingency for unforeseen expenses.

\courssub{Case Study 2: Interconnecting Three Buildings of a College Campus}

\begin{experience}[Campus network with fibre backbone]
\textbf{Situation.} A college has three buildings 80--150~m apart: Administration (10 stations), Classroom block (2 laboratories of 30 PCs), Library (15 stations + Wi-Fi).

\textbf{Choice of topology.} Copper Ethernet is limited to about 100~m per segment, so inter-building links use \cle{optical fibre}. A main switch in the Administration equipment room connects to a distribution switch in each building: physically a \cle{tree (hierarchical) topology}, each building internally wired as a star.

\textbf{Sizing.} Administration: $10 + 1$ uplink $= 11$ ports used; $11 \times 1.2 = 13.2$, so at least 14 ports $\Rightarrow$ one 16-port switch. Classroom block: $60 + 2$ uplinks $= 62$ ports used; $62 \times 1.2 = 74.4$, so at least 75 ports $\Rightarrow$ four 24-port switches (96 ports). Library: $15 + 1$ AP $+ 1$ uplink $= 17$ ports used; $17 \times 1.2 = 20.4$, so at least 21 ports $\Rightarrow$ one 24-port switch.
\end{experience}

\begin{popfigure}
\begin{tikzpicture}[scale=0.9, every node/.style={font=\small}]
% Main building
\draw[fill=popblueL, draw=popblue, thick, rounded corners] (4.6,3.6) rectangle (8.0,5.2);
\node[align=center] at (6.3,4.4) {\textbf{Administration}\\ main switch + router\\ 16-port switch};
% Building 2
\draw[fill=popgreenL, draw=popgreen, thick, rounded corners] (0,0) rectangle (3.6,1.6);
\node[align=center] at (1.8,0.8) {\textbf{Classroom block}\\ 4 $\times$ 24-port switches\\ 2 labs (60 PCs)};
% Building 3
\draw[fill=popgreenL, draw=popgreen, thick, rounded corners] (9.0,0) rectangle (12.6,1.6);
\node[align=center] at (10.8,0.8) {\textbf{Library}\\ 24-port switch\\ 15 PCs + Wi-Fi AP};
% Fibre links
\draw[very thick, poporange] (5.3,3.6) -- (2.6,1.6);
\draw[very thick, poporange] (7.3,3.6) -- (10.0,1.6);
\node[poporange, rotate=52] at (3.4,2.85) {fibre};
\node[poporange, rotate=-52] at (9.2,2.85) {fibre};
% Internet cloud (drawn as an ellipse to avoid extra libraries)
\draw[thick, poppurple, fill=white] (6.3,6.5) ellipse (1.5 and 0.55);
\node[poppurple] at (6.3,6.5) {Internet};
\draw[thick, poppurple] (6.3,5.2) -- (6.3,5.95);
% Internal stars hint
\foreach \x in {0.5,1.2,1.9,2.6,3.3}{\draw[popline] (1.8,0.35) -- (\x,-0.5);}
\node[font=\footnotesize\itshape] at (1.8,-0.75) {star inside each lab};
\foreach \x in {9.5,10.2,10.9,11.6,12.3}{\draw[popline] (10.8,0.35) -- (\x,-0.5);}
\node[font=\footnotesize\itshape] at (10.8,-0.75) {star in the library};
\end{tikzpicture}
\legende{Three-building campus: tree topology with an optical-fibre backbone and a star inside each building.}
\end{popfigure}

\courssub{The 5-Step Equipment Selection Procedure}

\begin{methode}[Selecting equipment for any organisation]
\begin{enumerate}
\item \textbf{Analyse needs:} users, services, traffic, growth, building layout.
\item \textbf{Choose the topology:} star for one room, tree with fibre backbone for several buildings.
\item \textbf{Size the equipment:} switch ports (apply the 20\% rule) and speed, router capacity, wireless coverage.
\item \textbf{Budget:} cables, devices, labour, operation; compare performance vs cost using TCO.
\item \textbf{Justify every choice:} each line of the equipment list must be defended by a figure from steps 1--3.
\end{enumerate}
\end{methode}

\begin{exoresolu}[Sizing the switches of an accounting firm]
An accounting firm in Douala occupies one floor with 18 workstations, 2 network printers and 1 server, all wired, plus 1 wireless access point. The firm expects to hire 6 more staff within 3 years. Determine the number and size of switches required, justifying your answer.
\tcblower
\textbf{Step 1 — count ports used today:} $18 + 2 + 1 + 1 = 22$ devices, plus 1 uplink port to the router $= 23$ ports.

\textbf{Step 2 — apply the 20\% rule:} $23 \times 1.2 = 27.6$, so at least 28 ports are required. The planned growth (6 staff) is covered: $23 + 6 = 29 \leq 48$.

\textbf{Step 3 — choose:} one 24-port switch is insufficient ($24 < 28$). Two 24-port Gigabit switches give 48 ports, satisfying the rule and the growth forecast, and providing partial redundancy. Gigabit is justified because accountants exchange large files with the server.

\textbf{Conclusion:} 2 switches of 24 Gigabit ports each.
\end{exoresolu}

\begin{exercice}[A hospital network]
A district hospital has: Reception (4 PCs, 1 printer), Consultation wing (8 PCs), Laboratory (6 PCs, 2 printers), all in one building, plus a pharmacy in a separate block 120~m away (3 PCs, 1 printer). Internet access and a patient-record server are required. (a) Perform the needs analysis. (b) Choose a topology and cabling for the inter-building link, justifying your choice. (c) Size the switches using the 20\% rule. (d) Sketch the network diagram.
\end{exercice}

\begin{corrige}[Exercise: a hospital network]
(a) Main building: $5 + 8 + 8 = 21$ devices; pharmacy: 4 devices; plus 1 server and 1 AP $\approx 27$ devices in total. Services: Internet, shared patient records, printing; growth moderate; two buildings 120~m apart.

(b) Tree topology: a star in each block, optical fibre between them (copper Ethernet is limited to about 100~m per segment; 120~m exceeds it, and fibre is immune to electrical interference from medical equipment).

(c) Main building: $21$ devices $+ 1$ server $+ 1$ AP $+ 1$ uplink $= 24$ ports used; $24 \times 1.2 = 28.8$, so at least 29 ports $\Rightarrow$ two 16-port switches (32 ports) or one 24-port + one 8-port switch. Pharmacy: $4 + 1$ uplink $= 5$ ports used; $5 \times 1.2 = 6$, so at least 6 ports $\Rightarrow$ one 8-port switch.

(d) Diagram: router + main switch at reception, fibre link to the pharmacy's 8-port switch, stars to each room's wall sockets.
\end{corrige}

\courssub{Integration: Producing a Complete Network Dossier}

The final deliverable of any network project — and of any GCE scenario question — is a \cle{network dossier} containing four elements:

\begin{enumerate}
\item \textbf{A labelled diagram} showing the topology, buildings, switches, router, servers and cabling types.
\item \textbf{An equipment list in a table}, each line justified by a calculation (port count, distance, bandwidth).
\item \textbf{A setup plan}: order of work (cabling $\to$ equipment room $\to$ device configuration $\to$ IP addressing $\to$ tests).
\item \textbf{A test checklist}: link lights on every port, \texttt{ping} between stations and to the router, Internet access, file and printer sharing, wireless coverage in every room, and a failover check (disconnect one cable and confirm only one station is affected).
\end{enumerate}

\begin{aretenir}[Exam tip]
In scenario questions, always structure your answer as: \textbf{diagram + equipment table + justification of each choice}. Examiners award marks for the \emph{reasoning} (needs analysis, 20\% rule, distance limits), not for brand names or prices.
\end{aretenir}

\begin{aretenir}[Key points of the section]
\begin{itemize}
\item Equipment selection starts from a \cle{needs analysis}: users, services, traffic, growth, building layout.
\item Size switches by port count (with \textbf{at least 20\% spare}) and speed; size routers by throughput; size wireless by coverage.
\item Budget with \cle{total cost of ownership}, balancing performance against cost.
\item Multi-room organisations use \cle{structured cabling}: backbone + horizontal cabling + patch panels + equipment room.
\item Inter-building links beyond about 100~m require \cle{optical fibre} and a \cle{tree topology}.
\item The professional deliverable is a complete \cle{network dossier}: diagram, justified equipment list, setup plan and test checklist.
\end{itemize}
\end{aretenir}

\newpage

\coursec{Integration activity}

\coursec{Integration Activity: Designing and Implementing a School Network}

\courssub{Aims of the Activity}

This integration activity mobilises, in one realistic project, all the competences developed in this chapter:

\begin{itemize}
\item Identify and compare \cle{physical topologies} (bus, star, ring, mesh, tree) and justify the choice of a topology for a given organisation;
\item Draw a clear, labelled \cle{network diagram} from a site plan;
\item Identify and select \cle{network equipment} (switches, router, modem, access point, cables, connectors) according to needs and budget;
\item Prepare and test a real \cle{straight-through RJ45 cable} (crimping, wire order T568B, continuity test);
\item Configure \cle{IP addresses} on hosts and verify connectivity with \texttt{ping};
\item Work in a team, manage a budget in FCFA, and defend a technical dossier orally before a decision-making audience.
\end{itemize}

\begin{aretenir}[Success Criteria]
Your dossier is successful if: the topology is justified (not just named), the diagram is complete and labelled, every equipment item is quantified and priced, the total stays within the budget, at least one crimped cable passes the tester, and two configured PCs successfully ping each other.
\end{aretenir}

\begin{definition}[Physical Network Topology]
The \cle{physical topology} of a network defines the geometric arrangement of the cables, devices and interconnections that make up the local area network (LAN). The five classic topologies are \textbf{bus}, \textbf{star}, \textbf{ring}, \textbf{mesh} and \textbf{tree} (hierarchical star). Each has distinct cost, reliability, scalability and maintenance characteristics.
\end{definition}

\begin{propriete}[Comparison Criteria for Topology Selection]
When choosing a topology for an organisation, evaluate:
\begin{enumerate}
\item \textbf{Cost:} cable length, number of ports, active equipment.
\item \textbf{Reliability / Fault tolerance:} impact of a single cable or port failure.
\item \textbf{Ease of maintenance / troubleshooting:} isolation of faults, simplicity of moves/adds/changes.
\item \textbf{Scalability:} ease of adding new hosts or segments.
\item \textbf{Performance:} collision domains, bandwidth sharing, latency.
\end{enumerate}
\end{propriete}

\courssub{Situation: Networking Garrisson Secondary School}

\textbf{Garrisson Secondary School} is a growing Anglophone secondary school in the North West Region of Cameroon. The school has two buildings separated by a courtyard:

\begin{itemize}
\item \textbf{Building A} contains Computer Lab 1 (20 PCs), Computer Lab 2 (20 PCs) and the Library (3 PCs used for cataloguing and research).
\item \textbf{Building B}, located \SI{60}{m} away from Building A, contains the Administration Office (5 PCs: Principal, Vice-Principal, Bursar and two Secretaries).
\end{itemize}

The Principal has just obtained an Internet subscription from a local ISP (a CAMTEL fibre offer, or a 4G router from MTN/Orange where fibre is unavailable). She wants:

\begin{enumerate}
\item All 48 computers to share files, a central printer in the Administration Office and the single Internet connection;
\item A Wi-Fi hotspot in the Library for students' laptops and phones;
\item A network that is easy to maintain and easy to extend (a third lab is planned next year);
\item The whole project to cost at most \textbf{1\,500\,000 FCFA}, including cabling and equipment.
\end{enumerate}

The Principal has contacted your class. Working in teams of 3 to 5 students, you must produce a complete technical dossier, carry out a partial practical implementation in the school lab, and defend your proposal before the class, which plays the role of the school administration.

\begin{center}
\begin{tabularx}{\linewidth}{|l|X|X|}\hline
\rowcolor{popblueL} \textbf{Data} & \textbf{Building A} & \textbf{Building B} \\ \hline
Hosts & 20 + 20 + 3 = 43 PCs & 5 PCs + 1 shared printer \\ \hline
Distance to other building & \multicolumn{2}{c|}{\SI{60}{m} (aerial or buried cable across the courtyard)} \\ \hline
Special need & Wi-Fi hotspot in the Library & Shared printer for all users \\ \hline
Budget & \multicolumn{2}{c|}{1\,500\,000 FCFA maximum} \\ \hline
\end{tabularx}
\end{center}

\begin{savaistu}[Cameroon Context]
In Cameroon, most schools and SMEs rely on \textbf{MTN} or \textbf{Orange} for 4G/5G Internet backup, while \textbf{CAMTEL} provides fibre backbone in urban areas. The \textbf{ANTIC} (Agence Nationale des Technologies de l'Information et de la Communication) regulates cybersecurity and promotes ICT in education. A typical school LAN uses UTP Cat6 copper for indoor runs (\le \SI{100}{m}) and fibre for inter-building links to avoid lightning damage — a real concern during the rainy season in the North West Region.
\end{savaistu}

\courssub{Tasks}

\textbf{Task 1 — Site Plan and Topology Choice.}\par
Draw a simple site plan of the two buildings. Choose a physical topology for each building and for the link between them. Justify your choice by comparing at least two topologies on the criteria of cost, reliability, ease of maintenance and scalability. Explain why a bus or a ring would (or would not) be appropriate here.

\textbf{Task 2 — Labelled Network Diagram.}\par
Produce the complete logical diagram of the network showing: the ISP modem, the router, the switches (with the number of ports used on each), the access point, the inter-building link and the hosts (you may represent groups of PCs symbolically, indicating the quantity). Label every device and every cable type.

\textbf{Task 3 — Equipment List and Budget.}\par
Establish a table of all equipment to buy: switches (state the port count of each and justify: e.g.\ a 24-port switch for a 20-PC lab leaves room for growth), router, modem, wireless access point, UTP cable (estimate the length: allow runs of \SI{30}{m} average per PC inside a lab, plus the \SI{60}{m} inter-building link — should you use UTP or fibre for this link? Justify), RJ45 connectors, patch panels or trunking, and a network printer if needed. Use realistic local prices (survey a shop or an online seller in your town) and check that the total does not exceed 1\,500\,000 FCFA.

\textbf{Task 4 — Practical Cabling.}\par
In the lab, crimp at least one straight-through RJ45 cable following the T568B wire order. Test it with a cable tester. Record the wire order you used and the test result in your dossier.

\begin{methode}[Crimping a Straight-Through RJ45 Cable (T568B)]
\begin{enumerate}
\item Cut the UTP cable to the required length (+ \SI{10}{cm} margin).
\item Strip \SI{2.5}{cm} of the outer jacket without nicking the inner wires.
\item Untwist the four pairs and arrange them in the T568B order:\par
\textbf{1. White-Orange, 2. Orange, 3. White-Green, 4. Blue, 5. White-Blue, 6. Green, 7. White-Brown, 8. Brown.}
\item Flatten and align the wires; trim them to \SI{1.3}{cm} from the jacket.
\item Insert the wires fully into the RJ45 plug (tab facing down, pins 1–8 left to right).
\item Crimp firmly with a ratchet crimping tool.
\item Repeat on the other end (same T568B order for straight-through).
\item Test with a cable tester: all 8 LEDs must light in sequence on both master and remote units.
\end{enumerate}
\end{methode}

\textbf{Task 5 — Configuration and Connectivity Test.}\par
Propose an IP addressing plan for the school (network address, subnet mask, ranges per room, gateway address). Then, on two real PCs in the lab, configure static IP addresses from your plan and verify connectivity with \texttt{ping}. Note the commands used and the results.

\begin{methode}[Static IP Configuration and Ping Test (Windows)]
\begin{enumerate}
\item Open \texttt{Control Panel \textgreater Network and Internet \textgreater Network and Sharing Center \textgreater Change adapter settings}.
\item Right-click the Ethernet adapter \textgreater \texttt{Properties} \textgreater select \texttt{Internet Protocol Version 4 (TCP/IPv4)} \textgreater \texttt{Properties}.
\item Choose \texttt{Use the following IP address} and enter:
\begin{itemize}
\item IP address (e.g. \texttt{192.168.10.11})
\item Subnet mask (\texttt{255.255.255.0})
\item Default gateway (\texttt{192.168.10.1})
\item Preferred DNS server (\texttt{8.8.8.8} or ISP DNS)
\end{itemize}
\item Click \texttt{OK} twice.
\item Open \texttt{Command Prompt} and run \texttt{ping <other\_PC\_IP>}. Verify replies with \texttt{TTL=128} (Windows) or \texttt{TTL=64} (Linux) and \texttt{time<1ms}.
\end{enumerate}
\end{methode}

\textbf{Task 6 — Presentation and Defence.}\par
Present your dossier in 10 minutes before the class (the ``school administration''). Be ready to defend: your topology choice, your budget lines, and your plan for the future third lab.

\begin{attention}[Common Mistakes to Avoid]
\begin{itemize}
\item Do not choose a topology ``because it is the best'' in general — justify it \emph{for this situation}.
\item Do not forget the inter-building link in the budget.
\item Do not buy a 24-port switch for 24 hosts: keep spare ports for growth and uplink.
\item Remember that UTP is limited to about \SI{100}{m} per run (including patch cords).
\item An IP plan where two rooms share the same host addresses, or where the gateway is missing, will fail the \texttt{ping} test.
\item Do not mix T568A on one end and T568B on the other for a straight-through cable (that creates a crossover).
\end{itemize}
\end{attention}

\courssub{Model Answer}

\textbf{Task 1 — Topology Choice.}\par
Inside each room we choose a \cle{star topology}: every PC is connected by its own UTP cable to a central switch. Compared with a \textbf{bus} (one shared coaxial cable), the star costs slightly more in cable, but a single cable failure disconnects only one PC, faults are easy to localise, and adding a PC is trivial — essential since a third lab is planned. A \textbf{ring} would complicate maintenance and is not supported by standard Ethernet equipment (no commercial Ethernet ring switches for LANs). The two building switches are linked to a central (core) switch in Building A, giving an overall \cle{tree (hierarchical star) topology}, which is the standard structure of organisational LANs.

\begin{exemplebox}[Why Not Bus or Ring?]
\textbf{Bus:} A single break in the backbone cable brings down the entire segment. Troubleshooting requires checking the whole cable. Adding a host means tapping into the backbone (vampire taps or T-connectors) — impractical for UTP. \textbf{Ring:} Token Ring / FDDI are obsolete. Modern Ethernet does not use physical rings; a break in a ring also partitions the network unless a dual-ring (redundant) is used, which doubles cable cost.
\end{exemplebox}

\textbf{Task 2 — Network Diagram.}\par

\begin{popfigure}
\begin{tikzpicture}[scale=0.95, every node/.style={font=\small}]
\node[draw, rounded corners, fill=popblueL, minimum width=1.9cm] (isp) at (0,4.2) {ISP Modem};
\node[draw, rounded corners, fill=popblueL, minimum width=1.9cm] (rtr) at (3,4.2) {Router};
\node[draw, rounded corners, fill=popgreenL, minimum width=2.2cm] (sw0) at (6.2,4.2) {Core Switch};
\node[draw, rounded corners, fill=popgreenL] (sw1) at (1.5,1.8) {Lab 1 Switch};
\node[draw, rounded corners, fill=popgreenL] (sw2) at (5,1.8) {Lab 2 Switch};
\node[draw, rounded corners, fill=popgreenL] (ap) at (8.5,1.8) {Wi-Fi AP};
\node[draw, rounded corners, fill=popgreenL] (sw3) at (11.5,1.8) {Admin Switch};
\node[draw, fill=popblueL] (pc1) at (1.5,0) {20 PCs};
\node[draw, fill=popblueL] (pc2) at (5,0) {20 PCs};
\node[draw, fill=popblueL] (lib) at (8.5,0) {3 PCs};
\node[draw, fill=popblueL] (adm) at (11.5,0) {5 PCs + Printer};
\draw[thick] (isp) -- (rtr) -- (sw0);
\draw[thick] (sw0) -- (sw1); \draw[thick] (sw0) -- (sw2);
\draw[thick] (sw0) -- (ap); \draw[thick] (sw0) -- (sw3);
\draw[thick] (sw1) -- (pc1); \draw[thick] (sw2) -- (pc2);
\draw[thick] (ap) -- (lib); \draw[thick] (sw3) -- (adm);
\node[font=\scriptsize, align=center, text width=2.2cm] at (9.2,3.1) {60 m link\\ (fibre or outdoor UTP)};
\end{tikzpicture}
\legende{Proposed tree topology for Garrisson Secondary School.}
\end{popfigure}

\textbf{Task 3 — Equipment List and Budget (Indicative Prices from Bamenda/Buea Market Survey, 2024).}\par

\begin{center}
\begin{tabularx}{\linewidth}{|X|l|l|l|}\hline
\rowcolor{popblueL} \textbf{Item} & \textbf{Qty} & \textbf{Unit Price (FCFA)} & \textbf{Total (FCFA)} \\ \hline
Switch 24 ports Gigabit (Labs 1 \& 2) & 2 & 45\,000 & 90\,000 \\ \hline
Switch 8 ports Gigabit (Admin) & 1 & 15\,000 & 15\,000 \\ \hline
Core Switch 8 ports Gigabit (Building A) & 1 & 15\,000 & 15\,000 \\ \hline
Router (Gigabit WAN/LAN, VPN capable) & 1 & 35\,000 & 35\,000 \\ \hline
ISP Modem / 4G Router (MTN/Orange/CAMTEL) & 1 & 25\,000 & 25\,000 \\ \hline
Wireless Access Point (Ceiling mount, PoE) & 1 & 30\,000 & 30\,000 \\ \hline
UTP Cable Cat6, Box of \SI{305}{m} (1000 ft) & 3 & 40\,000 & 120\,000 \\ \hline
Outdoor Fibre Link \SI{60}{m} + 2 Media Converters & 1 & 80\,000 & 80\,000 \\ \hline
RJ45 Connectors (Bag of 100) & 1 & 5\,000 & 5\,000 \\ \hline
Trunking, Patch Cords, Crimping Tool, Tester & 1 & 40\,000 & 40\,000 \\ \hline
Network Printer (Admin, Wi-Fi + Ethernet) & 1 & 150\,000 & 150\,000 \\ \hline
Installation Labour \& Contingency (10\%) & 1 & 100\,000 & 100\,000 \\ \hline
\rowcolor{popgreenL} \textbf{Total} & & & \textbf{705\,000} \\ \hline
\end{tabularx}
\end{center}

The total of about \textbf{705\,000 FCFA} stays well within the 1\,500\,000 FCFA budget, leaving a comfortable margin for the future third lab (additional 24-port switch + cable box ≈ 85\,000 FCFA). For the \SI{60}{m} inter-building link, UTP Cat6 is technically acceptable (under the \SI{100}{m} limit) if run in a protective conduit; however, \textbf{outdoor fibre with media converters} is strongly preferred for immunity to lightning, electromagnetic interference and ground loops — a critical consideration in the North West Region.

\textbf{Task 4 — Cabling.}\par
Straight-through cable, T568B order on both ends: White-Orange, Orange, White-Green, Blue, White-Blue, Green, White-Brown, Brown. After crimping, the tester lights wires 1 to 8 in sequence on both ends: the cable passes.

\begin{exoresolu}[Cable Test Interpretation]
\textbf{Statement:} A cable tester shows LEDs 1–8 lighting on the master unit, but on the remote unit LED 3 and LED 6 are swapped. What is the fault?
\tcblower
\textbf{Solution:} The wire map shows that pins 3 and 6 are crossed. In T568B, pin 3 is White-Green and pin 6 is Green. This indicates that the Green pair (pins 3 \& 6) has been reversed at one end. Re-crimp that end ensuring the correct order.
\end{exoresolu}

\textbf{Task 5 — IP Plan and Test.}\par
Private network \texttt{192.168.10.0/24} (mask \texttt{255.255.255.0}). Gateway (router LAN interface): \texttt{192.168.10.1}. Address allocation:

\begin{center}
\begin{tabular}{|l|l|l|}\hline
\rowcolor{popblueL} \textbf{Area} & \textbf{IP Range} & \textbf{Notes} \\ \hline
Lab 1 (20 PCs) & 192.168.10.11 – 192.168.10.30 & Static or DHCP pool \\ \hline
Lab 2 (20 PCs) & 192.168.10.31 – 192.168.10.50 & Static or DHCP pool \\ \hline
Library (3 PCs + Wi-Fi) & 192.168.10.51 – 192.168.10.60 & Wi-Fi clients via AP (DHCP) \\ \hline
Admin (5 PCs) & 192.168.10.61 – 192.168.10.65 & Static for servers/printers \\ \hline
Network Printer & 192.168.10.70 & Static \\ \hline
Reserved / Future & 192.168.10.71 – 192.168.10.254 & Growth \\ \hline
\end{tabular}
\end{center}

\textbf{Test:} PC1 set to \texttt{192.168.10.11}, PC2 to \texttt{192.168.10.31}. From PC1: \texttt{ping 192.168.10.31} replies successfully (4 packets sent, 4 received, 0\% loss, time <1 ms), confirming Layer 3 connectivity across the switches and core.

\textbf{Task 6 — Defence.}\par
The team presents the diagram, budget and live ping demonstration, and explains that the tree-of-stars design lets the future third lab be added by simply connecting one more 24-port switch to an available port on the core switch, with no re-cabling of existing labs. The budget margin of $\approx$ 795\,000 FCFA covers the third lab and unforeseen expenses.

\begin{exercice}[Extension: VLAN Segmentation (Optional for Advanced Students)]
The school later decides to separate student traffic (labs, library) from admin traffic for security. Propose a VLAN plan (VLAN IDs, IP subnets, trunk ports) and explain which device must perform inter-VLAN routing.
\end{exercice}

\begin{corrige}[Extension]
\begin{itemize}
\item VLAN 10 (Students): 192.168.10.0/24 — ports on Lab 1, Lab 2, Library switches + AP.
\item VLAN 20 (Admin): 192.168.20.0/24 — ports on Admin switch.
\item Trunk links: Core–Lab1, Core–Lab2, Core–Admin, Core–AP (tagged).
\item Inter-VLAN routing: Router (router-on-a-stick) or Layer 3 core switch. The router's LAN interface is configured with sub-interfaces for VLAN 10 and VLAN 20, each with its own gateway IP (e.g. 192.168.10.1 and 192.168.20.1).
\end{itemize}
\end{corrige}

\newpage

\coursec{Methods and worked examples}

\coursec{Physical Network Topologies}

\begin{definition}[Physical Network Topology]
The \cle{physical topology} of a network describes the actual geometric layout of cables, devices, and connection points. It differs from the \cle{logical topology}, which describes how signals travel and how devices communicate (e.g., a physical star wired to a hub behaves logically as a shared bus).
\end{definition}

\begin{propriete}[Six Core Physical Topologies]
The syllabus requires mastery of these six fundamental topologies:
\begin{itemize}
\item \textbf{Bus (Linear Bus):} All nodes connect via drop lines/taps to a single backbone cable (historically coaxial 10BASE2/10BASE5, later UTP in a daisy-chain). Requires terminators at both ends to absorb signals and prevent reflections.
\item \textbf{Star:} Every node connects to a central device (hub, switch, or wireless access point) via a dedicated point-to-point link. The central device manages all traffic.
\item \textbf{Ring:} Each node connects to exactly two neighbours, forming a closed loop. Data travels unidirectionally (single ring) or bidirectionally (dual ring for redundancy). Token passing or similar media access control is typical.
\item \textbf{Mesh:} Nodes interconnect with multiple redundant paths. \textbf{Full mesh:} every node connects directly to every other node ($n(n-1)/2$ links). \textbf{Partial mesh:} only critical nodes have multiple interconnections.
\item \textbf{Tree (Hierarchical/Extended Star):} Groups of star-configured networks connected to a linear bus backbone or to a higher-level central device. Forms a hierarchy: core $\rightarrow$ distribution $\rightarrow$ access.
\item \textbf{Hybrid:} A combination of two or more basic topologies (e.g., star-bus, star-ring, star-mesh). Modern enterprise networks are almost always hybrid.
\end{itemize}
\end{propriete}

\begin{methode}[Identifying and Describing Physical Network Topologies]
\begin{enumerate}
\item \textbf{List the six core topologies:} Bus, Star, Ring, Mesh, Tree (hierarchical), Hybrid.
\item \textbf{Recognise each by its physical layout:}
      \begin{itemize}
      \item \textbf{Bus:} Single backbone cable (coaxial or twisted pair) with drop lines/taps to each node.
      \item \textbf{Star:} Every node connects to a central device (hub/switch) via dedicated cable.
      \item \textbf{Ring:} Each node connects to exactly two neighbours forming a closed loop; data travels in one direction (or both in dual-ring).
      \item \textbf{Mesh:} Every node connects directly to every other node (full mesh) or to several others (partial mesh).
      \item \textbf{Tree:} Groups of star-configured networks connected to a linear bus backbone (hierarchy).
      \item \textbf{Hybrid:} Combination of two or more basic topologies (e.g., star-ring, star-bus).
      \end{itemize}
\item \textbf{Key parameters to note for each:} Number of cables/ports, cable length, single point of failure, ease of adding/removing nodes, typical media (UTP, fibre, coaxial).
\item \textbf{Reflex:} When shown a diagram, trace the physical cables — not the logical signal path — to name the topology.
\end{enumerate}
\end{methode}

\begin{exemplebox}[Topology Identification Practice]
\textbf{Scenario:} A network diagram shows 8 computers. Four computers on the left connect to Switch A. Four computers on the right connect to Switch B. Switch A and Switch B are linked by a single fibre cable. No other connections exist.
\begin{enumerate}
\item What is the physical topology?
\item How many cable links are needed for the 8 computers?
\item Identify all single points of failure.
\end{enumerate}
\textbf{Solution:}
\begin{enumerate}
\item \textbf{Hybrid (Star-Bus / Star-Fibre Backbone)} — two stars interconnected by a backbone link.
\item Each computer $\rightarrow$ its switch: 8 UTP links. Inter-switch: 1 fibre link. Total = 9 distinct cable links.
\item Each switch is a local SPOF for its segment. The inter-switch fibre link is a global SPOF: if cut, the two stars cannot communicate (though each remains internally functional).
\end{enumerate}
\end{exemplebox}

\begin{exoresolu}[Identifying Topologies from Diagrams and Descriptions]
\textbf{Statement:} 
Three physical network layouts are described below. For each, (i) name the topology, (ii) state the number of cable links required for 6 computers, (iii) identify the single point of failure (if any), and (iv) give one typical use case in a Cameroonian context.

\begin{enumerate}
\item All computers connect to a central switch using UTP patch cords. The switch is powered by mains electricity with no UPS.
\item Computers are daisy-chained along a single coaxial cable terminated at both ends with 50\,$\Omega$ terminators. A T-connector attaches each NIC to the backbone.
\item Two separate star networks (each with its own switch) are interconnected by a single fibre link between the two switches. Each star serves a different department.
\end{enumerate}
\tcblower
\textbf{Detailed Solution:}

\textbf{Layout 1:}
\begin{enumerate}
\item \textbf{Topology:} \cle{Star} (physical star, logical star).
\item \textbf{Cable links for 6 computers:} Each computer needs one patch cord to the switch $\rightarrow$ 6 links. The switch requires at least 6 ports (an 8-port switch is typical).
\item \textbf{Single point of failure:} The \textbf{central switch} (and its power supply). If the switch fails or loses power, the entire segment goes down. No UPS exacerbates this.
\item \textbf{Cameroonian use case:} A typical \textbf{cyber café in Yaoundé or Douala} with 10–20 PCs connected to a central switch; also a school computer lab (e.g., Lycée Bilingue d'Essos) where the teacher's station manages the class via the switch.
\end{enumerate}

\textbf{Layout 2:}
\begin{enumerate}
\item \textbf{Topology:} \cle{Bus} (linear bus, historically ``Ethernet 10BASE2'' or ``Thinnet'').
\item \textbf{Cable links for 6 computers:} One continuous backbone cable with 6 T-connectors $\rightarrow$ effectively \textbf{1 backbone segment} + 6 drop cables (very short). In terms of distinct cable runs: 1 main run + 6 drops.
\item \textbf{Single point of failure:} A \textbf{break anywhere on the backbone} (or a missing/loose terminator) brings down the whole network. Also, a faulty NIC that jams the bus can disrupt all.
\item \textbf{Cameroonian use case:} Legacy installation in a small \textbf{administrative office in Buea} using old coaxial cabling; or a temporary exam registration centre where laptops are quickly daisy-chained with BNC T-connectors.
\end{enumerate}

\textbf{Layout 3:}
\begin{enumerate}
\item \textbf{Topology:} \cle{Hybrid} (specifically \textbf{Star-Bus} or \textbf{Star-Fibre Backbone}). Two stars linked by a high-speed backbone.
\item \textbf{Cable links for 6 computers (3 per star):} Each star: 3 UTP links to its switch. Backbone: 1 fibre link between switches. Total distinct links = 3 + 3 + 1 = 7.
\item \textbf{Single point of failure:} Each star has its own switch as a local SPOF. The \textbf{inter-switch fibre link} is a global SPOF: if cut, the two departments cannot communicate (though each star remains internally functional).
\item \textbf{Cameroonian use case:} \textbf{University campus (e.g., University of Buea)} where Faculty of Engineering (star) and Faculty of Science (star) are linked by fibre; or a bank with two branches in Douala (Akwa and Bonanjo) each wired as a star, connected via CAMTEL fibre lease line.
\end{enumerate}

\textbf{Examiner's Tips:}
\begin{itemize}
\item Always distinguish \emph{physical} topology (cable layout) from \emph{logical} topology (signal flow). A physical star wired to a hub behaves logically as a bus (shared medium).
\item For bus topology, remember the \textbf{terminators} — they absorb signals to prevent reflections. Missing terminator $\rightarrow$ network errors.
\item Hybrid topologies are the norm in modern organisations; pure bus or ring are rare in new installations.
\end{itemize}
\end{exoresolu}

\begin{aretenir}[Key Points: Physical Topologies]
\begin{itemize}
\item Physical topology = cable layout; Logical topology = signal flow.
\item Star is the dominant LAN topology today (switched Ethernet).
\item Bus and Ring are legacy (10BASE2, Token Ring, FDDI) — know them for exams but do not recommend for new installs.
\item Mesh provides highest redundancy but at highest cost (cables, ports).
\item Tree/Hierarchical = standard for multi-floor, multi-building, or multi-department sites.
\item Hybrid = reality. Always describe the \emph{components} (e.g., ``star of stars with fibre backbone'').
\end{itemize}
\end{aretenir}

\coursec{Comparing and Choosing a Topology}

\begin{methode}[Comparing and Choosing a Network Topology]
\begin{enumerate}
\item \textbf{Define the decision criteria:} Cost (cable, ports, devices), Scalability (ease of adding nodes), Fault Tolerance (impact of a link/node failure), Performance (bandwidth sharing, collisions), Installation/Maintenance complexity.
\item \textbf{Create a comparison table} with topologies as rows and criteria as columns. Fill with qualitative ratings (High/Medium/Low) or quantitative estimates (e.g., cable length $\approx n \times d$ for star vs. $\approx (n-1)\times d$ for bus).
\item \textbf{Match scenario constraints to topology strengths:}
      \begin{itemize}
      \item \textbf{Low budget, few nodes, temporary:} Bus (if legacy) or Star with a cheap unmanaged switch.
      \item \textbf{High reliability required (bank, hospital):} Mesh (core) or Star with redundant switches/links.
      \item \textbf{Multiple departments, hierarchical management:} Tree / Extended Star.
      \item \textbf{Existing infrastructure (old building):} Hybrid — reuse whatever cabling exists.
      \end{itemize}
\item \textbf{Justify your choice in writing:} ``I recommend \emph{Extended Star} because...'' citing at least two criteria from the table.
\item \textbf{Reflex:} Never recommend a pure bus or ring for a new installation in 2025 — they are obsolete for LANs. Always assume switched Ethernet (star) as the baseline.
\end{enumerate}
\end{methode}

\begin{exoresolu}[Choosing a Topology for a New Office in Bamenda]
\textbf{Statement:} 
A new NGO office in Bamenda needs a LAN for 25 workstations, 3 network printers, and 2 servers. The building has two floors. Requirements:
\begin{itemize}
\item Budget is tight (donor-funded), but reliability is important for donor reporting.
\item The IT officer (only one) must be able to troubleshoot quickly.
\item Future expansion to 40 workstations is expected within 2 years.
\item The building has conduit trunking already installed on each floor.
\end{itemize}
\begin{enumerate}
\item Draw a comparison table (Star vs. Tree vs. Mesh) for the criteria: Cable Cost, Port Count, Fault Isolation, Ease of Expansion, Troubleshooting Simplicity.
\item Recommend the best topology and justify with three bullet points.
\item Sketch a labelled physical diagram (text description acceptable) showing switches, uplink, servers, and workstations.
\end{enumerate}
\tcblower
\textbf{Detailed Solution:}

\textbf{1. Comparison Table}

\begin{center}
\begin{tabularx}{\linewidth}{|l|X|X|X|}\hline
\rowcolor{popblueL}
\textbf{Criterion} & \textbf{Star (Single Switch)} & \textbf{Tree / Extended Star} & \textbf{Full Mesh} \\ \hline
Cable Cost & Medium (each node to central switch) & Medium-High (backbone + horizontal) & Very High ($n(n-1)/2$ links) \\ \hline
Port Count & 30+ ports on one switch (needs 48-port) & Two 24-port switches + uplink ports & 29 ports per node (impossible) \\ \hline
Fault Isolation & Excellent (one cable = one node) & Excellent (same) + backbone risk & Excellent (multiple paths) \\ \hline
Ease of Expansion & Limited by switch ports; need new switch & Add switch on new floor/segment easily & Impractical (rewire all) \\ \hline
Troubleshooting & Simple: check link lights on one switch & Simple per segment; backbone separate & Complex: many interconnections \\ \hline
\end{tabularx}
\end{center}

\textbf{2. Recommendation:} \textbf{Tree / Extended Star (Two-floor hierarchical star)}.

\textbf{Justification:}
\begin{itemize}
\item \textbf{Scalability:} Each floor gets its own 24-port access switch (one per floor). Adding a third floor or more ports only requires adding another access switch and uplinking to a core switch — no rewiring of existing drops.
\item \textbf{Fault Isolation \& Troubleshooting:} A faulty cable or NIC affects only that workstation. The IT officer can see link lights on the floor switch immediately. The uplink fibre between floors is a single monitored link.
\item \textbf{Cost-Effectiveness:} Uses two affordable 24-port unmanaged/managed switches (e.g., TP-Link TL-SG1024 or Cisco SG250) and one fibre uplink (or UTP if distance $<100$ m). Full mesh is financially impossible; single 48-port switch would require long cable runs from both floors to a central cabinet, exceeding 100 m UTP limit for ground-floor nodes.
\end{itemize}

\textbf{3. Physical Diagram Description (Text):}
\begin{center}
\begin{tabular}{|p{0.9\linewidth}|}\hline
\texttt{[Server Rack - Ground Floor] \hspace{2cm} [First Floor]}\\
\texttt{\hspace{0.5cm}| \hspace{4.5cm} |}\\
\texttt{\hspace{0.5cm}| (UTP Cat6, $<$90m) \hspace{2cm} | (UTP Cat6, $<$90m)}\\
\texttt{\hspace{0.5cm}v \hspace{4.5cm} v}\\
\texttt{+----------+ \hspace{2cm} +----------+}\\
\texttt{| 24-port  | \hspace{2cm} | 24-port  |}\\
\texttt{| Access   | \hspace{2cm} | Access   |}\\
\texttt{| Switch   | \hspace{2cm} | Switch   |}\\
\texttt{| (SW-GF)  | \hspace{2cm} | (SW-F1)  |}\\
\texttt{+----+-----+ \hspace{2cm} +----+-----+}\\
\texttt{\hspace{0.5cm}| \hspace{4.5cm} |}\\
\texttt{\hspace{0.5cm}| (Fibre Uplink / UTP if $<$100m) \hspace{1.5cm} |}\\
\texttt{\hspace{0.5cm}+---------------------------+---------------------------+}\\
\texttt{\hspace{3.5cm}|}\\
\texttt{\hspace{3.5cm}v}\\
\texttt{\hspace{2cm}+--------+--------+}\\
\texttt{\hspace{2cm}|  Core Switch    |}\\
\texttt{\hspace{2cm}|  (Optional, or   |}\\
\texttt{\hspace{2cm}|   use SW-GF as   |}\\
\texttt{\hspace{2cm}|   core)          |}\\
\texttt{\hspace{2cm}+--------+--------+}\\
\texttt{\hspace{3.5cm}|}\\
\texttt{\hspace{1.5cm}+------------+------------+}\\
\texttt{\hspace{1.5cm}|            |            |}\\
\texttt{\hspace{1.5cm}[Server 1] [Server 2] [Router/Firewall]}\\
\texttt{\hspace{1.5cm}|            |            |}\\
\texttt{\hspace{1.5cm}+------------+------------+}\\
\texttt{\hspace{3.5cm}|}\\
\texttt{\hspace{3.5cm}[Internet / WAN]}\\\hline
\end{tabular}
\end{center}

\textbf{Notes:}
\begin{itemize}
\item If budget is extremely tight, SW-GF can act as core (uplink from SW-F1 plugs into SW-GF). This saves one switch but makes SW-GF a SPOF for inter-floor traffic.
\item Use \textbf{Cat6 UTP} for all horizontal runs (max 90 m in conduit + 10 m patch cords). Use \textbf{multimode fibre (OM3/OM4)} for inter-floor uplink if distance $>100$ m or to avoid EMI in electrical trunking.
\item Label every cable at both ends (e.g., ``GF-01'', ``F1-12'') and document in a spreadsheet — essential for the lone IT officer.
\end{itemize}
\end{exoresolu}

\begin{attention}[Common Mistakes in Topology Selection]
\begin{itemize}
\item Recommending \textbf{Bus} or \textbf{Ring} for a new LAN — they are obsolete, lack fault isolation, and do not support modern switched Ethernet speeds.
\item Ignoring the \textbf{100 m UTP distance limit} — placing a single switch in a central cabinet for a multi-floor building often exceeds this limit for remote nodes.
\item Confusing \textbf{physical} and \textbf{logical} topology — a star wired to a hub is physically a star but logically a shared bus (half-duplex, collisions).
\item Forgetting \textbf{growth} — a 24-port switch filled today leaves no room for tomorrow's hires.
\item Over-engineering: proposing \textbf{full mesh} for a 20-node LAN — cost and complexity are unjustifiable.
\end{itemize}
\end{attention}

\coursec{Network Equipment}

\begin{definition}[Network Equipment by OSI Layer]
Network devices operate at specific layers of the OSI model. Understanding the layer determines the device's function, addressing scheme, and configuration complexity.
\end{definition}

\begin{propriete}[Equipment Classification by OSI Layer]
\begin{itemize}
\item \textbf{Layer 1 (Physical):} Cables (UTP Cat5e/6/6a, Fibre OM3/OM4, Coaxial), Connectors (RJ-45, LC/SC/ST, BNC), Patch panels, Wall plates, Repeaters/Hubs (obsolete), Media converters.
\item \textbf{Layer 2 (Data Link):} NIC (Network Interface Card), \textbf{Switch} (multi-port bridge), Wireless Access Point (WAP), Bridge. Forward frames based on MAC addresses.
\item \textbf{Layer 3 (Network):} \textbf{Router}, Layer-3 Switch, Firewall (often L3/L4), Modem (ISP demarcation). Forward packets based on IP addresses; perform routing, NAT, ACL, DHCP.
\item \textbf{Layer 4–7:} Load balancer, Proxy, WAF, Application Delivery Controller. Inspect payload, session, or application data.
\end{itemize}
\end{propriete}

\begin{methode}[Identifying Network Equipment and Their Roles]
\begin{enumerate}
\item \textbf{Categorise equipment by OSI layer} (see table above).
\item \textbf{Key functional distinctions:}
      \begin{itemize}
      \item \textbf{Hub vs. Switch:} Hub = shared bandwidth, half-duplex, one collision domain. Switch = dedicated bandwidth per port, full-duplex, separate collision domains, learns MAC addresses.
      \item \textbf{Switch vs. Router:} Switch forwards frames within same LAN (MAC). Router forwards packets between networks (IP), performs NAT, ACL, DHCP.
      \item \textbf{Modem vs. Router:} Modem converts ISP signal (DSL, Cable, Fibre ONT, 4G/5G) to Ethernet. Router does IP routing. Many ISP boxes (MTN/Orange Flybox, CAMTEL GPON ONT+Router) combine both.
      \item \textbf{Access Point vs. Router:} AP bridges wireless to wired (Layer 2). Wireless Router = Router + Switch + AP in one box.
      \end{itemize}
\item \textbf{Reflex:} When asked ``which device...'', first decide the layer/function: \emph{connect PCs in one room} $\rightarrow$ Switch; \emph{connect LAN to internet} $\rightarrow$ Router; \emph{extend Wi-Fi} $\rightarrow$ AP (not another router in NAT mode).
\end{enumerate}
\end{methode}

\begin{exemplebox}[Device Selection Quick Reference]
\begin{center}
\begin{tabularx}{\linewidth}{|l|l|X|}\hline
\rowcolor{popblueL}
\textbf{Need} & \textbf{Device} & \textbf{Key Feature} \\ \hline
Connect 20 PCs in a lab & 24-port Gigabit Switch (L2) & MAC learning, full-duplex, VLAN support \\ \hline
Connect LAN to CAMTEL fibre & Router / Gateway & PPPoE client, NAT, Firewall, DHCP server \\ \hline
Provide Wi-Fi in a hall & Ceiling AP (PoE) + PoE Switch/Injector & 802.11ax (Wi-Fi 6), VLAN tagging, roaming \\ \hline
Connect two buildings 500 m apart & Media Converters + Fibre (or Fibre SFP in switches) & OM3/OM4 multimode or OS2 single-mode \\ \hline
Load balance two ISP links & Dual-WAN Router / SD-WAN appliance & Failover, load balancing, policy routing \\ \hline
\end{tabularx}
\end{center}
\end{exemplebox}

\begin{exoresolu}[Selecting the Right Device for Specific Functions]
\textbf{Statement:} 
A small hotel in Limbe wants to provide guest Wi-Fi and wired internet for the reception desk, manager's office, and conference room. They have a CAMTEL fibre subscription (GPON ONT provided). The manager buys the following equipment from a local vendor:
\begin{itemize}
\item Device A: 8-port Gigabit Unmanaged Switch (TP-Link TL-SG108)
\item Device B: Wireless Router (TP-Link Archer C6) — labelled ``Router/AP/Repeater''
\item Device C: PoE Injector (for a single ceiling AP)
\item Device D: USB 4G Dongle (Huawei E3372) — ``for backup''
\end{itemize}
\begin{enumerate}
\item For each device, state its primary OSI layer(s) and its role in this network.
\item The manager plugs the ONT into Device B's WAN port, connects Device A to a LAN port of Device B, and plugs the reception PC, office PC, and conference room PC into Device A. He sets Device B to ``Wireless Router'' mode. Guests complain they can access the reception PC's shared folder. Explain why and how to fix it using the existing hardware.
\item The manager wants to add a ceiling-mounted AP in the garden for better coverage. He has a PoE-capable AP (Ubiquiti U6-Lite) and Device C (PoE Injector). Describe the correct cabling from Device A to the AP.
\item If the CAMTEL fibre fails, how can Device D be used for automatic failover? What additional feature must Device B support?
\end{enumerate}
\tcblower
\textbf{Detailed Solution:}

\textbf{1. Device Roles and OSI Layers}

\begin{center}
\begin{tabularx}{\linewidth}{|l|l|X|}\hline
\rowcolor{popblueL}
\textbf{Device} & \textbf{Primary OSI Layer(s)} & \textbf{Role in This Network} \\ \hline
A: 8-port Switch & Layer 2 (Data Link) & Expands wired ports: connects reception, office, conference PCs to the router's LAN. Creates separate collision domains, learns MACs. \\ \hline
B: Wireless Router & Layer 3 (Network) + Layer 2 (Switch/AP) & \textbf{Edge router}: terminates CAMTEL ONT (WAN), does NAT, DHCP, firewall for entire hotel LAN. \textbf{Built-in AP}: provides guest Wi-Fi. \textbf{Built-in 4-port switch}: connects Device A and other wired devices. \\ \hline
C: PoE Injector & Layer 1 (Physical) & Injects DC power (48V) onto Ethernet cable to power the ceiling AP (Device C has no data processing role). \\ \hline
D: 4G USB Dongle & Layer 1 + Layer 2 (Modem) & Backup WAN link: provides cellular internet (MTN/Orange 4G) when fibre fails. Appears as Ethernet or RNDIS interface to router. \\ \hline
\end{tabularx}
\end{center}

\textbf{2. Guest Access to Reception PC — Cause and Fix}
\begin{itemize}
\item \textbf{Cause:} Device B is in \textbf{``Wireless Router'' mode} (default). Its built-in switch ports (LAN1–LAN4) and Wi-Fi interfaces are \textbf{bridged together in the same LAN subnet} (e.g., 192.168.0.0/24). Guests on Wi-Fi and wired PCs on Device A (connected to LAN port of B) are in the \textbf{same broadcast domain}. No firewall rule isolates them $\rightarrow$ guests can SMB-browse to reception PC.
\item \textbf{Fix using existing hardware:} Enable \textbf{Guest Network} (or ``AP Isolation'' / ``Client Isolation'') on Device B's Wi-Fi settings. This creates a separate virtual SSID (e.g., ``Hotel\_Guest'') with its own subnet (e.g., 192.168.1.0/24) and firewall rules blocking access to the main LAN. \textbf{Do not} plug the AP into a separate router (double NAT) — the Archer C6 supports Guest Network natively.
\item \textbf{Alternative (if no Guest Network):} Put Device B in \textbf{Access Point mode}, buy a separate router for the main LAN, and use VLANs — but that requires new hardware. Guest Network is the correct answer here.
\end{itemize}

\textbf{3. Cabling for Ceiling AP with PoE Injector}
\begin{enumerate}
\item Run a single \textbf{Cat6 UTP cable} (max 90 m solid + 5 m patch each end) from a free port on \textbf{Device A (Switch)} to the \textbf{PoE Injector's ``Data In'' / ``LAN'' port} (located indoors, near switch).
\item Connect a short patch cord from Injector's \textbf{``PoE Out'' / ``Data \& Power'' port} to the \textbf{AP's Ethernet port} (the cable continues through ceiling conduit to AP).
\item \textbf{Critical:} The injector must be \textbf{802.3af/at compliant} (the U6-Lite needs 802.3af). The TP-Link injector (TL-POE150S or similar) works. Do \textbf{not} use a passive 24V injector — it will damage the Ubiquiti AP.
\item Power the injector from a mains socket (preferably on UPS).
\end{enumerate}

\textbf{4. 4G Failover with Device D}
\begin{itemize}
\item \textbf{Requirement:} Device B (Archer C6) must support \textbf{WAN Failover / Dual WAN / Backup WAN} with a \textbf{USB WAN} or \textbf{Ethernet WAN} port configurable for failover.
\item \textbf{Setup:} Plug Device D (4G dongle) into Device B's USB port. In Device B's web UI: \textit{Network $\rightarrow$ Internet $\rightarrow$ Backup WAN} $\rightarrow$ select ``USB Modem'', configure APN (e.g., ``mtncameroon.net'' or ``orange.cm''), enable ``Failover Mode'' (primary = WAN port / fibre, backup = USB). Set detection method (ping to 8.8.8.8 every 30 s).
\item \textbf{Note:} The Archer C6 v3+ supports USB modem failover. If the model lacks this, a separate router with dual WAN (e.g., MikroTik hAP ac², Ubiquiti EdgeRouter X) is needed — Device D alone cannot trigger failover.
\end{itemize}
\end{exoresolu}

\begin{savaistu}[Cameroon ISP Equipment Context]
\begin{itemize}
\item \textbf{CAMTEL GPON:} Provides an ONT (Optical Network Terminal) — typically Huawei or ZTE — with one or more Gigabit Ethernet ports. The ONT is Layer 1/2; you connect your router's WAN port to it. PPPoE credentials are required.
\item \textbf{MTN/Orange Fibre:} Often provide a combined ONT+Router (``box''). You can request ``bridge mode'' to use your own router, or use DMZ/port forwarding.
\item \textbf{4G/5G Backup:} MTN and Orange sell USB dongles (Huawei, ZTE) and MiFi routers. For business, ask for a ``Fixed Wireless Access'' (FWA) router with external antenna ports — more stable than a dongle.
\item \textbf{ANTIC Compliance:} All network equipment handling personal data must comply with Law No. 2010/012 on Cybersecurity and Cybercrime. Logs must be kept for 12 months; incident reporting to ANTIC is mandatory.
\end{itemize}
\end{savaistu}

\coursec{Setting Up a Computer Network}

\begin{methode}[Setting Up a Small Computer Network (Step-by-Step)]
\begin{enumerate}
\item \textbf{Phase 1 — Planning:} Draw topology, list devices, assign IP scheme (private range, subnet mask, gateway, DNS), label cables.
\item \textbf{Phase 2 — Physical Installation:}
      \begin{enumerate}
      \item Install switches/router in ventilated, secure cabinet/rack.
      \item Run horizontal cables (Cat6) from wall plates to patch panel (or directly to switch). \textbf{Respect 100 m max (90 m + 10 m patches).}
      \item Terminate cables: TIA/EIA-568B colour code on both ends (patch panel + RJ-45 plugs). Test every link with cable tester (wiremap, length, attenuation).
      \item Connect patch cords: Patch panel $\rightarrow$ Switch; Wall plate $\rightarrow$ PC/Printer.
      \item Connect WAN: ONT/Modem $\rightarrow$ Router WAN.
      \item Power on: Modem $\rightarrow$ Router $\rightarrow$ Switches $\rightarrow$ APs.
      \end{enumerate}
\item \textbf{Phase 3 — Logical Configuration:}
      \begin{enumerate}
      \item Access router web UI (default IP, admin/admin). Change admin password immediately.
      \item Configure WAN: PPPoE (CAMTEL), DHCP (MTN/Orange), or Static IP.
      \item Configure LAN: Private subnet (e.g., 192.168.10.0/24), enable DHCP server (pool .100–.200), set DNS (1.1.1.1, 8.8.8.8).
      \item Configure Wi-Fi: SSID, WPA3-Personal (or WPA2-AES), strong passphrase, hide SSID optional.
      \item Configure Guest Network (isolated subnet, bandwidth limit).
      \item Assign static IPs / DHCP reservations for servers, printers, APs.
      \end{enumerate}
\item \textbf{Phase 4 — Testing \& Documentation:}
      \begin{enumerate}
      \item Ping gateway from each PC. Ping 8.8.8.8. Browse web.
      \item Test printer, file share, speed test (iperf3 or fast.com).
      \item Document: IP plan, device passwords (in sealed envelope), cable map, Wi-Fi channels.
      \end{enumerate}
\item \textbf{Reflex:} Always test \textbf{before} closing trunking/ceilings. Label \textbf{both ends} of every cable. Use a \textbf{UPS} for router + core switch.
\end{enumerate}
\end{methode}

\begin{exoresolu}[Complete Setup of a 10-PC Office Network in Douala]
\textbf{Statement:} 
You are contracted to set up a LAN for a law firm in Douala (Akwa). Requirements:
\begin{itemize}
\item 10 workstations (Windows 11), 1 network printer (HP LaserJet), 1 file server (Ubuntu Server).
\item Internet via MTN Fibre (GPON ONT provided, PPPoE credentials given).
\item Wired connections for all fixed devices; Wi-Fi for lawyers' laptops and guests.
\item Budget allows: 1× 16-port Gigabit Switch (managed), 1× Wireless Router (dual-band, VLAN support), 1× Ceiling AP (PoE), Cat6 cable, RJ-45, patch panel, 6U wall cabinet, UPS (1000 VA).
\end{itemize}
\begin{enumerate}
\item Propose an IP addressing plan (subnet, gateway, DHCP pool, static reservations, VLANs for Guest/Staff).
\item List the exact cabling steps from ONT to switch to devices, including patch panel termination standard.
\item Write the router configuration checklist (WAN, LAN, VLANs, DHCP, Firewall, Wi-Fi, Admin security).
\item Describe how you would test and document the installation for handover.
\end{enumerate}
\tcblower
\textbf{Detailed Solution:}

\textbf{1. IP Addressing Plan (VLAN-Segmented)}

\begin{center}
\begin{tabularx}{\linewidth}{|l|l|l|X|}\hline
\rowcolor{popblueL}
\textbf{VLAN / Network} & \textbf{Subnet / CIDR} & \textbf{Gateway} & \textbf{Usage / DHCP Pool} \\ \hline
VLAN 10 — Staff Wired & 192.168.10.0/24 & 192.168.10.1 & Workstations, Printer, Server. DHCP: .100–.150. Reservations: Server .10, Printer .11. \\ \hline
VLAN 20 — Staff Wi-Fi & 192.168.20.0/24 & 192.168.20.1 & Lawyers' laptops, phones. DHCP: .100–.150. Same firewall rules as VLAN 10 (inter-VLAN routing allowed). \\ \hline
VLAN 30 — Guest Wi-Fi & 192.168.30.0/24 & 192.168.30.1 & Guests only. DHCP: .100–.200. \textbf{No access to VLAN 10/20}. Internet only. Bandwidth limit: 5 Mbps/user. \\ \hline
Management (untagged on switch) & 192.168.99.0/24 & 192.168.99.1 & Switch, Router, AP management IPs. No DHCP. Static only. \\ \hline
\end{tabularx}
\end{center}

\textbf{2. Cabling Steps (Physical Installation)}

\begin{enumerate}
\item \textbf{Cabinet mounting:} Fix 6U cabinet on wall in secure room (server room / storage). Install patch panel (16-port), switch, router, UPS inside.
\item \textbf{Horizontal runs:} Pull Cat6 UTP from cabinet to each wall outlet (10 workstations + 1 printer + 1 server + 2 AP locations = 14 drops). Use conduit/trunking. Leave 1 m slack at each end.
\item \textbf{Termination at patch panel (cabinet end):} Strip jacket, untwist pairs \textbf{$<$13 mm}, terminate on 110-type IDC blocks following \textbf{TIA/EIA-568B} colour code:
      \begin{center}
      \begin{tabular}{|c|c|}\hline
      \textbf{Pin} & \textbf{Colour (568B)} \\ \hline
      1 & White/Orange \\ \hline
      2 & Orange \\ \hline
      3 & White/Green \\ \hline
      4 & Blue \\ \hline
      5 & White/Blue \\ \hline
      6 & Green \\ \hline
      7 & White/Brown \\ \hline
      8 & Brown \\ \hline
      \end{tabular}
      \end{center}
      Punch down with Krone tool. Label each port (WS-01 … WS-10, PRN-01, SRV-01, AP-01, AP-02).
\item \textbf{Termination at wall plates (room end):} Same 568B on keystone jacks. Snap into faceplates.
\item \textbf{Patch cords (cabinet):} Factory-made Cat6 patch cords (1 m) from patch panel ports $\rightarrow$ Switch ports. Match labels.
\item \textbf{Device patch cords (rooms):} 2–3 m Cat6 patch cords from wall plate $\rightarrow$ PC/Printer/Server.
\item \textbf{AP cabling:} Run one drop to ceiling AP location. At cabinet, patch panel port $\rightarrow$ PoE injector ``Data In'' $\rightarrow$ PoE injector ``Data+Power Out'' $\rightarrow$ patch cord to AP (or direct if switch is PoE+). Second AP (router's built-in) needs no extra cable.
\item \textbf{WAN:} ONT (GPON) $\rightarrow$ Router WAN port (Cat6 patch cord).
\item \textbf{Power:} Router, Switch, PoE Injector $\rightarrow$ UPS battery-backed outlets. ONT $\rightarrow$ UPS if possible.
\end{enumerate}

\textbf{3. Router Configuration Checklist (Managed Router, e.g., MikroTik hAP ax³ or Ubiquiti Dream Machine / ER-X + Switch + AP)}

\begin{enumerate}
\item \textbf{Admin Security:} Change default user `admin` password to 16+ char. Create limited `tech` user. Disable WinBox/SSH from WAN. Enable 2FA if supported.
\item \textbf{WAN (PPPoE):} Interface = ether1 (WAN). PPPoE client: username/password from MTN. Add default route. MTU 1492 (PPPoE overhead).
\item \textbf{VLANs on Switch Ports (Bridge):} Create bridge `bridge-LAN`. Add VLAN filtering. Define VLANs 10, 20, 30, 99. Assign access ports: ether2–ether11 = VLAN 10 (untagged). ether12 (to AP) = trunk: tagged VLAN 20, 30, 99. ether13 (to ceiling AP) = trunk: tagged VLAN 20, 30, 99.
\item \textbf{IP Addresses on VLAN Interfaces:} Add IP 192.168.10.1/24 on VLAN10 interface, 192.168.20.1/24 on VLAN20, 192.168.30.1/24 on VLAN30, 192.168.99.1/24 on VLAN99.
\item \textbf{DHCP Servers:} One per VLAN (10, 20, 30) with pools/gateways/DNS as in table. VLAN 99: no DHCP.
\item \textbf{Firewall Filter Rules (IPv4):}
      \begin{itemize}
      \item Input chain: Allow established/related, ICMP, DHCP, DNS from LAN. Drop all from WAN.
      \item Forward chain: Allow VLAN10 $\leftrightarrow$ VLAN20 (staff). Allow VLAN10/20 $\rightarrow$ WAN. Allow VLAN30 $\rightarrow$ WAN. \textbf{Drop VLAN30 $\rightarrow$ VLAN10/20/99}. Drop invalid.
      \end{itemize}
\item \textbf{NAT (Masquerade):} Outbound on WAN interface for all LAN subnets (srcnat chain).
\item \textbf{Wi-Fi (Router built-in + Ceiling AP):}
      \begin{itemize}
      \item SSID ``LawFirm\_Staff'' — WPA3-Personal (or WPA2-AES), VLAN 20, 5 GHz preferred, 2.4 GHz enabled.
      \item SSID ``LawFirm\_Guest'' — WPA3-Personal (or open + captive portal), VLAN 30, bandwidth limit.
      \item Channel: Manual non-overlapping (1, 6, 11 on 2.4; 36, 40, 44, 48 on 5). TX power: medium.
      \end{itemize}
\item \textbf{Time/NTP:} Set timezone Africa/Douala, NTP servers (pool.ntp.org).
\item \textbf{Backup:} Export config to file, save to USB/email.
\end{enumerate}

\textbf{4. Testing \& Handover Documentation}

\begin{enumerate}
\item \textbf{Cable Certification:} Fluke DSX-5000 or similar — print PDF report for each run (Pass/Fail Cat6).
\item \textbf{Connectivity Tests:}
      \begin{itemize}
      \item Ping gateway from each VLAN (PC, Laptop, Phone).
      \item Ping 8.8.8.8 and resolve google.com (DNS).
      \item Print test page from each workstation to network printer.
      \item Access file server share (SMB) from staff VLANs only.
      \item Verify Guest Wi-Fi cannot ping 192.168.10.10 (server).
      \item Speed test: iperf3 between server and workstation ($>900$ Mbps), internet speed test ($>90\%$ of subscribed plan).
      \item Failover test: Disconnect ONT $\rightarrow$ verify 4G backup (if configured) or note manual procedure.
      \end{itemize}
\item \textbf{Handover Pack (Printed + Digital):}
      \begin{itemize}
      \item Network diagram (Visio/draw.io) with IP plan, VLANs, port mapping.
      \item Cable test reports.
      \item Device list: Model, Serial, Firmware version, Management IP, Admin credentials (sealed envelope).
      \item Wi-Fi heatmap (Ekahau/NetSpot screenshot).
      \item Maintenance schedule: Firmware check quarterly, UPS battery test 6-monthly, password rotation yearly.
      \item Contact: Your phone/email, MTN Business support number, CAMTEL fault line.
      \end{itemize}
\end{enumerate}
\end{exoresolu}

\begin{attention}[Installation Pitfalls]
\begin{itemize}
\item \textbf{Wrong colour code:} Mixing 568A and 568B creates crossover cables — use 568B everywhere in Cameroon.
\item \textbf{Untested cables:} Always test \emph{before} closing trunking. A single bad punch-down costs hours later.
\item \textbf{No labels:} Unlabelled cables make troubleshooting impossible for the next technician.
\item \textbf{Default passwords:} Leaving `admin/admin` on the router is a security violation (ANTIC Law).
\item \textbf{No UPS:} Power fluctuations in Cameroon (ENEO) will corrupt switch/router configs and shorten hardware life.
\item \textbf{Ignoring 100 m limit:} Measure runs! Use fibre for $>90$ m horizontal or inter-building links.
\end{itemize}
\end{attention}

\coursec{Selecting Equipment for Organisational Networks}

\begin{methode}[Selecting Equipment for Organisational Networks]
\begin{enumerate}
\item \textbf{Gather requirements:} Number of users (wired/wireless), bandwidth per user, applications (VoIP, video, ERP), growth forecast (2–3 years), budget, physical constraints (building layout, cabling distance), redundancy needs (uptime SLA), management capability (CLI vs. cloud vs. GUI).
\item \textbf{Define network blocks:} Core, Distribution, Access (Three-tier) or Collapsed Core (Two-tier) or Leaf-Spine (Data Centre). For typical Cameroonian SME: \textbf{Collapsed Core (Core = Distribution)} + Access switches.
\item \textbf{Select per block:}
      \begin{itemize}
      \item \textbf{Core/Distribution:} Layer 3 switch (routing, VLANs, QoS, stacking). Ports: 24/48 × 10G/1G SFP+ uplinks. Redundant PSU. Licensing (base vs. advanced).
      \item \textbf{Access:} PoE+ (802.3at) for APs/Phones/Cameras. 24/48 ports. Stacking or uplink to core. Fanless for quiet offices.
      \item \textbf{Wireless:} Wi-Fi 6/6E (802.11ax) APs. Controller: Cloud (Ubiquiti, Meraki, Omada) or On-prem (WLC). Density: 1 AP per 2–3 rooms / 30–50 users. Site survey mandatory.
      \item \textbf{Edge Router/Firewall:} Throughput $\geq$ ISP speed (e.g., 1 Gbps fibre $\rightarrow$ 2 Gbps firewall throughput). Features: VPN (IPsec/WireGuard), IDS/IPS, SSL inspection, SD-WAN if multi-site.
      \item \textbf{Cabling:} Cat6A for new builds (10G ready, 100 m). OM4 fibre for backbone ($>100$ m or 10G/25G/40G). Patch panels, cable management, labelling.
      \item \textbf{Power:} UPS (Online Double Conversion) for core + edge. Runtime 30–60 min. PDU with monitoring.
      \end{itemize}
\item \textbf{Vendor strategy:} Single vendor (simpler support, licensing) vs. Best-of-breed. In Cameroon: \textbf{Ubiquiti (UniFi)} — popular, cloud-managed, good value. \textbf{MikroTik} — powerful, cheap, CLI-heavy. \textbf{Cisco/Aruba/HPE} — enterprise, expensive, local partner support (e.g., \textbf{Cameroon Telecom / MTN Business / Orange Business} resell). \textbf{TP-Link Omada} — budget-friendly, cloud controller.
\item \textbf{Reflex:} Always size \textbf{uplinks} at 10× access port speed (1G access $\rightarrow$ 10G uplink). Plan for \textbf{stacking} or \textbf{MLAG} for core redundancy. Don't forget \textbf{spare ports} (20\% headroom).
\end{enumerate}
\end{methode}

\begin{exoresolu}[Equipment Selection for a Growing Cameroonian SME]
\textbf{Statement:} 
``KolaTech'', a software startup in Yaoundé, currently has 30 employees in one open-plan office (single floor). They plan to hire 20 more in 18 months and open a second floor in the same building. Current pain points: Wi-Fi dead zones, slow file transfers to on-prem GitLab server, single ISP (Orange Fibre 100 Mbps) with no backup. They need a new network design and equipment list. Budget: 3,500,000 FCFA (≈ 5,300 EUR) for hardware only.
\begin{enumerate}
\item Propose a two-tier topology (Collapsed Core + Access) with VLANs. Draw a logical diagram (text description).
\item Select specific models (with approximate FCFA prices) for: Core Switch, Access Switch(es), Wireless APs (Wi-Fi 6), Router/Firewall, UPS, Cabling (Cat6A, Fibre). Justify each choice against requirements.
\item Show the budget breakdown and confirm it fits within 3.5M FCFA.
\item Explain how the design accommodates the second floor expansion without rewiring the core.
\end{enumerate}
\tcblower
\textbf{Detailed Solution:}

\textbf{1. Two-Tier Topology (Collapsed Core) — Logical Diagram}
\begin{center}
\begin{tabular}{|p{0.9\linewidth}|}\hline
\texttt{[Orange Fibre ONT] --- (PPPoE) ---> [Edge Router/Firewall]}\\
\texttt{\hspace{3.5cm}|}\\
\texttt{\hspace{3.5cm}| (10G SFP+)}\\
\texttt{\hspace{3.5cm}v}\\
\texttt{+------------------+}\\
\texttt{|  CORE SWITCH     |}\\
\texttt{|  (L3, Routing,   |}\\
\texttt{|   VLANs, QoS)    |}\\
\texttt{|  24x 1G/10G SFP+ |}\\
\texttt{+--------+---------+}\\
\texttt{\hspace{1.5cm}|}\\
\texttt{\hspace{0.5cm}+---------------------+---------------------+---------------------+}\\
\texttt{\hspace{0.5cm}| (10G Fibre)         | (10G Fibre)         | (1G Uplink)         |}\\
\texttt{\hspace{0.5cm}v                     v                     v}\\
\texttt{+--------------+      +--------------+      +--------------+}\\
\texttt{| ACCESS SW 1  |      | ACCESS SW 2  |      | ACCESS SW 3  |}\\
\texttt{| (Ground Floor)|      | (1st Floor)  |      | (Server/     |}\\
\texttt{| 48-port PoE+ |      | 24-port PoE+ |      |  Mgmt)       |}\\
\texttt{| Stack Member |      | Stack Member |      | 12-port      |}\\
\texttt{+------+-------+      +------+-------+      +------+-------+}\\
\texttt{\hspace{0.5cm}|                     |                     |}\\
\texttt{\hspace{0.5cm}|                     |                     |}\\
\texttt{[AP-1]           [AP-2]    [PC/Phone]  [AP-3]               [GitLab Server]}\\
\texttt{(Wi-Fi 6)        (Wi-Fi 6)  (VLAN10)   (Wi-Fi 6)             [NAS Backup]}\\
\texttt{\hspace{0.5cm}|                     |                     |                     |}\\
\texttt{\hspace{0.5cm}+---------------------+---------------------+---------------------+}\\
\texttt{\hspace{3.5cm}|}\\
\texttt{\hspace{3.5cm}VLANs: 10(Staff), 20(VoIP), 30(Guest), 40(IoT), 99(Mgmt)}\\
\texttt{\hspace{3.5cm}|}\\
\texttt{\hspace{3.5cm}Inter-VLAN Routing on Core Switch}\\
\texttt{\hspace{3.5cm}|}\\
\texttt{\hspace{3.5cm}Firewall: NAT, VPN, IDS/IPS, Failover to 4G}\\\hline
\end{tabular}
\end{center}

\textbf{2. Equipment Selection \& Justification (Prices approximate, Yaoundé market 2024–2025)}

\begin{center}
\begin{tabularx}{\linewidth}{|l|X|l|l|X|}\hline
\rowcolor{popblueL}
\textbf{Item} & \textbf{Model (Example)} & \textbf{Qty} & \textbf{Unit Price (FCFA)} & \textbf{Justification} \\ \hline
Core Switch & \textbf{MikroTik CRS326-24G-2S+RM} (24×1G RJ45 + 2×10G SFP+) & 1 & 420,000 & L3 switching, VLAN routing, 10G uplinks to access switches, rackmount, dual PSU option, RouterOS (powerful, no licence fees). Cheaper than Cisco/Ubiquiti Enterprise. \\ \hline
Access Switch 1 (GF) & \textbf{Ubiquiti USW-Pro-48-PoE} (48×1G PoE+ 600W + 4×10G SFP+) & 1 & 950,000 & 48 PoE+ ports for 30 current + 20 future PCs, 4 APs, IP phones. 10G uplinks to core. UniFi cloud management (easy for small IT team). Stackable via 10G. \\ \hline
Access Switch 2 (1st Floor) & \textbf{Ubiquiti USW-Pro-24-PoE} (24×1G PoE+ 400W + 2×10G SFP+) & 1 & 580,000 & Same OS, stacking, PoE for 2nd floor APs/phones. 10G uplink to core. \\ \hline
Access Switch 3 (Server) & \textbf{MikroTik CRS309-1G-8S+IN} (1G mgmt + 8×10G SFP+) & 1 & 350,000 & Pure SFP+ switch for server room: connect GitLab, NAS, Backup at 10G. No PoE needed. Managed by same RouterOS. \\ \hline
Wireless APs (Wi-Fi 6) & \textbf{Ubiquiti U6-Pro} (Wi-Fi 6, 2×2:2 5GHz + 2×2:2 2.4GHz, PoE+) & 4 & 180,000 & 4 APs: 2 ground, 2 first floor. Ceiling mount. Centralised UniFi controller (on Cloud Key or self-hosted). Supports VLANs, band steering, airtime fairness. \\ \hline
Router/Firewall & \textbf{MikroTik CCR2004-1G-12S+2XS} (1×1G RJ45, 12×10G SFP+, 2×25G SFP28) & 1 & 650,000 & Wire-speed 10G routing, IPsec, WireGuard, BGP, OSPF, powerful firewall, no subscription. Fits budget better than UXG-Pro. \\ \hline
UPS (Core+Edge) & \textbf{Eaton 5P 1550i Rackmount} (1550 VA / 1100 W, 2U) + Network Card & 1 & 550,000 & Online double-conversion, 30 min runtime for core switch, router, 2 access switches, ONT. Network card for graceful shutdown. \\ \hline
Cabling: Cat6A UTP (305m box) & \textbf{Belden / Legrand / Local certified} & 4 boxes & 85,000 & 4 boxes = 1220 m. Enough for 50 drops × 20 m avg = 1000 m + slack. Cat6A supports 10G to 100 m. \\ \hline
Cabling: OM4 Fibre (100m drum) & \textbf{Corning / FS.com} 12-core OM4 & 1 & 120,000 & Backbone: Core $\leftrightarrow$ Access1 (GF), Core $\leftrightarrow$ Access2 (1st Floor), Core $\leftrightarrow$ Server SW. 6 fibres used (3 pairs), 6 spare. \\ \hline
Fibre Patch Cords (LC-LC, 3m) & OM4 Duplex & 6 & 8,000 & Core-to-Access uplinks (3 links × 2 ends). \\ \hline
SFP+ Modules (10G SR) & \textbf{FS.com / MikroTik} 10GBASE-SR & 12 & 12,000 & 3 uplinks × 2 ends × 2 modules (Tx/Rx) = 12. Cheap compatible optics. \\ \hline
Patch Panel (24-port Cat6A) & \textbf{Legrand / Panduit} & 3 & 45,000 & One per floor + server room. \\ \hline
Wall Plates + Keystones (Cat6A) & & 50 & 2,500 & 50 drops. \\ \hline
Cabinet (12U Wall Mount) & \textbf{Local fabrication / APC} & 1 & 180,000 & For core/router/UPS in server room. \\ \hline
Cable Management, Labels, Tools & Consumables & 1 lot & 150,000 & Velcro, cable ties, label printer tape, RJ45 crimper, fibre cleaner. \\ \hline
\multicolumn{4}{|r|}{\textbf{TOTAL ESTIMATE}} & \textbf{3,386,000 FCFA} \\ \hline
\end{tabularx}
\end{center}

\textbf{Budget Status:} 3,386,000 FCFA $<$ 3,500,000 FCFA. \textbf{Remaining ~114,000 FCFA} for contingencies (extra SFP+, longer fibre, tax).

\textbf{3. Expansion to Second Floor Without Rewiring Core}
\begin{itemize}
\item \textbf{Core switch} already has free 10G SFP+ ports (24 total, using 3 for uplinks). Adding 2nd floor Access Switch (SW2) only requires:
      \begin{enumerate}
      \item Pulling \textbf{one 12-core OM4 fibre} from core cabinet to 2nd floor cabinet (through existing building riser/conduit).
      \item Terminating 2 fibres (1 pair) on LC patch panels at each end.
      \item Inserting 2× 10G SR SFP+ modules (one in core, one in SW2).
      \item Configuring the new port on core as trunk carrying VLANs 10,20,30,40,99.
      \item Adopting SW2 in UniFi controller (or RouterOS) — zero-touch provisioning if pre-configured.
      \end{enumerate}
\item \textbf{No changes} to core switch routing, VLANs, firewall rules, DHCP pools. The new switch simply extends the existing Layer 2 domains.
\item \textbf{Wireless:} Add 2 more U6-Pro APs on 2nd floor, plug into SW2 PoE ports. They adopt to same UniFi controller, inherit WLAN configs (VLAN mapping). Done.
\end{itemize}

\textbf{Key Advantage:} The \textbf{Collapsed Core with 10G uplinks} makes the core a non-blocking fabric. Adding floors is just adding leaf switches — classic ``spine-leaf'' simplicity without full spine-leaf cost.
\end{exoresolu}

\begin{aretenir}[Equipment Selection Checklist]
\begin{itemize}
\item Match port count + PoE budget to device count (PCs, APs, phones, cameras) + 20\% headroom.
\item Uplinks = 10× access speed (1G $\rightarrow$ 10G, 10G $\rightarrow$ 100G).
\item Choose management style: Cloud (UniFi, Omada, Meraki) for ease; CLI (MikroTik, Cisco) for control.
\item Verify local vendor support \& warranty in Cameroon (MTN Business, Orange Business, CAMTEL, local integrators).
\item Budget for: UPS (online double-conversion), fibre (OM4/OS2), SFP+ modules, patch panels, cabinet, labelling.
\item Plan for redundancy: dual PSU, stacking/MLAG, dual ISP (SD-WAN/failover).
\end{itemize}
\end{aretenir}

\coursec{Integration Activities}

\begin{methode}[Integrated Network Design and Troubleshooting (Synthesis)]
\begin{enumerate}
\item \textbf{Read the scenario twice:} Identify \textbf{constraints} (budget, space, existing gear, ISP, regulations), \textbf{requirements} (users, apps, uptime, security, growth), and \textbf{deliverables} (diagram, config, test plan, bill of materials).
\item \textbf{Apply the ``Divide and Conquer'' design loop:}
      \begin{enumerate}
      \item \textbf{Physical Layer:} Cabling routes, distances, media (Cat6A/OM4), cabinet locations, power/UPS, grounding, labelling scheme.
      \item \textbf{Data Link Layer:} VLAN plan, STP/RSTP/MSTP root placement, port profiles (access/trunk, PoE, storm control, BPDU guard), switch stacking/MLAG.
      \item \textbf{Network Layer:} IP addressing (VLSM, summarisation), routing (static, OSPF, BGP), gateway redundancy (VRRP/HSRP), NAT, VPN.
      \item \textbf{Transport/Application:} QoS (DSCP marking for VoIP/Video), ACLs, firewall zones, DNS/DHCP/NTP servers, monitoring (SNMP, Syslog, NetFlow).
      \item \textbf{Wireless:} Site survey (predictive $\rightarrow$ validation), channel plan, TX power, SSID/VLAN mapping, roaming (802.11r/k/v), guest isolation.
      \item \textbf{Security:} 802.1X (RADIUS), MAC auth bypass, port security, DHCP snooping, DAI, ND inspection, segmentation (DMZ, Guest, Mgmt).
      \item \textbf{Resilience:} Dual ISP (SD-WAN/failover), dual core, redundant PSUs, link aggregation (LACP), backups (config, firmware).
      \end{enumerate}
\item \textbf{Troubleshooting Methodology (GCE style):}
      \begin{enumerate}
      \item \textbf{Identify symptoms:} Who, what, when, where, how many affected.
      \item \textbf{Establish scope:} Single host $\rightarrow$ cable/NIC/IP. Single VLAN $\rightarrow$ access switch. Multiple VLANs $\rightarrow$ core/router/firewall. Wireless only $\rightarrow$ AP/Controller/Channel.
      \item \textbf{Layer-by-layer test:} Physical (link light, cable tester) $\rightarrow$ Data Link (MAC table, VLAN, STP) $\rightarrow$ Network (ping, traceroute, route table) $\rightarrow$ Transport (port open, firewall) $\rightarrow$ Application (DNS, HTTP, credentials).
      \item \textbf{Use tools:} ping, traceroute, nslookup, ipconfig/ifconfig, arp -a, netstat, Wireshark (capture), cable tester, Wi-Fi analyser.
      \item \textbf{Document fix \& update docs.}
      \end{enumerate}
\item \textbf{Reflex:} In exams, always \textbf{draw a diagram} first. Label IPs, VLANs, interfaces. Write config snippets (CLI or GUI steps). Mention \textbf{Cameroon context}: ANTIC cybersecurity law (data localisation, incident reporting), MTN/Orange/CAMTEL ISP specifics, local vendor support.
\end{enumerate}
\end{methode}

\begin{exoresolu}[Integrated Case Study: Design \& Troubleshoot a School Network in Bafoussam]
\textbf{Statement:} 
``Lycée Technique de Bafoussam'' has two buildings (Admin + Classrooms) 80 m apart. Current network: One old 24-port 10/100 switch in Admin, daisy-chained to a 16-port hub in Classrooms via a single Cat5 cable. 40 admin PCs, 60 lab PCs, 5 IP cameras, 1 server (grades, library). Wi-Fi: one old router in Admin (2.4 GHz only). Problems: Frequent disconnects, slow file access, cameras drop frames, no guest Wi-Fi, students access admin shares.
The school has a new CAMTEL GPON 200 Mbps subscription (ONT in Admin). Budget: 2,000,000 FCFA. You are the network consultant.
\begin{enumerate}
\item Design a complete solution (topology, equipment, IP plan, VLANs, Wi-Fi, security) within budget. Provide a priced BoM (Bill of Materials).
\item Write the core switch configuration (CLI-style) for VLANs, Inter-VLAN routing, DHCP, and basic firewall (protect Admin VLAN from Student VLAN).
\item During commissioning, lab PCs get IP 169.254.x.x. Admin PCs work fine. Systematic troubleshooting steps to isolate and fix.
\item The principal wants to monitor student internet usage (sites visited, bandwidth). Propose a solution using existing/proposed hardware (no new budget).
\end{enumerate}
\tcblower
\textbf{Detailed Solution:}

\textbf{1. Complete Solution Design}

\textbf{Topology:} Collapsed Core (Admin Building) + Access Switch (Classroom Building) linked by Fibre.
\begin{itemize}
\item \textbf{Core Switch (Admin):} Managed L2+/L3, 24-port PoE+ (for Admin PCs, Cameras, AP, Uplink).
\item \textbf{Access Switch (Classrooms):} Managed, 48-port PoE+ (for 60 Lab PCs — use 2 switches stacked or one 48-port + one 12-port; we'll use 2× 24-port stacked for budget).
\item \textbf{Backbone:} 2-core OM3 fibre (80 m) between buildings — immune to lightning/EMI, 10G capable.
\item \textbf{Wireless:} 3× Wi-Fi 6 APs (1 Admin, 2 Classrooms) — PoE powered.
\item \textbf{Router/Firewall:} MikroTik hEX (RB750Gr3) or similar — 5× Gigabit, RouterOS, cheap, powerful.
\item \textbf{UPS:} 1500 VA for Core + Router + ONT in Admin; 1000 VA for Classroom switch stack.
\end{itemize}

\textbf{VLAN Plan:}

\begin{center}
\begin{tabular}{|l|l|l|l|}\hline
\rowcolor{popblueL}
\textbf{VLAN ID} & \textbf{Name} & \textbf{Subnet} & \textbf{Devices} \\ \hline
10 & ADMIN & 10.10.10.0/24 & Admin PCs, Server, Cameras, Admin AP (SSID: Admin\_WiFi) \\ \hline
20 & LABS & 10.10.20.0/24 & Lab PCs, Classroom APs (SSID: Lab\_WiFi) \\ \hline
30 & GUEST & 10.10.30.0/24 & Guest SSID (Admin\_Guest, Lab\_Guest) — Internet only \\ \hline
40 & CAMERA & 10.10.40.0/24 & IP Cameras (isolated, no internet) \\ \hline
99 & MGMT & 10.10.99.0/24 & Switches, Router, APs, Server (iDRAC/iLO) \\ \hline
\end{tabular}
\end{center}

\textbf{Bill of Materials (Estimated FCFA, Yaoundé/Bafoussam prices):}

\begin{center}
\begin{tabularx}{\linewidth}{|l|X|c|l|}\hline
\rowcolor{popblueL}
\textbf{Item} & \textbf{Model} & \textbf{Qty} & \textbf{Total (FCFA)} \\ \hline
Core Switch (L3, 24p PoE+) & MikroTik CRS328-24P-4S+RM & 1 & 550,000 \\ \hline
Access Switches (24p PoE+, stackable) & MikroTik CRS326-24G-2S+RM (non-PoE) + PoE Injectors? Better: \textbf{Ubiquiti USW-Lite-24-PoE} (24p PoE+) & 2 & 2 × 380,000 = 760,000 \\ \hline
Wi-Fi 6 APs & Ubiquiti U6-Lite & 3 & 3 × 150,000 = 450,000 \\ \hline
Router/Firewall & MikroTik hEX (RB750Gr3) & 1 & 85,000 \\ \hline
Fibre: OM3 2-core outdoor (100m) + termination & Local supplier & 1 lot & 120,000 \\ \hline
SFP+ Modules (10G SR) & FS.com compatible & 4 & 4 × 10,000 = 40,000 \\ \hline
Cat6A Cable (305m) & & 3 boxes & 3 × 85,000 = 255,000 \\ \hline
Patch Panels (24p Cat6A) & & 2 & 2 × 45,000 = 90,000 \\ \hline
Wall Plates + Keystones & & 110 & 110 × 2,500 = 275,000 \\ \hline
UPS (1500VA + 1000VA) & Eaton 5P / APC Smart-UPS & 2 & 350,000 \\ \hline
Cabinet (6U Admin, 6U Classroom) & Local & 2 & 2 × 80,000 = 160,000 \\ \hline
Consumables (labels, ties, tools) & & 1 lot & 100,000 \\ \hline
\multicolumn{3}{|r|}{\textbf{TOTAL}} & \textbf{3,135,000} \\ \hline
\end{tabularx}
\end{center}

\textbf{Budget Alert:} 3.135M $>$ 2M. \textbf{Adjustments to fit 2M FCFA:}
\begin{itemize}
\item Core Switch: Use \textbf{MikroTik CRS326-24G-2S+RM} (non-PoE, 240k) + separate PoE injector for cameras/AP (4 ports × 15k = 60k). Save 250k.
\item Access Switches: Use \textbf{TP-Link TL-SG2428P} (24-port PoE+, L2+, 180k each) × 2 = 360k. Save 400k.
\item APs: \textbf{TP-Link EAP610 (Wi-Fi 6)} × 3 = 3 × 110k = 330k. Save 120k.
\item Fibre: Use \textbf{single-mode OS2} (cheaper per metre) + 1G SFP (LX) instead of 10G. 1G is enough for 200 Mbps ISP + internal traffic. Save 80k.
\item UPS: 1× 1500VA (Admin) only; Classroom switch on surge protector (mains stable-ish). Save 150k.
\item Cable: Cat6 (not 6A) — 65k/box. Save 60k.
\end{itemize}
\textbf{Revised Total $\approx$ 1,950,000 FCFA.} Fits budget.

\textbf{2. Core Switch Configuration (MikroTik RouterOS / SwitchOS CLI)}

\begin{center}
\begin{tabular}{|p{0.9\linewidth}|}\hline
\texttt{\# --- VLANs on Bridge ---}\\
\texttt{/interface bridge}\\
\texttt{add name=bridge1 vlan-filtering=yes}\\
\texttt{}\\
\texttt{/interface bridge vlan}\\
\texttt{add bridge=bridge1 tagged=sfp-sfpplus1,sfp-sfpplus2,ether1-ether24 vlan-ids=10,20,30,40,99}\\
\texttt{}\\
\texttt{\# --- Access Ports (untagged) ---}\\
\texttt{/interface bridge port}\\
\texttt{add bridge=bridge1 interface=ether1 pvid=10  \# Admin PC 1}\\
\texttt{add bridge=bridge1 interface=ether2 pvid=10  \# Admin PC 2}\\
\texttt{... (repeat for all admin ports)}\\
\texttt{add bridge=bridge1 interface=ether20 pvid=40 \# Camera 1}\\
\texttt{add bridge=bridge1 interface=ether21 pvid=40 \# Camera 2}\\
\texttt{...}\\
\texttt{add bridge=bridge1 interface=ether24 pvid=99 \# Mgmt port to Router}\\
\texttt{}\\
\texttt{\# --- Trunk to Classroom Switch (sfp-sfpplus1) ---}\\
\texttt{/interface bridge port}\\
\texttt{add bridge=bridge1 interface=sfp-sfpplus1}\\
\texttt{}\\
\texttt{\# --- IP Addresses on VLAN Interfaces (L3 Routing) ---}\\
\texttt{/interface vlan}\\
\texttt{add name=vlan10 interface=bridge1 vlan-id=10}\\
\texttt{add name=vlan20 interface=bridge1 vlan-id=20}\\
\texttt{add name=vlan30 interface=bridge1 vlan-id=30}\\
\texttt{add name=vlan40 interface=bridge1 vlan-id=40}\\
\texttt{add name=vlan99 interface=bridge1 vlan-id=99}\\
\texttt{}\\
\texttt{/ip address}\\
\texttt{add address=10.10.10.1/24 interface=vlan10 comment="Admin GW"}\\
\texttt{add address=10.10.20.1/24 interface=vlan20 comment="Labs GW"}\\
\texttt{add address=10.10.30.1/24 interface=vlan30 comment="Guest GW"}\\
\texttt{add address=10.10.40.1/24 interface=vlan40 comment="Camera GW"}\\
\texttt{add address=10.10.99.1/24 interface=vlan99 comment="Mgmt GW"}\\
\texttt{}\\
\texttt{\# --- DHCP Servers ---}\\
\texttt{/ip pool}\\
\texttt{add name=pool10 ranges=10.10.10.100-10.10.10.200}\\
\texttt{add name=pool20 ranges=10.10.20.100-10.10.20.200}\\
\texttt{add name=pool30 ranges=10.10.30.100-10.10.30.200}\\
\texttt{}\\
\texttt{/ip dhcp-server}\\
\texttt{add name=dhcp10 interface=vlan10 address-pool=pool10 lease-time=1h}\\
\texttt{add name=dhcp20 interface=vlan20 address-pool=pool20 lease-time=1h}\\
\texttt{add name=dhcp30 interface=vlan30 address-pool=pool30 lease-time=30m}\\
\texttt{}\\
\texttt{/ip dhcp-server network}\\
\texttt{add address=10.10.10.0/24 gateway=10.10.10.1 dns-server=1.1.1.1,8.8.8.8}\\
\texttt{add address=10.10.20.0/24 gateway=10.10.20.1 dns-server=1.1.1.1,8.8.8.8}\\
\texttt{add address=10.10.30.0/24 gateway=10.10.30.1 dns-server=1.1.1.1,8.8.8.8}\\
\texttt{}\\
\texttt{\# --- FIREWALL: Protect Admin from Labs/Guest ---}\\
\texttt{/ip firewall filter}\\
\texttt{add chain=forward action=accept connection-state=established,related comment="Allow established"}\\
\texttt{add chain=forward action=drop connection-state=invalid}\\
\texttt{add chain=forward action=accept in-interface=vlan10 out-interface=vlan20 comment="Admin -> Labs (if needed)"}\\
\texttt{add chain=forward action=drop in-interface=vlan20 out-interface=vlan10 comment="BLOCK Labs -> Admin"}\\
\texttt{add chain=forward action=drop in-interface=vlan30 out-interface=vlan10 comment="BLOCK Guest -> Admin"}\\
\texttt{add chain=forward action=drop in-interface=vlan30 out-interface=vlan20 comment="BLOCK Guest -> Labs"}\\
\texttt{add chain=forward action=drop in-interface=vlan40 comment="BLOCK Cameras -> anywhere"}\\
\texttt{add chain=forward action=accept in-interface=vlan10 out-interface=wan comment="Admin -> Internet"}\\
\texttt{add chain=forward action=accept in-interface=vlan20 out-interface=wan comment="Labs -> Internet"}\\
\texttt{add chain=forward action=accept in-interface=vlan30 out-interface=wan comment="Guest -> Internet"}\\
\texttt{add chain=forward action=drop comment="Drop all other inter-VLAN"}\\\hline
\end{tabular}
\end{center}

\textbf{3. Troubleshooting: Lab PCs get 169.254.x.x (APIPA)}

\textbf{Symptom:} Lab PCs (VLAN 20) fail to get DHCP. Admin PCs (VLAN 10) work.

\textbf{Systematic Steps:}

\begin{enumerate}
\item \textbf{Verify Physical Link:} Check link lights on Classroom Access Switch ports for Lab PCs. \textbf{Result:} Lights ON (Green/Amber) $\rightarrow$ Layer 1 OK.
\item \textbf{Verify VLAN Assignment on Access Ports:} On Classroom Switch, check `interface bridge port` for lab ports: `pvid=20`. \textbf{Result:} Correct.
\item \textbf{Verify Trunk to Core:} On Classroom Switch, uplink port (sfp1) must be \textbf{tagged} member of VLAN 20 (and 10,30,40,99). Check `/interface bridge vlan`. \textbf{Result:} Missing VLAN 20 on trunk! Only VLAN 10,99 tagged.
\item \textbf{Fix:} Add VLAN 20 as tagged on trunk port (sfp1) on \textbf{both} Classroom Switch and Core Switch.
      \begin{center}
      \begin{tabular}{|p{0.8\linewidth}|}\hline
      \texttt{/interface bridge vlan add bridge=bridge1 tagged=sfp-sfpplus1 vlan-ids=20}\\\hline
      \end{tabular}
      \end{center}
\item \textbf{Verify DHCP Relay (if DHCP server on Router, not Core):} In this design, Core Switch \textbf{is} the DHCP server (L3). So no relay needed. But if Router was DHCP server, we'd need `/ip dhcp-relay` on Core for vlan20 pointing to Router IP.
\item \textbf{Test:} Release/renew on Lab PC (`ipconfig /release \&\& ipconfig /renew`). PC gets 10.10.20.105. Ping 10.10.20.1 (GW) $\rightarrow$ OK. Ping 10.10.10.10 (Server) $\rightarrow$ OK (firewall allows Lab$\rightarrow$Admin? No, firewall blocks Lab$\rightarrow$Admin. But ping to Server (VLAN10) from Lab (VLAN20) is blocked by design. That's correct. Ping 8.8.8.8 $\rightarrow$ OK.)
\item \textbf{Document:} ``Trunk port missing VLAN 20 tag on both ends. Added. DHCP now works.''
\end{enumerate}

\textbf{4. Internet Usage Monitoring (No New Budget)}

\textbf{Proposal:} Use \textbf{MikroTik Router (hEX) + Syslog + NetFlow/Traffic-Flow} + \textbf{Free Open-Source Analyzer}.

\begin{enumerate}
\item \textbf{Enable Traffic-Flow (NetFlow v9/IPFIX) on MikroTik:}
      \begin{center}
      \begin{tabular}{|p{0.8\linewidth}|}\hline
      \texttt{/ip traffic-flow}\\
      \texttt{set enabled=yes interfaces=all cache-entries=4k}\\
      \texttt{/ip traffic-flow target}\\
      \texttt{add address=10.10.99.10 (Syslog/Collector VM) port=2055 version=9}\\\hline
      \end{tabular}
      \end{center}
\item \textbf{Deploy a Collector/Analyzer VM on the existing Server (Ubuntu):}
      \begin{itemize}
      \item Install \textbf{ntopng} (free community edition) or \textbf{Elastic Stack (Filebeat + Logstash + Kibana)} or \textbf{Graylog}.
      \item ntopng listens on NetFlow (port 2055) and/or sFlow. Provides web dashboard: Top talkers, applications (nDPI), websites (SNI from TLS), bandwidth per VLAN/IP.
      \item Configure ntopng to map VLAN 20 IPs to ``Students'', VLAN 10 to ``Admin''.
      \end{itemize}
\item \textbf{DNS Logging (for domain names):} On MikroTik, enable DNS query logging to same syslog server:
      \begin{center}
      \begin{tabular}{|p{0.8\linewidth}|}\hline
      \texttt{/system logging action}\\
      \texttt{add name=remote\_syslog target=remote remote=10.10.99.10 remote-port=514}\\
      \texttt{/system logging}\\
      \texttt{add action=remote\_syslog topics=dns}\\\hline
      \end{tabular}
      \end{center}
\item \textbf{HTTPS Inspection (Optional, advanced):} MikroTik can do TLS Host (SNI) extraction without full MITM. ntopng shows SNI. For full URL, need MITM proxy (Squid + SSL Bump) — heavy for hEX CPU, not recommended.
\item \textbf{Reporting:} Schedule weekly PDF reports from ntopng to principal's email.
\item \textbf{Cost:} 0 FCFA (uses existing Server RAM/CPU, free software). Meets requirement.
\end{enumerate}

\textbf{Examiner's Note:} This solution demonstrates \textbf{integration} of all chapter topics: topology choice (collapsed core), equipment selection (budget constraints), VLAN design, inter-VLAN routing, DHCP, firewall segmentation, fibre backbone, Wi-Fi 6, troubleshooting methodology (Layer 1$\rightarrow$3), and monitoring with existing gear. Always mention \textbf{ANTIC} compliance: logging must respect privacy laws — inform users via captive portal / Acceptable Use Policy.
\end{exoresolu}

\begin{exercice}[Practice Questions for GCE Revision]
\begin{itemize}
\item A company has two buildings 150 m apart. They need to connect their LANs. Which media is most appropriate and why? (Consider cost, bandwidth, immunity to interference.)
\item Draw a labelled diagram of a \textbf{Tree topology} for a 3-floor building with 30 PCs per floor. Show core, distribution, access layers, uplink media, and VLAN segmentation for Admin, Staff, Guest, and IoT.
\item A switch has 48 ports, 370W PoE budget. You need to power 24 APs (15W each), 10 IP phones (7W each), and 5 cameras (12W each). Is the PoE budget sufficient? Show calculation.
\item Explain the difference between a \textbf{Layer 2 switch} and a \textbf{Layer 3 switch}. When would you use each in a school network?
\item A student connects a home router (LAN port) to a wall jack in the school lab, creating a rogue DHCP server. What switch feature prevents this? Configure it on a MikroTik switch port.
\item Your school's CAMTEL GPON ONT has 4 Ethernet ports. Port 1 is used for the main router. Can you plug a second router into Port 2 for a separate network? Explain the constraints (PPPoE sessions, VLANs on ONT).
\item Design a Wi-Fi channel plan for a 2.4 GHz and 5 GHz deployment in a 4-classroom block (20 m × 40 m). Show channel numbers, widths, and AP placement.
\item A network printer has a static IP 192.168.10.50/24. The DHCP pool is 192.168.10.100–200. A teacher's laptop gets 192.168.10.50 via DHCP and causes an IP conflict. Why did this happen? How to prevent it permanently?
\item Write a \textbf{firewall rule set} (MikroTik style) that: (a) allows Admin VLAN to access Server VLAN on TCP 443 (HTTPS) and TCP 22 (SSH); (b) blocks Student VLAN from Admin VLAN entirely; (c) allows both to access Internet; (d) logs all dropped packets.
\item List 5 items that must appear in a \textbf{network handover document} for a Cameroonian SME.
\end{itemize}
\end{exercice}

\begin{corrige}[Exercise 1]
\textbf{Media:} \textbf{Multimode Fibre (OM3/OM4)} or \textbf{Single-mode Fibre (OS2)}.
\textbf{Why:} 150 m exceeds UTP 100 m limit. Fibre is immune to EMI/lightning (critical in Cameroon). OM3 supports 10G up to 300 m; OS2 supports 10G/100G up to 10 km. Cost is moderate (fibre + 2 media converters or SFP modules). Wireless bridge (60 GHz or 5 GHz) is an alternative but requires line-of-sight and is weather-sensitive.
\end{corrige}

\begin{corrige}[Exercise 3]
\textbf{Calculation:}
\begin{align*}
\text{APs:} & \quad 24 \times 15\text{W} = 360\text{W} \\
\text{Phones:} & \quad 10 \times 7\text{W} = 70\text{W} \\
\text{Cameras:} & \quad 5 \times 12\text{W} = 60\text{W} \\
\text{Total:} & \quad 360 + 70 + 60 = 490\text{W}
\end{align*}
\textbf{Result:} 490W $>$ 370W budget. \textbf{Insufficient.} Solutions: (1) Use PoE+ (30W/port) switch with 720W budget (e.g., USW-Pro-48-PoE 600W — still tight). (2) Use PoE++ (60W/port) or injectors for cameras. (3) Reduce AP count or use lower-power APs. (4) Power cameras via separate PoE switch/injectors.
\end{corrige}

\begin{corrige}[Exercise 5]
\textbf{Feature:} \textbf{DHCP Snooping} (with \textbf{Trusted/Untrusted ports}).
\textbf{Configuration (MikroTik):}
\begin{center}
\begin{tabular}{|p{0.8\linewidth}|}\hline
\texttt{/interface bridge port}\\
\texttt{set [find interface=ether1] trusted=yes \# Uplink to core (trusted)}\\
\texttt{set [find interface=ether2] trusted=no \# Student port (untrusted)}\\
\texttt{/ip dhcp-server}\\
\texttt{set dhcp1 dhcp-snooping=yes}\\\hline
\end{tabular}
\end{center}
Only trusted ports can send DHCP OFFER/ACK. Rogue DHCP on untrusted port is blocked.
\end{corrige}

\begin{corrige}[Exercise 8]
\textbf{Cause:} The printer's static IP (192.168.10.50) is \textbf{inside the DHCP pool range} (100–200? No, 50 is outside 100–200. Wait — the pool is .100–.200, so .50 is outside. But the teacher got .50. This means the DHCP server \textbf{assigned an address outside its pool} — impossible unless the pool was misconfigured (e.g., .50–.200) or the printer's IP was not excluded.
\textbf{Correction:} The DHCP pool must \textbf{exclude} all static IPs. Define pool as 192.168.10.100–200 \textbf{AND} add exclusion for .10 (server), .11 (printer), .50 (printer if different). Better: use a separate subnet for infrastructure (VLAN 99) or reserve via DHCP reservation (MAC-based) instead of static on device.
\textbf{Prevention:} (1) Keep static IPs \textbf{outside} DHCP pool. (2) Use \textbf{DHCP reservations} (MAC $\rightarrow$ IP) on the server — single source of truth. (3) Enable \textbf{conflict detection} (ping before offer) on DHCP server.
\end{corrige}

\begin{corrige}[Exercise 9]
\begin{center}
\begin{tabular}{|p{0.9\linewidth}|}\hline
\texttt{/ip firewall filter}\\
\texttt{add chain=forward action=accept connection-state=established,related}\\
\texttt{add chain=forward action=drop connection-state=invalid log=yes log-prefix="INVALID"}\\
\texttt{\# Admin -> Server (HTTPS, SSH)}\\
\texttt{add chain=forward action=accept in-interface=vlan10 out-interface=vlan50 protocol=tcp dst-port=443,22}\\
\texttt{\# Block Student -> Admin}\\
\texttt{add chain=forward action=drop in-interface=vlan20 out-interface=vlan10 log=yes log-prefix="STU->ADM"}\\
\texttt{\# Admin -> Internet}\\
\texttt{add chain=forward action=accept in-interface=vlan10 out-interface=wan}\\
\texttt{\# Student -> Internet}\\
\texttt{add chain=forward action=accept in-interface=vlan20 out-interface=wan}\\
\texttt{\# Drop all other inter-VLAN}\\
\texttt{add chain=forward action=drop log=yes log-prefix="DROP-ALL"}\\\hline
\end{tabular}
\end{center}
Order matters: specific allows before general drops. Logging on drops for audit (ANTIC).
\end{corrige}

\begin{corrige}[Exercise 10]
\begin{enumerate}
\item Network topology diagram (physical + logical) with IP/VLAN/port mapping.
\item Device inventory: Model, Serial, Firmware, Management IP, Licenses.
\item Cable test reports (certification) + cable map (patch panel $\leftrightarrow$ wall plate).
\item Configuration backups (encrypted) + admin credentials (sealed envelope).
\item Maintenance schedule: firmware updates, UPS battery tests, log review, password rotation, ANTIC compliance checklist.
\item\item Acceptable Use Policy (AUP) signed by users + Wi-Fi captive portal terms.
\end{enumerate}
\end{corrige}

\begin{aretenir}[Chapter ICT05 — Final Revision Checklist]
\begin{itemize}
\item \textbf{Topologies:} Identify 6 physical topologies; know SPOF, cable count, use cases. Star = baseline; Bus/Ring = legacy; Mesh = core redundancy; Tree/Hybrid = real world.
\item \textbf{Comparison:} Build tables (Cost, Scalability, Fault Tolerance, Performance, Maintenance). Justify choice with \emph{at least two criteria} from scenario.
\item \textbf{Equipment:} Classify by OSI layer (L1: cables, L2: switch/AP, L3: router/firewall). Hub vs Switch vs Router vs Modem vs AP — know the difference.
\item \textbf{Setup:} 4 phases (Plan, Physical, Logical, Test/Document). 568B termination, 100 m limit, UPS, labels, change default passwords, VLANs, DHCP, Firewall, Guest isolation.
\item \textbf{Selection:} Requirements $\rightarrow$ Blocks (Core/Access) $\rightarrow$ Models (ports, PoE, uplink 10×, management). Budget in FCFA, local vendor support, redundancy (dual PSU, stacking, dual ISP).
\item \textbf{Integration:} Design loop (L1$\rightarrow$L7), Troubleshooting (Layer-by-layer), Monitoring (NetFlow/Syslog + free tools), ANTIC compliance (logs, privacy, incident reporting).
\item \textbf{Cameroon Context:} CAMTEL GPON (PPPoE), MTN/Orange Fibre/4G, ANTIC Law 2010/012, local prices (FCFA), lightning/EMI $\rightarrow$ fibre, ENEO power $\rightarrow$ UPS.
\end{itemize}
\end{aretenir}

\newpage

\coursec{Exercises}

\courssub{I check my knowledge}

\begin{exercice}[MCQ: Topology identification]
Which physical topology connects every node directly to a central device?
\begin{enumerate}
\item Bus
\item Ring
\item Star
\item Mesh
\end{enumerate}
\end{exercice}

\begin{exercice}[True/False with justification]
State whether each statement is true or false and justify your answer.
\begin{enumerate}
\item In a full mesh topology with $n$ nodes, the number of links required is $n(n-1)$.
\item A switch operates at the Physical Layer (Layer 1) of the OSI model.
\item A router connects two or more logical networks and forwards packets based on IP addresses.
\item UTP Cat 6 cable supports a maximum data rate of 100 Mbps over 100 m.
\end{enumerate}
\end{exercice}

\begin{exercice}[Definitions]
Define the following terms in the context of computer networks:
\begin{enumerate}
\item Network topology
\item Collision domain
\item Broadcast domain
\item Straight-through cable (T568B)
\item Patch panel
\end{enumerate}
\end{exercice}

\begin{exercice}[Direct calculation]
A small business decides to cable its LAN as a full mesh using UTP cables. If there are 6 computers:
\begin{enumerate}
\item Calculate the total number of cable links required.
\item If each link costs 2\,500 FCFA, what is the total cabling cost?
\item How many ports must each network interface card (NIC) have?
\end{enumerate}
\end{exercice}

\courssub{I practise}

\begin{exercice}[Diagram analysis]
The diagram below shows a network installed in a secondary school in Yaoundé. 
\begin{center}
\begin{tikzpicture}[scale=0.8, every node/.style={font=\small}]
% Central switch
\node[draw, circle, fill=popblue, text=white, minimum size=1.2cm] (SW) at (0,0) {Switch};
% PCs
\foreach \i/\angle in {1/90, 2/30, 3/-30, 4/-90, 5/-150, 6/150} {
    \node[draw, rectangle, fill=popgreen, text=white, minimum width=1.5cm, minimum height=0.8cm] (PC\i) at (\angle:3) {PC\i};
    \draw[thick, popdark] (SW) -- (PC\i);
}
% Router
\node[draw, rectangle, fill=poporange, text=white, minimum width=1.5cm, minimum height=0.8cm] (R) at (0,-5) {Router};
\draw[thick, popdark] (SW) -- (R);
% Internet cloud (replaced by ellipse to avoid shapes.symbols library)
\node[draw, ellipse, fill=poppink, minimum width=2.5cm, minimum height=1.5cm] (NET) at (0,-7) {Internet};
\draw[thick, popdark] (R) -- (NET);
\end{tikzpicture}
\end{center}
\begin{enumerate}
\item Identify the physical topology of the LAN segment connecting PC1–PC6.
\item Name the devices labelled ``Switch'' and ``Router'' and state the OSI layer at which each primarily operates.
\item State one advantage and one disadvantage of the identified topology for this school environment.
\item If PC3 fails, what is the impact on communication between PC1 and PC5? Justify.
\end{enumerate}
\end{exercice}

\begin{exercice}[Cabling standards and testing]
A technician in Douala is asked to produce a straight-through UTP Cat 5e cable terminated with RJ45 connectors using the T568B wiring standard.
\begin{enumerate}
\item Write the colour order of the 8 wires from pin 1 to pin 8 for T568B.
\item Explain the difference between a straight-through cable and a crossover cable. Give one typical use case for each.
\item After crimping, the technician tests the cable with a cable tester. The tester shows ``Open'' on pin 3. List three possible physical causes.
\item What does a ``Short'' indication on the tester mean?
\end{enumerate}
\end{exercice}

\begin{exercice}[Equipment matching]
Match each device (A–J) to its correct description (1–10). Write your answers as pairs (e.g., A–3).
\begin{center}
\begin{tabularx}{\linewidth}{|l|X|}\hline
\textbf{Device} & \textbf{Description} \\ \hline
A. NIC & 1. Connects multiple network segments and filters traffic based on MAC addresses \\ \hline
B. Hub & 2. Provides wireless connectivity to wired LAN \\ \hline
C. Switch & 3. Interconnects different networks and routes packets by IP address \\ \hline
D. Router & 4. Amplifies and regenerates signals to extend cable length \\ \hline
E. Modem & 5. Interface between computer and network medium \\ \hline
F. Access Point & 6. Converts digital signals to analogue for transmission over telephone lines \\ \hline
G. Repeater & 7. Terminates cables in a wiring closet for organisation \\ \hline
H. Bridge & 8. Multiport repeater; broadcasts all frames to every port \\ \hline
I. Patch Panel & 9. Connects two LAN segments and filters by MAC address \\ \hline
J. Cable Tester & 10. Verifies continuity, wiring order and faults in cables \\ \hline
\end{tabularx}
\end{center}
\end{exercice}

\begin{exercice}[Troubleshooting scenario]
After setting up a 10‑computer lab network in a Star topology (Cat 6, Gigabit switch), PC7 cannot communicate with any other PC. The link LED on the switch port for PC7 is off. List four distinct possible physical‑layer causes and, for each, the specific test you would perform to confirm or eliminate it.
\end{exercice}

\courssub{I prepare for the GCE}

\begin{exercice}[GCE Paper 2 style: Topology design and justification — 15 marks]
The Ministry of Public Health plans to network a new 3‑storey building in Yaoundé. Each floor has 10 offices. Each office requires 2 wired data points (for PC and VoIP phone) and wireless coverage. The network must:
\begin{itemize}
\item Guarantee that failure of any single office's equipment does not affect other offices.
\item Support future upgrade to 10 Gbps backbone.
\item Allow centralised management and VLAN segmentation per department.
\item Stay within a reasonable budget for a public institution.
\end{itemize}
\begin{enumerate}
\item Propose a physical topology (or combination) for the LAN. Draw a labelled diagram showing the topology for one floor and the inter‑floor backbone. [4 marks]
\item Recommend the central connectivity equipment for the floor and the core, justifying your choice with reference to collision domains, broadcast domains and manageability. [4 marks]
\item Specify the cable category you would use for horizontal cabling and for the backbone. Justify each choice with bandwidth and distance considerations. [3 marks]
\item Explain how you would implement redundancy for the core switch to meet the high‑availability requirement. [2 marks]
\item List two Cameroonian regulatory or standards bodies whose guidelines should be consulted during installation. [2 marks]
\end{enumerate}
\end{exercice}

\begin{exercice}[GCE Paper 2 style: Equipment selection and cost analysis — 12 marks]
A startup in Buea needs to connect 25 workstations in a single large room. Two proposals are made:
\begin{itemize}
\item \textbf{Proposal A:} One 24‑port unmanaged Gigabit switch + one 8‑port unmanaged Gigabit switch (cascaded).
\item \textbf{Proposal B:} One 48‑port managed Gigabit switch with PoE+ support.
\end{itemize}
\begin{enumerate}
\item For each proposal, calculate the number of available ports for workstations after inter‑switch links (if any). [2 marks]
\item Compare the two proposals in terms of: (i) collision domains, (ii) broadcast domains, (iii) VLAN capability, (iv) QoS for future VoIP, (v) power budget for PoE devices. Present as a table. [5 marks]
\item The startup expects to add 15 more workstations within 12 months. Which proposal scales better? Justify. [2 marks]
\item Assuming the following unit prices in Cameroon: unmanaged 24‑port = 85\,000 FCFA, unmanaged 8‑port = 25\,000 FCFA, managed 48‑port PoE+ = 420\,000 FCFA. Calculate the initial equipment cost for each proposal. [1 mark]
\item Beyond cost, give two technical reasons why the ministry might still prefer Proposal B for a similar deployment. [2 marks]
\end{enumerate}
\end{exercice}

\begin{exercice}[GCE Paper 2 style: Practical configuration and troubleshooting — 13 marks]
You are given a network simulation (Packet Tracer or physical lab) with the following addressing scheme:
\begin{itemize}
\item LAN A: 192.168.10.0/24, gateway 192.168.10.1
\item LAN B: 192.168.20.0/24, gateway 192.168.20.1
\item Serial link between routers: 10.0.0.0/30
\end{itemize}
The routers are connected via a serial DCE/DTE cable. PC1 (LAN A) cannot ping PC2 (LAN B).
\begin{enumerate}
\item List the sequence of commands you would enter on Router A to configure its LAN interface, serial interface (DCE side), and a static route to LAN B. Assume Cisco IOS syntax. [5 marks]
\item The serial link shows ``Protocol down''. Name three Layer 1/2 issues that could cause this and the command to verify each. [3 marks]
\item After fixing the serial link, PC1 pings PC2 successfully but cannot browse a web server on PC2 (port 80). The Windows Firewall on PC2 is off. Suggest two possible causes at Layer 4 or above and how to test each. [3 marks]
\item Write the T568B pinout colour order used for the straight-through patch cables connecting PCs to the access switch. [2 marks]
\end{enumerate}
\end{exercice}

\newpage

\coursec{Solutions to the exercises}

\coursec{I check my knowledge}

\begin{corrige}[Exercise 1 — MCQ: Topology identification]
\textbf{Correct answer: item 3 — Star.}\par
In a \cle{star topology}, every node (computer, printer, etc.) is connected by its own dedicated cable directly to a \cle{central device} (a hub or, in modern networks, a switch). All communication between two nodes passes through this central device.
\begin{itemize}
\item \textbf{Bus} is wrong: in a bus, all nodes share a single common backbone cable terminated at both ends; there is no central device.
\item \textbf{Ring} is wrong: each node is connected to exactly two neighbours, forming a closed loop; there is no central device.
\item \textbf{Mesh} is wrong: nodes are interconnected directly with each other (point-to-point links), not through a central device.
\end{itemize}
\end{corrige}

\begin{corrige}[Exercise 2 — True/False with justification]
\begin{enumerate}
\item \textbf{False.} In a full mesh, each of the $n$ nodes connects to the $n-1$ other nodes, which gives $n(n-1)$ ordered connections; but each physical link is counted twice (the link A–B is the same as B–A). The correct number of links is
\[ L = \frac{n(n-1)}{2}. \]
\item \textbf{False.} A switch operates primarily at the \cle{Data Link Layer (Layer 2)} of the OSI model: it reads MAC addresses to forward frames only to the destination port. A device working at Layer 1 (Physical) is a hub or a repeater, which simply regenerates electrical signals.
\item \textbf{True.} A router operates at Layer 3 (Network layer). It interconnects two or more distinct logical networks (different IP networks/subnets) and forwards packets by examining the destination \cle{IP address} and consulting its routing table.
\item \textbf{False.} UTP Cat 6 supports up to \SI{1}{Gbps} (1000 Mbps) over the full \SI{100}{m} channel, and even \SI{10}{Gbps} over shorter distances (up to about \SI{55}{m}). The 100 Mbps limit over 100 m corresponds to older categories such as Cat 5 in Fast Ethernet (100BASE-TX).
\end{enumerate}
\end{corrige}

\begin{corrige}[Exercise 3 — Definitions]
\begin{enumerate}
\item \textbf{Network topology:} the arrangement or layout of the nodes (computers, devices) and the communication links in a network. The \emph{physical} topology describes how devices are actually cabled and interconnected; the \emph{logical} topology describes how data actually flows between them.
\item \textbf{Collision domain:} the set of network segments (or ports) within which two frames transmitted simultaneously can collide. On a hub, all ports form one single collision domain; on a switch, each port is a separate collision domain.
\item \textbf{Broadcast domain:} the set of all devices that receive a broadcast frame (e.g., an ARP request) sent by any one of them. Switches forward broadcasts to all ports, so a switched LAN is one broadcast domain; routers and VLANs delimit separate broadcast domains.
\item \textbf{Straight-through cable (T568B):} a UTP patch cable in which both RJ45 connectors are wired with the same pinout (here T568B on both ends), so pin 1 connects to pin 1, pin 2 to pin 2, etc. It is used to connect \emph{unlike} devices, e.g., a PC to a switch.
\item \textbf{Patch panel:} a passive panel, usually mounted in a rack in a wiring closet, on which the horizontal cabling runs from the offices are terminated (punched down). Short patch cords then connect the patch panel ports to the switch ports, making administration, labelling and re-patching easy.
\end{enumerate}
\end{corrige}

\begin{corrige}[Exercise 4 — Direct calculation]
\begin{enumerate}
\item In a full mesh with $n$ nodes, the number of links is $L = \dfrac{n(n-1)}{2}$. For $n=6$:
\[ L = \frac{6 \times 5}{2} = \frac{30}{2} = 15 \text{ links.} \]
\item Total cabling cost:
\[ 15 \times 2\,500 = 37\,500 \text{ FCFA.} \]
\item In a full mesh, each computer connects directly to the $n-1 = 5$ other computers, so each NIC must have \textbf{5 ports} (or equivalently each PC needs 5 network interfaces). This is precisely why full mesh is impractical for LANs and is reserved for small backbone interconnections.
\end{enumerate}
\end{corrige}

\coursec{I practise}

\begin{corrige}[Exercise 5 — Diagram analysis]
\begin{enumerate}
\item The LAN segment PC1–PC6 is a \cle{star topology}: each PC has its own dedicated link to the central switch.
\item \textbf{Switch:} a Layer 2 (Data Link) device that forwards Ethernet frames based on MAC addresses, creating one collision domain per port. \textbf{Router:} a Layer 3 (Network) device that interconnects different IP networks (here the school LAN and the Internet) and forwards packets based on IP addresses.
\item \textbf{Advantage:} easy administration and fault isolation — a failure of one PC or one cable affects only that PC; adding or removing a station is simple and does not disturb the others. \textbf{Disadvantage:} the central switch is a \emph{single point of failure}: if it breaks down, the whole LAN stops; also, cabling cost is higher than a bus because each node needs its own cable run.
\item \textbf{No impact.} PC1 and PC5 communicate through the switch using their own dedicated links; the link to PC3 is independent. The failure of PC3 (or its cable) does not interrupt traffic between any other pair of stations. This is the main resilience advantage of the star over bus and ring topologies.
\end{enumerate}
\end{corrige}

\begin{corrige}[Exercise 6 — Cabling standards and testing]
\begin{enumerate}
\item T568B colour order, pin 1 to pin 8:
\begin{center}
\begin{tabularx}{\linewidth}{|c|X|}\hline
\rowcolor{popblueL}\textbf{Pin} & \textbf{Wire colour} \\ \hline
1 & White/Orange \\ \hline
2 & Orange \\ \hline
3 & White/Green \\ \hline
4 & Blue \\ \hline
5 & White/Blue \\ \hline
6 & Green \\ \hline
7 & White/Brown \\ \hline
8 & Brown \\ \hline
\end{tabularx}
\end{center}
\item A \textbf{straight-through} cable has the same standard (both T568B, or both T568A) on both ends: pin 1 to pin 1, etc. It connects \emph{unlike} devices: PC to switch, switch to router. A \textbf{crossover} cable has T568A on one end and T568B on the other, so the transmit pair of one end lands on the receive pair of the other. It connects \emph{like} devices: PC to PC, switch to switch (on ports without auto-MDIX).
\item Three possible physical causes of an ``Open'' on pin 3 (white/green wire):
\begin{itemize}
\item The white/green wire was not pushed fully to the end of the RJ45 plug before crimping, so the contact blade did not pierce it.
\item The wire is broken or cut inside the cable (damaged during pulling or kinked sharply).
\item Poor crimping: the crimp tool did not press the contact for pin 3 down properly (or a defective connector).
\end{itemize}
\item A \textbf{``Short''} means that two (or more) wires that should be insulated from each other are electrically touching somewhere along the link — e.g., crushed cable, damaged insulation, or strands touching inside a badly stripped connector. The affected pairs must be re-terminated or the cable replaced.
\end{enumerate}
\end{corrige}

\begin{corrige}[Exercise 7 — Equipment matching]
\begin{center}
\begin{tabularx}{\linewidth}{|l|l|X|}\hline
\rowcolor{popblueL}\textbf{Pair} & \textbf{Device} & \textbf{Reason} \\ \hline
A--5 & NIC & It is the hardware interface between the computer and the network medium. \\ \hline
B--8 & Hub & A multiport repeater: it regenerates and broadcasts every frame out of all ports (single collision domain). \\ \hline
C--1 & Switch & Connects multiple segments/hosts and forwards frames selectively using MAC addresses. \\ \hline
D--3 & Router & Interconnects different networks and routes packets using IP addresses. \\ \hline
E--6 & Modem & Modulates/demodulates: converts digital signals to analogue (and back) for telephone lines. \\ \hline
F--2 & Access Point & Gives wireless (Wi-Fi) stations access to the wired LAN. \\ \hline
G--4 & Repeater & Regenerates and amplifies the signal to extend the maximum cable length. \\ \hline
H--9 & Bridge & Joins two LAN segments and filters traffic by MAC address (a two-port switch ancestor). \\ \hline
I--7 & Patch Panel & Terminates horizontal cable runs in the wiring closet for tidy organisation. \\ \hline
J--10 & Cable Tester & Checks continuity, wire map and faults (open, short, miswire). \\ \hline
\end{tabularx}
\end{center}
\textbf{Answers:} A–5, B–8, C–1, D–3, E–6, F–2, G–4, H–9, I–7, J–10.
\end{corrige}

\begin{corrige}[Exercise 8 — Troubleshooting scenario]
The link LED off indicates a \cle{Layer 1 (physical)} problem. Four distinct causes and their tests:
\begin{enumerate}
\item \textbf{Faulty or badly crimped patch cable (open/short/miswire).} \emph{Test:} test the cable with a cable tester (wire map), or simply replace it with a known-good cable and check whether the LED lights.
\item \textbf{Faulty switch port.} \emph{Test:} plug PC7's cable into another switch port known to work (e.g., the port of a functioning PC). If the LED lights, the original port is defective or disabled.
\item \textbf{Faulty or disabled NIC in PC7.} \emph{Test:} check the NIC's own link LED and its status in the operating system (enabled/disabled, driver present); connect a known-good PC to PC7's cable — if the LED lights, PC7's NIC is the problem.
\item \textbf{Cable not properly seated / damaged horizontal run or wall socket.} \emph{Test:} re-seat both RJ45 connectors until they click; inspect and re-punch the wall outlet/patch panel termination, or test the permanent link with the cable tester from patch panel to outlet.
\end{enumerate}
\end{corrige}

\coursec{I prepare for the GCE}

\begin{corrige}[Exercise 9 — Topology design and justification (15 marks)]
\textbf{Indicative mark scheme and model answer.}\par
\textbf{1. Proposed topology [4 marks].} An \cle{extended star (hierarchical star)}: on each floor, an \emph{extended star} with a floor switch (access switch) in the wiring closet, each office's 2 data points cabled individually to it; the three floor switches are linked by a backbone to a central core switch (star of stars). Wireless access points on each floor connect to the floor switch.
\begin{popfigure}
\begin{tikzpicture}[scale=0.75, every node/.style={font=\small}]
\node[draw, rectangle, fill=poporange, text=white, minimum width=2.2cm, minimum height=0.9cm] (core) at (0,0) {Core switch};
\node[draw, rectangle, fill=popblue, text=white, minimum width=2.2cm, minimum height=0.9cm] (f1) at (-5,-2.5) {Floor 1 switch};
\node[draw, rectangle, fill=popblue, text=white, minimum width=2.2cm, minimum height=0.9cm] (f2) at (0,-2.5) {Floor 2 switch};
\node[draw, rectangle, fill=popblue, text=white, minimum width=2.2cm, minimum height=0.9cm] (f3) at (5,-2.5) {Floor 3 switch};
\draw[thick, popdark] (core) -- (f1);
\draw[thick, popdark] (core) -- (f2);
\draw[thick, popdark] (core) -- (f3);
\node[draw, rectangle, fill=popgreen, text=white, minimum width=1.6cm, minimum height=0.7cm] (o1) at (-6.6,-5) {Office PCs};
\node[draw, rectangle, fill=popgreen, text=white, minimum width=1.6cm, minimum height=0.7cm] (o2) at (-4.4,-5) {VoIP phones};
\node[draw, rectangle, fill=poppink, minimum width=1.4cm, minimum height=0.7cm] (ap) at (-2.2,-5) {Wi-Fi AP};
\draw[thick, popdark] (f1) -- (o1);
\draw[thick, popdark] (f1) -- (o2);
\draw[thick, popdark] (f1) -- (ap);
\node[font=\small\itshape] at (-4.4,-6) {Repeated on each floor (10 offices, 2 points each)};
\end{tikzpicture}
\legende{Extended star (hierarchical) topology for a three-floor building}
\end{popfigure}
\emph{Marking: 1 mark for star per floor, 1 for backbone star to a core, 1 for labelled diagram, 1 for including APs and both data points per office.}
\par
\textbf{2. Central equipment [4 marks].} Use \textbf{managed switches} at both access (floor) and core levels — not hubs. Justification:
\begin{itemize}
\item Each switch port is a separate \cle{collision domain}, eliminating collisions and giving full bandwidth per office (a hub would create one large collision domain).
\item Managed switches support \cle{VLANs}, so each department gets its own \cle{broadcast domain}, improving security and reducing broadcast traffic.
\item Manageability: SNMP monitoring, port security, remote configuration — required for centralised management.
\end{itemize}
\emph{Marking: 1 mark for choosing managed switches, 1 for collision domains, 1 for broadcast domains/VLANs, 1 for manageability.}
\par
\textbf{3. Cabling [3 marks].}
\begin{itemize}
\item \textbf{Horizontal cabling (offices to floor closet):} UTP \textbf{Cat 6} (or Cat 6A) — supports 1 Gbps over the full \SI{100}{m} and 10 Gbps up to about \SI{55}{m}, cheap and easy to terminate; runs here are short (one floor).
\item \textbf{Backbone (floor to core):} \textbf{fibre optic} (e.g., OM3/OM4 multimode) — supports 10 Gbps and beyond over the inter-floor distances, immune to electromagnetic interference, and satisfies the ``future 10 Gbps backbone'' requirement. (Cat 6A copper is acceptable only if distances stay under 100 m, but fibre is the future-proof answer.)
\end{itemize}
\emph{Marking: 1 mark per correct choice with bandwidth/distance justification (3 marks total).}
\par
\textbf{4. Core redundancy [2 marks].} Deploy \textbf{two core switches} with each floor switch dual-homed (one uplink to each core), using link aggregation and a redundancy protocol (e.g., spanning tree with rapid convergence, or stacked/HSRP-style first-hop redundancy). If one core fails, traffic automatically flows through the second. \emph{(1 mark for dual core/dual uplinks, 1 for the failover mechanism.)}
\par
\textbf{5. Regulatory/standards bodies [2 marks].} Any two of:
\begin{itemize}
\item \textbf{ANTIC} (National Agency for Information and Communication Technologies) — ICT security and standards guidelines in Cameroon.
\item \textbf{ART} (Telecommunications Regulatory Board) — regulation of telecommunications installations.
\item International standards bodies also accepted if justified: \textbf{ISO/IEC} and \textbf{TIA/EIA} (structured cabling standards such as TIA-568).
\end{itemize}
\emph{(1 mark each, any two valid bodies.)}
\end{corrige}

\begin{corrige}[Exercise 10 — Equipment selection and cost analysis (12 marks)]
\textbf{1. Available ports [2 marks].}
\begin{itemize}
\item \textbf{Proposal A:} cascading uses 1 port on each switch for the inter-switch link: $(24-1)+(8-1)=23+7=\textbf{30 ports}$ for workstations.
\item \textbf{Proposal B:} no inter-switch link needed: \textbf{48 ports} available.
\end{itemize}
\textbf{2. Comparison table [5 marks — 1 per correct row].}
\begin{center}
\begin{tabularx}{\linewidth}{|l|X|X|}\hline
\rowcolor{popblueL}\textbf{Criterion} & \textbf{Proposal A (unmanaged)} & \textbf{Proposal B (managed PoE+)} \\ \hline
(i) Collision domains & One per switch port (switches), but no control; cascaded link is a bottleneck & One per port, with configurable settings \\ \hline
(ii) Broadcast domains & Single broadcast domain for all 25 PCs & Can be segmented into several via VLANs \\ \hline
(iii) VLAN capability & None (unmanaged) & Yes — departmental VLANs, security \\ \hline
(iv) QoS for VoIP & Not supported; voice and data compete equally & QoS queues prioritise voice traffic \\ \hline
(v) PoE power budget & None — phones/APs need separate power adapters & PoE+ powers phones and APs directly from the switch \\ \hline
\end{tabularx}
\end{center}
\textbf{3. Scalability [2 marks].} Proposal B scales better: $48-25=23$ free ports, enough for the 15 extra workstations with 8 spare, with no new equipment. Proposal A has only $30-25=5$ free ports, so a third switch (and another cascade hop, adding latency and a failure point) would be required. \emph{(1 mark for the port arithmetic, 1 for the conclusion.)}
\par
\textbf{4. Initial cost [1 mark].}
\begin{align*}
\text{Proposal A} &: 85\,000 + 25\,000 = 110\,000 \text{ FCFA} \\
\text{Proposal B} &: 420\,000 \text{ FCFA}
\end{align*}
\textbf{5. Two technical reasons to prefer B [2 marks — 1 each].} Any two of:
\begin{itemize}
\item \textbf{Manageability:} remote monitoring/configuration (SNMP), port security, faster fault diagnosis — essential for a public institution.
\item \textbf{VLAN segmentation:} separates departments/services into distinct broadcast domains, improving security and performance.
\item \textbf{PoE+:} powers VoIP phones and Wi-Fi access points without separate electrical installations.
\item \textbf{No cascade bottleneck/single point of failure} between two switches; one device is simpler and more reliable.
\end{itemize}
\end{corrige}

\begin{corrige}[Exercise 11 — Practical configuration and troubleshooting (13 marks)]
\textbf{1. Configuration of Router A [5 marks — 1 per logical block].}
\begin{itemize}
\item Enter privileged and global configuration modes:
\texttt{enable} then \texttt{configure terminal}
\item LAN interface (e.g., GigabitEthernet0/0):
\texttt{interface g0/0} ; \texttt{ip address 192.168.10.1 255.255.255.0} ; \texttt{no shutdown}
\item Serial interface, DCE side (must provide the clock):
\texttt{interface s0/0/0} ; \texttt{ip address 10.0.0.1 255.255.255.252} ; \texttt{clock rate 64000} ; \texttt{no shutdown}
\item Static route to LAN B (via Router B's serial address 10.0.0.2):
\texttt{ip route 192.168.20.0 255.255.255.0 10.0.0.2}
\item Save: \texttt{end} then \texttt{copy running-config startup-config}
\end{itemize}
\emph{Examiner expectations: correct subnet masks (/24 and /30), \texttt{no shutdown} on both interfaces, and \texttt{clock rate} only on the DCE side are the classic marking points.}
\par
\textbf{2. ``Protocol down'' on the serial link — three Layer 1/2 causes [3 marks — 1 each].}
\begin{itemize}
\item \textbf{Missing clock rate on the DCE side} — verify with \texttt{show controllers serial0/0/0} (shows DCE/DTE and clocking).
\item \textbf{Encapsulation mismatch} (HDLC on one side, PPP on the other) — verify with \texttt{show interfaces serial0/0/0} (line ``Encapsulation'').
\item \textbf{Interface shut down or faulty/loose serial cable} — verify with \texttt{show ip interface brief} (status ``administratively down'') and by physically inspecting/replacing the cable.
\end{itemize}
\textbf{3. Ping works but HTTP (port 80) fails — two Layer 4+ causes [3 marks].}
\begin{itemize}
\item \textbf{The web server service is not running} on PC2 (or listening on a different port). \emph{Test:} on PC2 run \texttt{netstat -an} and check for a listener on port 80, or from PC1 try \texttt{telnet 192.168.20.x 80} — connection refused confirms the service is down.
\item \textbf{An access control list (ACL) on a router blocks TCP port 80} while permitting ICMP. \emph{Test:} \texttt{show access-lists} and \texttt{show ip interface} on both routers to see applied ACLs; remove or correct the offending entry.
\end{itemize}
\emph{(1 mark per cause, 1 mark for valid tests; other valid answers, e.g., proxy misconfiguration, accepted.)}
\par
\textbf{4. T568B pinout for the straight-through patch cables [2 marks].}\par
Pin 1: White/Orange — Pin 2: Orange — Pin 3: White/Green — Pin 4: Blue — Pin 5: White/Blue — Pin 6: Green — Pin 7: White/Brown — Pin 8: Brown, identical on both ends. \emph{(2 marks for the complete correct order; 1 mark if one or two pairs are swapped.)}
\end{corrige}

\newpage

\coursec{Summary sheet}

\begin{popfigure}
\begin{tikzpicture}[scale=0.85, every node/.style={font=\small}]
% Central node
\node[circle, draw=popdark, fill=popblueL, text width=2.8cm, align=center, minimum height=1.6cm] (center) {ICT05: Network Topology and Setup};
% Branch nodes
\node[circle, draw=popblue, fill=popblue!20, text width=2.2cm, align=center, minimum height=1.2cm] (b1) at (90:4.2cm) {Physical Topologies};
\node[circle, draw=poporange, fill=poporange!20, text width=2.2cm, align=center, minimum height=1.2cm] (b2) at (18:4.2cm) {Comparing Topologies};
\node[circle, draw=popgreen, fill=popgreen!20, text width=2.2cm, align=center, minimum height=1.2cm] (b3) at (-54:4.2cm) {Network Equipment};
\node[circle, draw=poppurple, fill=poppurple!20, text width=2.2cm, align=center, minimum height=1.2cm] (b4) at (-126:4.2cm) {Setting Up a Network};
\node[circle, draw=poppink, fill=poppink!20, text width=2.2cm, align=center, minimum height=1.2cm] (b5) at (162:4.2cm) {Selecting Equipment};
% Connections
\draw[popdark, thick] (center) -- (b1);
\draw[popdark, thick] (center) -- (b2);
\draw[popdark, thick] (center) -- (b3);
\draw[popdark, thick] (center) -- (b4);
\draw[popdark, thick] (center) -- (b5);
\end{tikzpicture}
\legende{Mind map of Chapter ICT05}
\end{popfigure}

\begin{bilan}
\textbf{Key Definitions}
\begin{itemize}
\item \cle{Physical Topology}: The actual layout of cables and devices in a network.
\item \cle{Logical Topology}: The way data flows between devices, independent of physical layout.
\item \cle{Network Equipment}: Hardware devices (switch, router, access point, modem, firewall) that enable communication.
\item \cle{Structured Cabling}: Standardised cabling system (TIA/EIA-568) for reliable connectivity.
\end{itemize}

\textbf{Main Topologies \& Characteristics}
\begin{itemize}
\item \textbf{Bus}: Single backbone cable, easy to install, but a break disables whole network.
\item \textbf{Star}: Central switch/hub, fault isolation, scalable, most common in LANs.
\item \textbf{Ring}: Each node connects to two neighbours, token passing, deterministic but single point of failure.
\item \textbf{Mesh}: Every node connects to many others, high redundancy, costly, used in backbone/WAN.
\item \textbf{Hybrid}: Combination (e.g., star‑ring, star‑bus) to balance cost and reliability.
\end{itemize}

\textbf{Choosing a Topology – Decision Criteria}
\begin{itemize}
\item Cost (cabling, ports, equipment)
\item Scalability (adding nodes)
\item Fault tolerance (single point of failure)
\item Performance (bandwidth, latency)
\item Management complexity
\item Environment (office, campus, industrial)
\end{itemize}

\textbf{Key Network Equipment}
\begin{itemize}
\item \cle{Switch} (Layer 2): Forwards frames based on MAC address; creates collision domains per port.
\item \cle{Router} (Layer 3): Connects different networks, routes packets by IP address.
\item \cle{Access Point}: Provides wireless connectivity (Wi‑Fi 5/6/6E).
\item \cle{Modem}: Converts digital signals for ISP link (DSL, fibre, 4G/5G).
\item \cle{Firewall}: Filters traffic based on security policies.
\item \cle{Patch Panel / Keystone}: Terminates structured cabling.
\end{itemize}

\textbf{Setting Up a LAN – Step‑by‑Step Method}
\begin{enumerate}
\item Plan: Draw topology, list devices, estimate cable lengths.
\item Procure: Switch(es), router, cables (Cat6/6a), patch panels, rack, UPS.
\item Install structured cabling: Run horizontal cables to outlets, terminate on patch panel.
\item Connect devices: Patch panel → switch ports; switch uplink → router WAN.
\item Configure: Assign static IPs or DHCP scope, VLANs, port security, SSID/WPA3 on AP.
\item Test: Ping, traceroute, cable certifier, throughput test (iperf3).
\item Document: Label cables, update network diagram, record IP scheme.
\end{enumerate}

\textbf{Selecting Equipment for an Organisation}
\begin{itemize}
\item \textbf{Small office (≤20 users)}: Unmanaged switch, SOHO router, single AP, Cat6.
\item \textbf{Medium enterprise (20–200)}: Managed switches (VLAN, QoS, PoE+), dual‑WAN router, controller‑based APs, core switch with stacking.
\item \textbf{Large campus (>200)}: Core/distribution/access layer, Layer 3 switches, redundant links (STP/RSTP, VRRP), wireless controller cluster, firewall with IPS/IDS, network management (SNMP, NetFlow).
\end{itemize}

\textbf{Common Mistakes (GCE Traps)}
\begin{itemize}
\item Confusing physical vs. logical topology.
\item Thinking a hub is a switch (hub = shared collision domain).
\item Forgetting that a star topology still has a single point of failure at the central switch.
\item Misidentifying cable categories (Cat5e vs Cat6 vs Cat6a) and max lengths (100 m).
\item Neglecting to configure VLANs on a managed switch before connecting devices.
\item Assuming wireless is always slower; modern Wi‑6 can exceed 1 Gbps.
\end{itemize}

\textbf{GCE Examination Tips}
\begin{itemize}
\item Know the diagram symbols for each topology and be able to draw them.
\item Be ready to compare two topologies in a table (cost, fault tolerance, scalability).
\item Questions often ask: "Which topology for a school lab with 30 PCs?" → Star with managed switch.
\item Understand the OSI layers of each device: Switch (L2), Router (L3), Firewall (L3/L4), AP (L1/L2).
\item Practise writing a short network setup procedure (5–7 steps).
\item Remember Cameroon context: MTN/Orange fibre, CAMTEL backbone, ANTIC regulations for cybersecurity.
\end{itemize}
\end{bilan}

\findecours

\end{document}
